 \documentclass[10pt,a4paper]{article}
\usepackage[T1]{fontenc}
\usepackage[utf8]{inputenc}
\usepackage[french]{babel}
\usepackage{graphicx,fancyhdr,}
\usepackage{mathtools}
%\usepackage[left=1cm,right=1cm,top=1.5cm,bottom=1.5cm]{geometry}
\usepackage{amsfonts,amsmath,amssymb}
\usepackage{amsthm}
\usepackage{pstricks-add, pspicture} % pour les axces
\usepackage{thmtools}
\usepackage{xcolor}
\usepackage{nameref}
\usepackage{enumitem}
\usepackage{hyperref}
\usepackage{hyperref, fancybox, microtype,pifont,array, caption, bbding, multicol, longtable, indentfirst, tabularx, ,setspace, fourier, courier, color, multirow, lscape,lipsum,textcomp}
\usepackage{tikz}
\usepackage{systeme}
\usepackage{pgfplots}
\pgfplotsset{compat=1.15}
\usepackage{mathrsfs}
\usetikzlibrary{arrows}
\pagestyle{empty}


%<<<<<<<WARNING>>>>>>>
% PGF/Tikz doesn't support hatch filling very well
% Use PStricks for a perfect hatching export

\usetikzlibrary[patterns]

\usepackage[normalem]{ulem} % pour hachurer
\usepackage{colortbl} %Pour hachurer une case d'un tableau 
\usepackage[tikz]{bclogo}
\usepackage[most]{tcolorbox} 
\usepackage{longtable}
%\usepackage{background} %filigrane
\usepackage{draftwatermark} %filigrane
\usepackage{fancybox,fancyhdr}%encadrement et les titres des en-tete
%\usepackage{multicol, multirow}%fusion de colone et ligne
\usepackage[np]{numprint}%pour écrire correctement les nombres à virgules sans espaces entre la virgule et le nombre tout en utilisant le commande \np{}
%\usepackage{bookman, fourier}

%\setlength{\columnseprule}{0.5pt}
\setlength{\columnsep}{3pt}
\usepackage{graphicx}
\usepackage[left=1.00cm, right=1.00cm, top=1.5cm, bottom=1.5cm]{geometry}
\newcounter{exo}
\setcounter{exo}{0}
\newcommand{\exo}[1]{
\stepcounter{exo} { \textbf{\underline{Exercice de maison \No}\arabic{exo}} : \textbf{#1} \newline }}
%donnez une fraction aleatoire à simplifier debut:
\newcounter{Num} \newcounter{Den}
\newcounter{Coef}
\newcommand{\FractAleat}{%
	\reinitrand[first=1,last=10,counter=Num]\rand
	\reinitrand[first=2,last=10,counter=Den]\rand
	\reinitrand[first=2,last=12,counter=Coef]\rand
	\setcounter{Num}{\value{Num}*\value{Coef}}
	\setcounter{Den}{\value{Den}*\value{Coef}}
	$\dfrac{\theNum}{\theDen}$}
%fin
%les activités et les consignes debut
\newcounter{q}
\setcounter{q}{0}
\newcounter{qq}
\newcommand{\act}[1]
{ 
		 \setcounter{qq}{0}
	\addtocounter{q}{1} \doublebox {\textbf{Activité} \theq.}\space \textbf{{#1}}\\
 }
\newcommand{\conAct}[1]{
	\addtocounter{qq}{1} \ovalbox{\textbf{Consignes	\theq.\space\theqq.}}  \textbf{{#1}} \\
}
\newcommand{\resul}{ \begin{center}
		\rule{2cm}{1mm} 
		\shadowbox{Résolutions}
  \rule{2cm}{1mm}
\end{center}}
%definition et autre 
\newcommand{\defs}[1]{\begin{bclogo}[couleur=green!20,arrondi=5pt]
		{\textbf{Définition}}
		#1
	\end{bclogo} 	 
}
%nb ou vocabulaire 
\newcommand{\nbVo}[2]{\begin{bclogo}[couleur = lightgray!15, couleur =red!15, arrondi = 0.3, ombre = true, epOmbre =0.25, couleurOmbre= lightgray!90, logo=\bcrecyclage, epBarre=3  ]{\textbf{#1}}
		#2
\end{bclogo}}
%Propriété
\newcommand{\propr}[1]{  \begin{bclogo}[couleur = lightgray!15, arrondi = 0.3, ombre = true, epOmbre =0.25,  nobreak=false, couleurOmbre= lightgray!90, logo=\bcrecyclage, epBarre=3 ]{Propriété} #1
\end{bclogo}}
%suitepro
\newcommand{\prosuit}[1]{  \begin{bclogo}[couleur = lightgray!15, arrondi = 0.3, ombre = true, epOmbre =0.25, couleurOmbre= lightgray!90, logo=\bcrecyclage,  epBarre=3 ]{   } #1
\end{bclogo}}
%attention 
\newcommand{\atten}[1]{  \begin{bclogo}[couleur = red!25, arrondi = 0.3, ombre = true, epOmbre =0.25, couleurOmbre= lightgray!90, logo=\bcattention, epBarre=3 ]{\textcolor{red!100}{ATTENTION}} #1
\end{bclogo}}
%les retenons 
\newcommand{\ret}[1]{\begin{tcolorbox}[colback=violet!10,colframe=blue!40,title=Retenons]
		#1 
\end{tcolorbox}}
	\newcommand{\dure}{ \underline{Stratégies et Durées:} TI : \dots min TG: \dots min   TC : \dots min \\
	\ding{36}\dotfill{} \\
	\hspace*{1cm}	\rule{2cm}{1mm} 
		\shadowbox{Résolutions}
		\rule{2cm}{1mm} \\}
%titre des sequence 
	\newcommand{\dures}{ \underline{Stratégies et Durées:} TI : \dots min TG: \dots min   TC : \dots min \\
	\ding{36}\dotfill{} }
\newcounter{sequen}
\setcounter{sequen}{0}
\newcommand{\sequence}[1]{
	\stepcounter{sequen}{
		\begin{bclogo}[couleur=cyan!50,arrondi=0.3, ombre = true, epOmbre =0.25,logo=\bcsoleil,epBarre=3.5]{ \LARGE Séquence \no\thesequen : }
			\Large \textbf{#1}
\end{bclogo} } }
	\newcounter{nu}
\setcounter{nu}{0}
\newcommand{\seq}{\addtocounter{nu}{0} 
	\textbf{\fcolorbox{red}{yellow}{{\LARGE Séquence} \no\thenu} \hspace{0.5cm}}}
%titre de la SA
\newcounter{situa}
\setcounter{situa}{0}
\newcommand{\situ}[1]{
	\stepcounter{situa}{
 \begin{bclogo}[couleur=blue!30,arrondi=0.3, ombre = true, epOmbre =0.25,logo=\bcbook,epBarre=3.5]{ \LARGE Situation d'Apprentissage \no\thesitua : }
 	\Large \textbf{#1}
\end{bclogo} } }
\setlength{\columnsep}{30pt}
%inclu
\newcommand{\nincl}{$\subset  $ \hspace*{-0.4cm}/ }
\newcommand{\reto}{ 
		 {\Large RETOUR ET PROJECTION} \\
	\ding{117}	Dis ce que tu as appris ? \\
	\ding{117} Fait part de tes réussites puis de tes difficultés et la façon dont tu  les  as surmontées.\\
	\ding{117}	Dis ce en quoi ces résultats peuvent être utiles pour toi dans la vie courante\\
	\Ovalbox{\textbf{ Durée}:\dots }
}
\newcommand{\doc}{\begin{center}
		\section*{Documents utilisés}
	\end{center}
	\ding{117}	Guide et programme de la classe de troisième;\\
	\ding{117}	Réussir en mathématiques troisième;\\
	\ding{117} 	CIAM troisième.}	
\newcommand{\sph}{
 \begin{tikzpicture}[scale=0.7]
	\draw[color=red] (0,0) circle (1cm);
	\draw[dashed, color=red] (1,0) arc (0:180:1cm and 0.5cm);
	\draw[color=red] (-1,0) arc (180:360:1cm and 0.5cm);
	\draw[color=red] (0,-1) arc (270:90:0.5cm and 1cm);
	\draw[dashed, color=red!50] (0,1) arc (270:90:-0.5cm and -1cm);
\end{tikzpicture}}
\newcommand{\rect}{
\begin{tikzpicture}
	\draw[color=blue] (0,0)rectangle (0.7,0.7);
	\draw [color=blue] (0.7,0.7)--(1.5,1.5)--(1.5,0.8)--(0.7,0);
	\draw[color=blue] (1.5,1.5)--(0.8,1.5)--(0,0.7) ;
\end{tikzpicture}}
\newcommand{\cone}{
\begin{tikzpicture}[scale=0.7]
	\draw[dashed, color=orange] (0,0) arc (0:180:1cm and 0.5cm);
	\draw[color=orange] (-2,0) arc (0:180:-1cm and -0.5cm);
	\draw[color=orange](-2,0)--(-1,2)--(0,0);
\end{tikzpicture}}
\newcommand{\rapot}{
		%\draw (0,0) circle (5cm);
		%	\draw (0,0)--(5:5) node[right, font=\scriptsize] {5};
		\draw[line width=2pt] (0,0)--(0:5) node[right, font=\scriptsize]  {0};
		\draw (0,0)--(5:5) node[right, font=\scriptsize] {5} ;
		\draw (0,0)--(10:5) node[right, font=\scriptsize] {10};
		\draw (0,0)--(15:5) node[right, font=\scriptsize]{15} ;
		\draw (0,0)--(20:5) node[right, font=\scriptsize]{20} ;
		\draw (0,0)--(25:5) node[right, font=\scriptsize]{25} ;
		\draw (0,0)--(30:5) node[right, font=\scriptsize] {30};
		\draw (0,0)--(35:5) node[right, font=\scriptsize] {35};
		\draw (0,0)--(40:5) node[right, font=\scriptsize]{40} ;
		\draw (0,0)--(45.0:5) node[above right, font=\scriptsize]{45} ;
		\draw (0,0)--(50:5)node[above right, font=\scriptsize]{50} ;
		\draw (0,0)--(55.0:5)node[above right, font=\scriptsize]{55} ;
		\draw (0,0)--(60:5)node[above right, font=\scriptsize] {60};
		\draw (0,0)--(65.0:5)node[above right, font=\scriptsize] {65};
		\draw (0,0)--(70:5)node[above right, font=\scriptsize]{70} ;
		\draw (0,0)--(75.0:5) node[above, font=\scriptsize]{75};
		\draw (0,0)--(80:5)node[above , font=\scriptsize]{80};
		\draw (0,0)--(85:5)node[above, font=\scriptsize]{85};
		\draw[line width=2pt] (0,0)--(90:5)node[above, font=\scriptsize]{90};
		\draw (0,0)--(175:5) node[left, font=\scriptsize] {175} ;
		\draw (0,0)--(170:5) node[left, font=\scriptsize] {170};
		\draw (0,0)--(165:5) node[left, font=\scriptsize]{165} ;
		\draw (0,0)--(160:5) node[left, font=\scriptsize]{160} ;
		\draw (0,0)--(155:5) node[left, font=\scriptsize]{155} ;
		\draw (0,0)--(150:5) node[left, font=\scriptsize] {150};
		\draw (0,0)--(145:5) node[left, font=\scriptsize] {145};
		\draw (0,0)--(140:5) node[left, font=\scriptsize]{140} ;
		\draw (0,0)--(135.0:5) node[above left, font=\scriptsize]{135} ;
		\draw (0,0)--(130:5)node[above left, font=\scriptsize]{130} ;
		\draw (0,0)--(125.0:5)node[above left, font=\scriptsize]{125} ;
		\draw (0,0)--(120:5)node[above left, font=\scriptsize] {120};
		\draw (0,0)--(115.0:5)node[above left, font=\scriptsize] {115};
		\draw (0,0)--(110:5)node[above left, font=\scriptsize]{110} ;
		\draw (0,0)--(105.0:5) node[above, font=\scriptsize]{105};
		\draw (0,0)--(100:5)node[above , font=\scriptsize]{100};
		\draw (0,0)--(95:5)node[above, font=\scriptsize]{95};
		\draw [line width=2pt](0,0)--(180.0:5) node[left, font=\scriptsize] {180};
		\draw [fill=white](4,0) arc (0:180:4);
		\draw [](5,0) arc (0:180:5);
		\draw [line width=2pt] (0,0)--(90:4);
		\draw [line width=2pt] (0,0) rectangle (0.2,0.2);
}
\newcommand{\pavdroi}{ 
		\draw  (0,0)rectangle (4,2);
		\draw  (4.7,2.7)--(0.7,2.7);
		\draw[dashed,] (0,0)--(0.7,0.7);
		\draw  [color=black] (4,0)--(4.7,0.7) ;
		\draw  (0,2)--(0.7,2.7);
		\draw  (4,02)--(4.7,2.7) ;
		\draw [dashed] (0.7,0.7)--(0.7,2.7);
		\draw [dashed] (0.7,0.7)--(4.7,0.7);
		\draw  (0,0)--(0.7,0.7)--(4.7,0.7)--(4,0)--cycle;
		\draw  (0,2)--(0.7,2.7)--(4.7,2.7)--(4,2)--cycle;
		\draw  (4.7,0.7)-- (4.7,2.7) ;
		\draw (0,0)--(4,0);
}
\newcommand{\cube}{\begin{tikzpicture}[x={(-0.572cm,-0.416cm)},y={(1cm,0cm)},z={(0cm,1cm)}]
		\draw  (0,0,0)--(1,0,0)--(1,1,0)--(0,1,0)--cycle;
		\draw   (0,0,1)--(1,0,1)--(1,1,1)--(0,1,1)--cycle; 
		\draw (0,0,0)--(0,0,1);
		\draw (1,0,0)--(1,0,1);
		\draw (0,1,0)--(0,1,1); 
		\draw (1,1,0)--(1,1,1) ; 
		\draw  (1,0,0)--(1,1,0) ;
		\draw  (1,1,0)--(0,1,0);
	
\end{tikzpicture}}
 \SetWatermarkText{KAYCT}
 \pagestyle{fancy} % pour afficher les pieds et entêtes de page
 \fancyhead[L]{\textbf{Support d'activités  3\up{è}M3}}
  \fancyhead[C]{\textbf{CEG Atrokpocodji }}
 \fancyfoot[L]{\textbf{On acquiert en pratiquant}}
 \fancyfoot[R]{\textbf{BEPC: 99\% de travail et 1\% de chance}}
 
 \begin{document}
 	\newgeometry{left=0.0cm, right=0.00cm, top=0.0cm, bottom=0.0cm}
 	\begin{titlepage}
 		\SetWatermarkText{	\begin{tikzpicture}
 				%	\draw (-9,-14) grid (9,13);
 				%	\draw[->, line width=2pt] (0,-14)--(0,13);
 				%	\draw[->, line width=2pt] (-9,0)--(9,0);
 				%\draw [top color=red!30,bottom color=blue!30, line width=0pt, color=red!30] (-11,0)--(0,15)--(-11,15)--(-11,0) ;
 				%	\draw [top color=blue!30,bottom color=yellow, color=blue!30] (-9,3)--(9,-14)--(-9,-14)--(-9,3) ;	\draw [ball color =cyan!30, line width=0pt, color=red!30] (-9,0) arc (270:360:8);	\draw[ball color =cyan!30, line width=0pt, color=red!30] (-1,8) arc (270:360:5);	\draw[fill=red!30, line width=0pt, color=red!30] (-9,0)--(4,13)--(0,13)--(-9,0);
 				\draw (-11.5,-15) grid (11.5,15);
 				\draw[ball color=pink] (-11.5,0) --(0,15)--(-11.5,15)--(-11.5,0);
 				\draw[ball color=pink] (-11.5,5) --(0,-15)--(-11.5,-15)--(-11.5,5);
 				\draw[ball color=pink] (-11.5,5) --(0,-15)--(-11.5,-15)--(-11.5,5);
 				\draw[->, line width=2pt] (-9,0)--(9,0);
 				\draw[fill=white!50, color=white] (-10,2) --(0,15)--(11.5,15)--(11.5,0)--(-10,2);
 				\draw[fill=white!50, color=white] (-10,2)--(11.5,0)--(11.5,-15)--(0,-15)--(-10,2);
 				\draw [color=black, line width=1pt] plot (\x, \x/4+3);
 				\draw [color=black, line width=1pt] plot (\x, \x/4+6)  ;
 				\rapot
 				\draw[color=yellow, line width=20pt] (-8,8)--(-3,8)--(-8,13)--cycle;
 				%	\draw[ball color=pink] (-11.5,-15) --(11.5,-15)--(11.5,0)--(-11.5,0)--(-11.5,15);
 		\end{tikzpicture}}
 		%	\SetWatermarkText{\begin{tikzpicture}
 				%	\draw[color=red] (0,0) rectangle (3,3);
 				%	\end{tikzpicture}}
 		\SetWatermarkAngle{0}
 		\vspace*{1cm}
 		\begin{center}
 			\textcolor{blue}{\textbf{\resizebox{15cm}{1.5cm}{SUPPORT D'ACTIVITES}}}
 			\vspace*{0.5cm}
 			\textcolor{blue}{\textbf{\resizebox{12cm}{1.5cm}{DE LA CLASSE DE TROIXIÈME}}}
 		\end{center}
 		\begin{center}
 			 
 		\end{center}
 		\begin{center}
 			{	\Huge
 				\textcolor{red}{\textbf{\resizebox{15cm}{2cm}{MATHÉMATIQUES}}}} \vspace*{-1cm}
 			\textcolor{red!50}{\textbf{\resizebox{15cm}{-1cm}{\textit{MATHÉMATIQUES}}}}
 		\end{center}
 	
 		 
 	\definecolor{ffffff}{rgb}{1,1,1}
 	\definecolor{zzttqq}{rgb}{0.6,0.2,0}
  \hspace*{3cm}	\begin{tikzpicture}[scale=0.6]
 		\fill[line width=2pt,color=zzttqq,fill=zzttqq,fill opacity=0.03] (-11.052201811652381,0.5367821746980063) -- (-6.138647431701272,0.550659108986587) -- (-4.34465327034393,2.344653270343928) -- (-9.235387015725221,2.3535969706251674) -- cycle;
 		\fill[line width=1.2pt,color=ffffff,fill=ffffff,fill opacity=0.14] (-9.003438016933359,1) -- (-6.858008450642274,1.006059154495914) -- (-5.999429059667906,1.864638545470281) -- (-8.138799471463077,1.864638545470281) -- cycle;
 		\fill[line width=0.8pt,color=ffffff,fill=ffffff,fill opacity=0.1] (-13,-6) -- (-4,-6) -- (-2,-2) -- (-11,-2) -- cycle;
 		\fill[line width=2pt,color=zzttqq,fill=zzttqq,fill opacity=0.02] (1.5943737545226166,0.6215983936426223) -- (7.549683454967029,0.5868735557391562) -- (9.25120051223686,2.323115450912458) -- (3.313253230744182,2.323115450912458) -- cycle;
 		\fill[line width=2pt,color=zzttqq,fill=zzttqq,fill opacity=0.04] (-0.009368600793884364,-6.054526279032764) -- (8.722019707498214,-6.070493396470353) -- (10.768285267741785,-1.9779622759832103) -- (2.0165550165331974,-2.0026790443786004) -- cycle;
 		\fill[line width=0.8pt,color=zzttqq,fill=zzttqq,fill opacity=0.04] (-13,-6) -- (-3.1838877119220834,-6.048721025885665) -- (-2,-2) -- (-11,-2) -- cycle;
 		\draw [line width=0.8pt,dash pattern=on 1pt off 1pt] (-9,-3)-- (-6,-5);
 		\draw [line width=0.8pt,dash pattern=on 1pt off 1pt] (-11,-5)-- (-4,-3);
 		\draw [line width=0.8pt,dash pattern=on 1pt off 1pt] (-7.5,-4)-- (-7.5,5.5);
 		\draw [line width=0.8pt] (-7.5,5.5)-- (-4,-3);
 		\draw [line width=0.8pt] (-7.5,5.5)-- (-6,-5);
 		\draw [line width=0.8pt,dash pattern=on 1pt off 1pt] (-7.5,5.5)-- (-9,-3);
 		\draw [line width=0.8pt] (-7.5,5.5)-- (-11,-5);
 		\draw [line width=0.8pt,color=zzttqq] (-11.052201811652381,0.5367821746980063)-- (-6.138647431701272,0.550659108986587);
 		\draw [line width=0.8pt,color=zzttqq] (-6.138647431701272,0.550659108986587)-- (-4.34465327034393,2.344653270343928);
 		\draw [line width=0.4pt,color=zzttqq] (-4.34465327034393,2.344653270343928)-- (-9.235387015725221,2.3535969706251674);
 		\draw [line width=0.8pt,color=zzttqq] (-9.235387015725221,2.3535969706251674)-- (-11.052201811652381,0.5367821746980063);
 		\draw [line width=0.8pt,color=ffffff] (-9.003438016933359,1)-- (-6.858008450642274,1.006059154495914);
 		\draw [line width=0.8pt,color=ffffff] (-6.858008450642274,1.006059154495914)-- (-5.999429059667906,1.864638545470281);
 		\draw [line width=0.8pt,color=ffffff] (-13,-6)-- (-4,-6);
 		\draw [line width=0.8pt,color=ffffff] (-4,-6)-- (-2,-2);
 		\draw [line width=0.8pt,color=ffffff] (-11,-2)-- (-13,-6);
 		\draw [rotate around={0:(5.370407472006344,-4.010355913823057)},line width=0.8pt] (5.370407472006344,-4.010355913823057) ellipse (2.6180020600442737cm and 1.6893592828040032cm);
 		\draw [line width=0.4pt] (5.402895777180931,5.608754933471161)-- (2.752761970585546,-3.9824752754810837);
 		\draw [line width=0.8pt] (5.402895777180931,5.608754933471161)-- (7.9883613758387195,-4.000109398599252);
 		\draw [line width=0.4pt] (5.402895777180931,5.608754933471161)-- (6.940846645602726,-5.362018533797325);
 		\draw [rotate around={-0.5583183624421316:(5.38985656254413,1.1031124838499535)},line width=0.8pt] (5.38985656254413,1.1031124838499535) ellipse (1.239085586525175cm and 0.16042788298019578cm);
 		\draw [line width=0.8pt,dash pattern=on 1pt off 1pt] (4.161258725288073,1.115084926684446)-- (6.618454399800344,1.0911400410154597);
 		\draw [line width=0.8pt,dash pattern=on 1pt off 1pt] (5.387720306248826,1.1031333012448732)-- (6.054399470343201,0.9613376605730773);
 		\draw [line width=0.8pt,color=zzttqq] (1.5943737545226166,0.6215983936426223)-- (7.549683454967029,0.5868735557391562);
 		\draw [line width=0.8pt,color=zzttqq] (7.549683454967029,0.5868735557391562)-- (9.25120051223686,2.323115450912458);
 		\draw [line width=0.8pt,color=zzttqq] (9.25120051223686,2.323115450912458)-- (3.313253230744182,2.323115450912458);
 		\draw [line width=0.8pt,color=zzttqq] (3.313253230744182,2.323115450912458)-- (1.5943737545226166,0.6215983936426223);
 		\draw [line width=0.8pt,color=zzttqq] (-0.009368600793884364,-6.054526279032764)-- (8.722019707498214,-6.070493396470353);
 		\draw [line width=0.8pt,color=zzttqq] (8.722019707498214,-6.070493396470353)-- (10.768285267741785,-1.9779622759832103);
 		\draw [line width=0.8pt,color=zzttqq] (10.768285267741785,-1.9779622759832103)-- (2.0165550165331974,-2.0026790443786004);
 		\draw [line width=0.8pt,color=zzttqq] (2.0165550165331974,-2.0026790443786004)-- (-0.009368600793884364,-6.054526279032764);
 		\draw [line width=0.8pt,dash pattern=on 1pt off 1pt] (2.752761970585546,-3.9824752754810837)-- (7.9883613758387195,-4.000109398599252);
 		\draw [line width=0.8pt,dash pattern=on 1pt off 1pt] (5.402895777180931,5.608754933471161)-- (5.45254713434296,-3.9915684738496475);
 		\draw [line width=0.8pt,dash pattern=on 1pt off 1pt] (5.45254713434296,-3.9915684738496475)-- (6.940846645602726,-5.362018533797325);
 		\draw [line width=2pt] (-6.858008450642273,1.0060591544959137)-- (-9.003438016933359,1);
 		\draw [line width=2pt] (-6.858008450642273,1.0060591544959137)-- (-5.999429059667906,1.864638545470281);
 		\draw [line width=1.2pt,dash pattern=on 1pt off 1pt] (-5.999429059667906,1.864638545470281)-- (-8.138799471463077,1.864638545470281);
 		\draw [line width=1.2pt,dash pattern=on 1pt off 1pt] (-8.138799471463077,1.864638545470281)-- (-9.003438016933359,1);
 		\draw [line width=0.8pt] (-11,-5)-- (-6,-5);
 		\draw [line width=0.8pt] (-6,-5)-- (-4,-3);
 		\draw [line width=1.2pt,dash pattern=on 1pt off 1pt] (-4,-3)-- (-9,-3);
 		\draw [line width=1.2pt,dash pattern=on 1pt off 1pt] (-9,-3)-- (-11,-5);
 		\draw [line width=0.8pt,color=zzttqq] (-13,-6)-- (-3.1838877119220834,-6.048721025885665);
 		\draw [line width=0.8pt,color=zzttqq] (-3.1838877119220834,-6.048721025885665)-- (-2,-2);
 		\draw [line width=0.8pt,color=zzttqq] (-2,-2)-- (-11,-2);
 		\draw [line width=0.8pt,color=zzttqq] (-11,-2)-- (-13,-6);
 		\draw [line width=1.2pt,dash pattern=on 1pt off 1pt] (-8.140742300143643,1.8691269658526934)-- (-6.858008450642273,1.0060591544959137);
 		\draw [line width=1.2pt,dash pattern=on 1pt off 1pt] (-8.999996760383807,1.0000097188485797)-- (-6.002019718165304,1.8620478869728814);
 		\draw [line width=0.8pt] (-8.985529408392216,-6.394533579553936)-- (-6.820716071714553,-6.399792508027989);
 		\draw [line width=0.8pt] (-8.982492134065824,-8.15598445226318)-- (-6.811717667091587,-8.186775579454302);
 		\draw [line width=0.8pt] (-8.985529408392216,-6.394533579553936)-- (-8.982492134065824,-8.15598445226318);
 		\draw [line width=0.8pt] (-6.820716071714553,-6.399792508027989)-- (-6.811717667091587,-8.186775579454302);
 		\draw [line width=1.2pt,dash pattern=on 1pt off 1pt] (-8.998288674427224,1.0051339767183267)-- (-8.985529408392216,-6.394533579553936);
 		\draw [line width=1.2pt,dash pattern=on 1pt off 1pt] (-6.858008450642275,1.0060591544959125)-- (-6.820716071714553,-6.399792508027989);
 		\draw [line width=0.8pt] (-11,-9.94246507628567)-- (-6,-10.001714871312494);
 		\draw [line width=0.8pt] (-11,-15)-- (-6,-15);
 		\draw [line width=1.2pt,dash pattern=on 1pt off 1pt] (-11,-5)-- (-11,-9.94246507628567);
 		\draw [line width=1.2pt,dash pattern=on 1pt off 1pt] (-6,-5)-- (-6,-10.001714871312494);
 		\draw [line width=0.8pt] (-11,-9.94246507628567)-- (-11,-15);
 		\draw [line width=0.8pt] (-6,-10.001714871312494)-- (-6,-15);
 		\draw [line width=0.8pt] (5.37505079538736,-7.619736237782518) circle (1.2711599258284343cm);
 		\draw [line width=0.8pt] (5.483234618862127,-12.650284029358872) circle (2.861508144451029cm);
 		\draw [line width=0.8pt,dash pattern=on 1pt off 1pt] (2.752761970585546,-3.9824752754810837)-- (2.8138759613587063,-11.619384801320034);
 		\draw [line width=0.8pt,dash pattern=on 1pt off 1pt] (7.9883613758387195,-4.000109398599252)-- (7.996837511301049,-11.282793861247159);
 		\draw [line width=0.8pt,dash pattern=on 1pt off 1pt] (2.6217722352735624,-12.666466985430665)-- (8.344741630066874,-12.652830710660519);
 		\draw [line width=0.8pt,dash pattern=on 1pt off 1pt] (5.483234618862127,-12.650284029358872)-- (7.2649232034247575,-14.889438925385225);
 		\draw [line width=0.8pt,dash pattern=on 1pt off 1pt] (4.1038923926895325,-7.617768426089483)-- (6.646146134990301,-7.606922390903581);
 		\draw [line width=0.8pt,dash pattern=on 1pt off 1pt] (6.054399470343201,0.9613376605730773)-- (6.122024642307467,-8.648268020981037);
 		\draw [line width=0.8pt,dash pattern=on 1pt off 1pt] (6.622714145889668,1.0753087252622615)-- (6.515496103204996,-7.058280935101274);
 		\draw [line width=0.8pt,dash pattern=on 1pt off 1pt] (4.156849980560731,1.0991290207950106)-- (4.274094414426187,-6.9843288060876825);
 		\draw [line width=0.8pt,dash pattern=on 1pt off 1pt] (5.452547134342961,-3.9915684738496484)-- (5.483234618862127,-12.650284029358872);
 		\draw [line width=0.8pt,dash pattern=on 1pt off 1pt] (5.465378218286185,-7.611959909447801)-- (6.122024642307467,-8.648268020981037);
 		\draw [->,line width=0.8pt] (-1.2006136673959062,0.20814161325148525) -- (-4.7040039462370125,1.8067760123342225);
 		\draw [->,line width=0.8pt] (-1.2006136673959062,0.20814161325148525) -- (2.064682126475416,1.3305870423946837);
 		\draw [->,line width=0.8pt] (-1.6427891394826477,-7.274827914369838) -- (-4.057747487033313,-6.288436476637936);
 		\draw [->,line width=0.8pt] (-1.6427891394826477,-7.274827914369838) -- (0.6701287145095387,-6.186395983079464);
 		\draw (-1.5,0.2) node[anchor=north west] {$(\mathcal{P})$};
 		\draw (-3.1,-7.4) node[anchor=north west] {Plan de la base};
 		\draw (-11.9,-4.4) node[anchor=north west] {$A$};
 		\draw (-5.6,-4.6) node[anchor=north west] {$B$};
 		\draw (-3.9,-2.4) node[anchor=north west] {$C$};
 		\draw (-10,-2.4) node[anchor=north west] {$D$};
 		\draw (-10,1.7) node[anchor=north west] {$A'$};
 		\draw (-6.7,1.45) node[anchor=north west] {$B'$};
 		\draw (-6,2.4) node[anchor=north west] {$C'$};
 		\draw (-9.2,2.4) node[anchor=north west] {$D'$};
 		\draw (-8.2,6) node[anchor=north west] {$S$};
 		\draw (-11,-15.5) node[anchor=north west] {Figure 1};
 		\draw (4,-16) node[anchor=north west] {Figure 2};
 		\draw [line width=0.8pt,dash pattern=on 1pt off 1pt] (6.940846645602726,-5.362018533797325)-- (7.2649232034247575,-14.889438925385225);
 	\end{tikzpicture}
 	
 		\begin{center}
 			\resizebox{10cm}{5mm}{\textcolor{black}{\Large \textbf{Réalisé par M.}  \textbf{TCHACON Kayodé Carmel}}}
 		\end{center}
 		\vspace*{1cm}
 		\begin{center}
 		{\Large CEG Atrokpocodji  \today }
 		\end{center}
 	\end{titlepage}
 	\restoregeometry
  \begin{center}
  	\section*{PROGRAMME }
  \end{center}
  \section*{SITUATION D’APPRENTISSAGE \No1 : TRIANGLES}
  
  \underline{SÉQUENCE \No 1} : NOMBRES RÉELS
  
  \underline{SÉQUENCE \No 2} : VALEUR ABSOLUE
  
  \underline{SÉQUENCE \No 3} : TRIGONOMÉTRIE
  
  \underline{SÉQUENCE \No 4} : PROPRIÉTÉS DE THALÈS RELATIVES AU TRIANGLE
  
  \underline{SÉQUENCE \No 5}: TRIANGLES SEMBLABLES
  
  \underline{SÉQUENCE \No 6}: TRIANGLES RECTANGLES 
  
  \underline{SÉQUENCE \No 7}: ANGLES ET CERCLE
  
  \section*{SITUATION D’APPRENTISSAGE \No2 : CONFIGURATION DE L’ESPACE} 
  
  \underline{SÉQUENCE \No 1}: CÔNE
  
  \underline{SÉQUENCE \No 2}: SECTION PLANE
  
  \section*{SITUATION D’APPRENTISSAGE \No3 : CALCUL LITTÉRAL }
  
  \underline{SÉQUENCE \No 1} : POLYNÔMES
  
  \underline{SÉQUENCE \No 2} : ÉQUATIONS DE DROITES 
  
  \underline{SÉQUENCE \No 3} : ÉQUATIONS ET INÉQUATIONS 
  
  \underline{SÉQUENCE \No 4} :SOMME DE DEUX VECTEURS - PRODUIT D'UN VECTEUR PAR UN NOMBRE RÉEL 
  
  \underline{SÉQUENCE \No 5}: CALCUL SUR LES COORDONNÉES DE VECTEURS 
  
  \section*{SITUATION D’APPRENTISSAGE \No 4 : ORGANISATION DES DONNÉES}
  
  \underline{SÉQUENCE \No 1} : APPLICATION AFFINE
  
  \underline{SÉQUENCE \No 2} : APPLICATION LINÉAIRE
  
  \underline{SÉQUENCE \No 3} : STATISTIQUE
  
  \begin{tabular}{|l|c|c|c|c|c|}
  	\cline{2-6}
  	\multicolumn{1}{c|}{} & $SA_{1}$ & $SA_{2}$&$SA_{3}$&$SA_{4}$ &Total\\
  	\hline
  	\textbf{Durée} & 54 H& 18 H& 42 H& 18H& 132 H\\ 
  	\hline
  	\textbf{Nombre de séances} & 18 & 6& 14 & 6 &44\\
  	\hline 
  	\textbf{Nombre de semaines }& 9 & 3& 7& 3 & 22\\
  	\hline 
  \end{tabular}
  
  \vspace{0.5cm}
  Chaque séance est prévue pour une durée de 3 heures.
  
  \rhead{2026-2027}
  \newpage
  \begin{center}
  	\resizebox{18cm}{0.8cm}{ \textcolor{black}{CONTRAT  DIDACTIQUE}}
  \end{center}
  
  Le \dots \dots /\dots \dots/20\dots\dots, nous élèves de la classe de \dots \dots du CEG \dots\dots\dots\dots\dots\dots\dots\dots  avons pris contact avec notre professeur de Mathématiques. Pour une bonne marche des activités pédagogiques, nous avons convenu des points suivants:
  \begin{enumerate}
  	
  	\item  Nous devons balayer la classe, la cour et effacer le tableau avant l'arrivée du professeur;
  	\item Pas de retard au cours ni d'absence non justifiée;
  	\item Pas de bavardage ni de perturbation du cours;
  	\item Toute prise de parole est subordonnée à une autorisation préalable du professeur;
  	\item Éviter de se moquer de ses camarades lorsqu'ils ont la parole;
  	\item  Apprendre au jour le jour ses leçons et traiter ses exercices;
  	\item  Une interrogation ou un devoir non fait et non justifié est sanctionné par la note zéro;
  	\item Avoir une bonne hygiène vestimentaire et corporelle;
  	\item Nous devons éviter le port des portables  dans la classe et dans l'enceinte  de l'établissement;
  	\item Nous devons respecter nos camarades, les enseignants et les autorités du collège;
  	\item Nous devons éviter tout comportement déviant et non conforme à la charte du collège;
  	\item Nous ne devons pas oublier nos cahiers (exercice et cours) à la maison;
  	\item  Nous devons avoir nos supports à chaque cours.
  \end{enumerate}
  \textbf{En cas de non respect de l'un de ces engagements, je serai puni(e) conformément au règlement intérieur de l'établissement.}
  
  
  \vspace{1cm}
  \begin{center}
  	\textbf{Le Professeur} \hspace{11cm}    \textbf{L'élève}
  \end{center}
  \newpage
  \begin{multicols*}{2}
  		\SetWatermarkText{\resizebox{28cm}{0.35cm}{Support d'activités 3\up{ème} }}
  	\SetWatermarkScale{1}
  	\SetWatermarkColor{black!30}
  	\SetWatermarkAngle{90}
  	\situ{TRIANGLES}
  	\underline{\textbf{Situation de départ}} \\
  	\textbf{Texte}: Le recasement 
  	
  	
  	\hspace{0.5cm} La ville de NOUNAGNON est scindée en deux zones $A$ et $B$. Après l'installation des propriétaires terrains de la zone $A$, la ville est en manque d'espace pour certaines infrastructures. Pour pallier ce problème, lors du lotissement de la zone $B$, la mairie a décidé d'appliquer un coefficient de réduction à chaque parcelle. Suite à cette opération, il a été prélevé $\dfrac{2}{5}$ de chaque parcelle afin de construire certaines infrastructures comme jardins publics et centre de santé.
  	
  	Après le lotissement, le conseil municipal a fait aménager un jardin public de forme circulaire de 120 m de diamètre. Au centre de ce jardin, un pilonne a été érigé. 
  	Finagnon possédait un terrain rectangulaire de 20 m sur 30 m. A l'issue des travaux de recasement, il lui a été attribué un terrain carré. Dowa, l'une des fille de Finagnon ayant le niveau de la quatrième, souhaite connaître les dimensions de la parcelle attribuée à son père. Par ailleurs, impressionnée par la structure du pylône elle voudrait en connaître la hauteur.  
  	
  	
  	\Ovalbox{\textbf{Tâche}} : A partir cette situation, tu vas te construire de nouvelles connaissances en mathématique.
  	
  	Pour cela, tu auras tout au long de la situation d’apprentissage à :\\ 
  	\ding{43} Exprimer ta perception de chacun des problèmes posés ;\\ 
  	\ding{43} Analyser chacun des problèmes posés ; \\
  	\ding{43} Mathématiser chacun des problèmes posés ; \\
  	\ding{43} Opérer sur l’objet mathématique que tu as identifié pour chaque problème ; \\
  	\ding{43} Améliorer au besoin ta production.
  	\begin{center}
  		\section*{I-INTRODUCTION}
  	\end{center}  
  	\doublebox{\textbf{Activité 0}} \textbf{ Brainstorming}\\
  	\begin{enumerate}
  		\item  Lis le texte de la situation de départ 
  		\item  Reformule la situation-problème en tes 
  		propres termes 
  		\item  Formule toutes les idées ou questions que 
  		t’inspire la situation de départ.
  	\end{enumerate}
  	
  	\underline{Stratégie et durée}: 15 min
  	\begin{center}
  		\underline{\textbf{Éléments de réponses}}
  	\end{center}
  	(Voir cahiers de recherche.)
  	
  	\begin{center}
  		\section*{II- RÉALISATION}
  	\end{center}
\sequence{NOMBRES REELS}
  	\act{ Notion de nombres réels}
  	
  	
  	\hspace{0.5cm} Joseph, un voisin et ami de Finagnon, a reçu à l'issue du recasement un terrain carré de 2500 $m^{2}$ d'aire. Les deux amis se préoccupent de savoir les dimensions des parcelles qui leur ont été attribuées. Dowa se propose de les aider à cet effet. 
  	
  	\conAct{Découverte}
  	\begin{enumerate}
  		\item Détermine la longueur $C_{1}$ du côté du terrain de Joseph.
  		\item Détermine l'aire $\mathcal{A}$ du domaine initiale de Finagnon.
  		\item Détermine l'aire $\mathcal{A'}$ du terrain qu'à reçu  Finagnon après le lotissement.
  		\item Écris la relation qui lie la longueur $C$ du côté de la nouvelle parcelle de Finagnon et $\mathcal{A'}$ 
  		\item Recopie puis remplace les pointillés par \textbf{carré} ou \textbf{positif}
  		
  		<< Puisque $C>0$ et $C^2=360 $, alors il existe un nombre ....... dont le ....... est égal à 360 >>
  		
  		\bcinfo{\textbf{Information}: Le nombre positif dont le carré est égal à 360 est noté $\sqrt{360}$ et est lu << racine carré de 360 >>. On a $\left(\sqrt{360}\right)^2=360$}
  		\item Déduis-en la longueur $C$ du côté de la nouvelle parcelle de Finagnon. 
  	\end{enumerate} 
  	\dures
  	\ret{\hspace{0.5cm} Les nombres comme $\sqrt{2}$, $\sqrt{3}$, $\sqrt{5}$ et $\sqrt{360}$, par exemple, ne sont pas des nombres rationnels. Ces nombres et tous ceux 
  		étudiés jusqu'ici (nombres rationnels) appartiennent à un 
  		ensemble appelé : « ensemble des nombres réels » et noté $\mathbb{R}$
  		et se lit  <<  grand $\mathbb{R}$>>.
  		
  		Ainsi on a la chaîne d'inclusion suivante:
  		
  		$\mathbb{N}\subset \mathbb{Z}\subset \mathbb{D}\subset \mathbb{Q}\subset \mathbb{R}$
  		
  		Les sous-ensembles de $\mathbb{R}$ sont : 
  		
  		$\mathbb{R}_{+}$ est l’ensemble des nombres réels positifs;
  		
  		$\mathbb{R}_{-}$ est l’ensemble des nombres réels négatifs;
  		
  		$\mathbb{R}^{*}$ est l’ensemble des nombres réels non nuls;
  		
  		$\mathbb{R}^{*}_{+}$ est l’ensemble des nombres réels positifs non nuls;
  		
  		$\mathbb{R}^{*}_{-}$ est l’ensemble des nombres réels négatifs non nuls. 
  	} 
  	
  	
  	\defs{On appelle « racine carrée d’un nombre réel positif $x$ », le nombre réel positif 
  		noté $\sqrt{x}$ et dont le carré est égal à $x$. << $\sqrt{x}$ >> se lit << racine carrée de $x$>>.
  		
  		Dans l’écriture << $\sqrt{x}$ >>, le symbole $\sqrt{}$  est 
  		appelé le  \textbf{radical} et le nombre réel positif $x$
  		est appelé le \textbf{ radicande}.}
  	\nbVo{Conséquence de la définition}{Si $a$ et $b$ sont des nombres réels positifs, on a :
  		
  		$\bullet$ $\sqrt{a}\geq0$ 
  		
  		$\bullet$ $ \left(\sqrt{a}\right)^2=a$ ou bien $\sqrt{a} \times \sqrt{a}=a $ 
  		
  		$\bullet$ $\sqrt{a}=b$ équivaut à $a=b^{2}$
  		
  		$\bullet$ $\sqrt{0}=0$ et $\sqrt{1}=1$}
  	
  	\fbox{Prise de notes: ........min}
  	
  	\conAct{ Construction d'un segment de longueur $\sqrt{a}$ $(a>0)$ (Escargot de Pythagore )}
  	On considère la figure suivante:
  	
  	\begin{minipage}{3cm}
  		\begin{tikzpicture}[scale=1.7]
  			%\clip(-0.5231000090689193,-0.2537795682037541) rectangle (5.121642894157984,3.4780878813309024);
  			\draw [line width=0.8pt] (1,2)-- (1,1);
  			\draw [line width=0.8pt] (1,1)-- (2,1);
  			\draw [line width=0.8pt] (2,1)-- (1,2);
  			\draw [line width=0.8pt] (1,2)-- (1.7026758571394138,2.7115101122214567);
  			\draw [line width=0.8pt] (1.7026758571394138,2.7115101122214567)-- (2,1);
  			\draw [line width=0.8pt] (1,1.1)-- (1.1,1.1);
  			\draw [line width=0.8pt] (1.1,1.1)-- (1.1,1);
  			\draw [line width=0.8pt] (1.077716992106855,1.922283007893145)-- (1.14672555859761,1.9951479086484307);
  			\draw [line width=0.8pt] (1.14672555859761,1.9951479086484307)-- (1.0714272545814891,2.0723252598002135);
  			\draw (0.7443266719051047,1.6062864404479535) node[anchor=north west] {$1$};
  			\draw (1.3780400123921166,1) node[anchor=north west] {$1$};
  			\draw (1.102257354957954,2.685946205722131) node[anchor=north west] {$1$};
  			\draw (0.6,1.1) node[anchor=north west] {$O$};
  			\draw (2,1.1) node[anchor=north west] {$A$};
  			\draw (0.7,2.1) node[anchor=north west] {$B$};
  			\draw (1.5716746442075926,2.967596579271916) node[anchor=north west] {$C$};
  		\end{tikzpicture}
  	\end{minipage} 
  	\begin{minipage}{5cm}
  		L'unité de longueur est le centimètre.
  	\end{minipage}
  	\begin{enumerate}
  		\item Calcule $AB$ et $AC$.
  		\item Reproduis et complète la figure ci-dessus jusqu'à obtenir un segment de longueur $\sqrt{5}$  
  		\item En remarquant que $41=5^2+4^2$, construire un segment de longueur $\sqrt{41}$ (Tu pourras construire un triangle rectangle dont les longueurs des côtés de l'angle droit sont 5 et 4).
  		\item En remarquant que $ 60=8^2-2^2 $, construire un segment de longueur $\sqrt{60}$ (Tu pourras construire dans un demi-cercle de diamètre 8, un triangle rectangle dont l'hypoténuse est égal au diamètre et dont l'un des côtés de l'angle droit est une corde de longueur 2)
  	\end{enumerate}	
  	\dures
  	 
  	\conAct{Puissance à exposant entier relatif}
  	\begin{enumerate}
  		\item 
  		\begin{enumerate}
  			\item  Compare $ \left(\sqrt{2}\right)^3$ et $\sqrt{2}\times \sqrt{2} \times \sqrt{2}$
  			\item  Compare $\left(\sqrt{5}\right)^4$ et $\sqrt{5}\times \sqrt{5} \times \sqrt{5}\times \sqrt{5}$
  			\item A partir de ce qui précède, complète l'égalité suivante.
  			$x^{...}= \underbrace{x\times x\times...\times x}_{n \hspace{0.2cm}facteurs}$, où $x$ est un nombre réel et $n$ un entier naturel plus grand que 1.	
  		\end{enumerate}
  		\item $a$ et $b$ sont deux nombres réels non nuls, $m$ et $n$ sont deux entiers relatifs. Complète les pointillés suivants:	
  		
  		$a^m\times a^n= a^{...}$; $(a^m)^n=a^{...}$; $a^{...}\times b^{n}= (ab)^n$; $\dfrac{a^n}{b^n}= \left(\dfrac{a}{b}\right)^{...}$;  $\dfrac{a^m}{a^n}=a^{...}$ et $a^{...}=\dfrac{1}{a^n}$.
  		\item Simplifie A et B tels que: 
  		$A=\sqrt{5^{26}}$ et $B=\sqrt{7^{23}}$
  	\end{enumerate}	
  	\dures
  	 
  	\defs{ $\bullet$ Si $x$ est un nombre réel non nul et $n$ un entier naturel plus grand que 1, on a
  		
  		$x^{n}= \underbrace{x\times x\times...\times x}_{n \hspace{0.2cm}facteurs}$ 
  		
  		$\bullet$ Si $x$  est un nombre réel non nul et $n$ un entier naturel non nul, 
  		
  		$x^{-n}=\dfrac{1}{x^n}$ et $a^n\times a^{-n}=1$. }
  	
  	\nbVo{Remarques}{ $\bullet$ $x^{0}=1$ avec $x\in \mathbb{R}^{*}$
  		
  		$\bullet$ $0^{0}$ n'est pas défini
  		
  		$\bullet$ Toutes les opérations (addition, soustraction, multiplication et division) dans $\mathbb{Q}$ sont valables dans $\mathbb{R}$ avec les mêmes propriétés.}
  	
  	\propr{ \ding{118} $x$ et $y$ sont deux nombres réels non nuls et $n$ et $m$ deux  entiers relatifs, on a:
  		
  		$\bullet$ $x^n\times y^n= (xy)^{n}$
  		
  		$\bullet$  $\dfrac{x^n}{y^n}= \left(\dfrac{x}{y}\right)^{n}$
  		
  		$\bullet$ $x^m\times x^n= x^{m+n}$
  		
  		$\bullet$	$(x^m)^n=x^{m\times n}$
  		
  		\ding{118} $x$ est un nombre réel positif et $n$ un nombre entier 
  		relatif, on a :
  		
  		$\bullet$ $\sqrt{x^{2n}}=x^n$
  		
  		$\bullet$ $\sqrt{x^{2n+1}}=x^n\sqrt{x}$. }
  	
  	\conAct{Règles de calcul sur les racines carrées}
  	La figure ci-dessous est celle de la parcelle rectangulaire reçue par un ami de Finagnon après le lotissement.
  	
  	\begin{tikzpicture}[scale=0.9]
  		\draw [line width=0.8pt] (2,2)-- (6,2);
  		\draw [line width=0.8pt] (6,2)-- (6,4);
  		\draw [line width=0.8pt] (6,4)-- (2,4);
  		\draw [line width=0.8pt] (2,4)-- (2,2);
  		\draw (0.8,3.1) node[anchor=north west] {$\sqrt{a}$};
  		\draw (3.7,4.9) node[anchor=north west] {$\sqrt{b}$};
  	\end{tikzpicture}
  	
  	$a$ et $b$ sont deux réels strictement positifs.
  	\begin{enumerate}
  		\item 
  		\begin{enumerate}
  			\item  Calcule l'aire $\mathcal{A}$ de cette parcelle.
  			\item Compare $\left(\sqrt{a} \times \sqrt{b}\right)^2$ et $\left(\sqrt{a b}\right)^2$.
  			\item  Compare $\sqrt{a} \times \sqrt{b}$ et $\sqrt{ab}$.
  		\end{enumerate}
  		\item 
  		\begin{enumerate}
  			\item  
  			Calcule $\left(\sqrt{\dfrac{a}{b}} \right)^2$ et $\left(\dfrac{\sqrt{a}}{\sqrt{b}} \right)^2$
  			\item Compare $\sqrt{\dfrac{a}{b}}$ et $\dfrac{\sqrt{a}}{\sqrt{b}}$
  		\end{enumerate}	
  		\item On suppose que $a=9$, $b=16$ $c=25$.
  		\begin{enumerate}
  			\item Calcule $\sqrt{a+b}$, $\sqrt{a}+\sqrt{b}$ puis déduis-en que  $\sqrt{a+b}\neq \sqrt{a}+\sqrt{b}$
  			\item Calcule  $\sqrt{c-b}$, $\sqrt{c}-\sqrt{b}$ puis déduis-en que  $\sqrt{c-b}\neq \sqrt{c}-\sqrt{b}$	
  		\end{enumerate}		
  	\end{enumerate}
  	\underline{Stratégie et durée}: TI... mn; TG...mn TC...mn
   	
 \propr{ 
  		
  		$P_{1}$ Si $x$ et $y$ sont des nombres réels tels que $x\geq 0$ et $y\geq 0$ alors:
  		
  		$\bullet$ $\sqrt{x} \times \sqrt{y} =\sqrt{xy}$
  		
  		$\bullet$ $\sqrt{x+y}\neq \sqrt{x}+\sqrt{y}$ 
  		
  		$P_{2}$  Si $x$ et $y$ sont des nombres réels tels que $x\geq 0$ et $y> 0$ alors:
  		
  		$\bullet$ $\sqrt{\dfrac{x}{y}}=\dfrac{\sqrt{x}}{\sqrt{y}} $.
  		
  		$P_{3}$  Si $x$ et $y$ sont des nombres réels tels que $x\geq 0$,  $y\geq 0$ et $x>y$  alors:
  		
  		$\bullet$ $\sqrt{x-y}\neq \sqrt{x}-\sqrt{y}$ }
  	
  	\conAct{Approfondissement}
  	\begin{enumerate}
  		\item Écris plus simplement
  		$A= \sqrt{25\times 121}$; $B= \sqrt{2}\times \sqrt{32}$; $C=\sqrt{\dfrac{400}{81}}\times \dfrac{\sqrt{98}}{\sqrt{2}}$ et $D= \sqrt{19-\sqrt{4+\sqrt{25}}}$
  		\item  	Écris plus simplement
  		$E=\dfrac{(x^2)^5y^{-11}}{x^{-20}y^{7}}$ et $F=\dfrac{2^7\times 5^{2}\times 3^{7}}{6^{8} \times 5}$
  	\end{enumerate}
  	\dures
  	
  	\act{ Rendre rationnel le dénominateur d'un quotient}
  	\conAct{Expression conjuguées d'un nombre réel}
  	On considère les nombres réels suivants.
  	$A=\left(\sqrt{3}-2 \right)\left(5\sqrt{3}+7 \right)$; $B=\left(\sqrt{5}-\sqrt{7} \right)\left(\sqrt{5}+\sqrt{7}\right)$; $C= \left(2\sqrt{6}-1 \right)\left(2\sqrt{6}+1 \right)$ et $D=\left(\sqrt{5}-2\sqrt{3} \right)^2$
  	\begin{enumerate}
  		\item Calcule les nombres $A$; $B$; $C$ et $D$
  		\item  Parmi les nombres calculés, donne ceux qui sont des des nombres rationnels.
  	\end{enumerate}
  	\dures
  	 
  	\ret{ De façon générale, si $a$, $b$,  $c$ et $d$ sont des nombres rationnels tels que $a$ et $b$ sont positifs, alors on a :
  		
  		$\bullet$ $\sqrt{a}-\sqrt{b}$  est une expression conjuguée de  $\sqrt{a}+\sqrt{b}$; on dit aussi que $\sqrt{a}-\sqrt{b}$ et $\sqrt{a}+\sqrt{b}$ sont des expressions conjuguées.
  		
  		$\bullet$  $c\sqrt{a}+d\sqrt{b}$ est une expression conjuguée de  $c\sqrt{a}-d\sqrt{b}$; on dit aussi que $c\sqrt{a}-d\sqrt{b}$ et $c\sqrt{a}+d\sqrt{b}$ sont des expressions conjuguées.
  		
  		$\bullet$  $b\sqrt{a}+c$ est une expression conjuguée de  $b\sqrt{a}-c$; on dit aussi que $b\sqrt{a}+c$ et $b\sqrt{a}-c$ sont des expressions conjuguées.
  	}
  	
  	\nbVo{Remarques}{ $\bullet$  Pour rendre rationnel (écris sans radical)	le dénominateur d'un nombre réel, on multiplie le numérateur et le dénominateur de ce nombre réel par une expression conjuguée du dénominateur.
  		
  		$\bullet$ Dans le cas où   le dénominateur est sous la forme $a\sqrt{b}$, on multiplie le numérateur et le dénominateur par $\sqrt{b}$. 	
  	}
  	
  	\conAct{Approfondissement}
  	
  	On pose $A=\dfrac{\sqrt{3}+\sqrt{6}}{2\sqrt{3}}$; $B=\left(-2\sqrt{3}+\sqrt{2} \right)^{-2}$ et $C= \dfrac{\sqrt{7-4\sqrt{3}}-\sqrt{7+4\sqrt{3}}}{\sqrt{7-4\sqrt{3}}+\sqrt{7+4\sqrt{3}}}$.
  	
  	Écris sans radical au dénominateur chacun des nombres $A$
  	, $B$ et $C$.
  	
  	\dures
  	
  	\act{ Comparaison des nombres réels }
  	\conAct{Propriétés des inégalités} 
  	
  	$a$, $b$ ,$c$ , $d$ et $k$ sont des nombres réels. Remplace les pointillés par le symbole qui convient.
  	\begin{enumerate}
  		\item Si $a<b$ alors $a+c ......... b+c$
  		\item Si $a<b$ et $c<d$ alors $a+c ......... b+d$
  		\item Si $a<b$ et si $k>0$ alors $ka ......... kb$
  		\item Si $a<b$ et si $k<0$ alors $ka ......... kb$
  		\item Si $a$, $b$ ,$c$ et $d$  sont des nombres réels strictement positifs.
  		
  		Si $a<b$ et $c<d$ alors $ac ......... bd$
  	\end{enumerate}
  	\dures
  	
  	\propr{\ding{118} Lorsqu'on ajoute un même nombre à chaque membre d’une inégalité, on obtient une nouvelle inégalité de même sens.
  		
  		\ding{118} Lorsqu'on additionne membres à membres des inégalités de même sens, on obtient une nouvelle inégalité de même sens. 
  		
  		\ding{118}  Lorsqu'on multiplie (ou on divise) les membres  d'une inégalité par un même nombre positif non nul, on obtient une nouvelle inégalité de même sens.
  
  \ding{118}  Lorsqu'on multiplie (ou on divise) les membres  d'une inégalité par un même nombre négatif non nul, on obtient une nouvelle inégalité de sens contraire.
  		
  		\ding{118} Lorsqu'on multiplie membres à membres des inégalités de même sens entre nombres négatifs non nuls, on obtient une nouvelle inégalité de sens contraire.  
  		
  		\ding{118} Lorsqu'on multiplie membres à membres des inégalités de même sens entre nombres positifs non nuls, on obtient une nouvelle inégalité de même sens.  
  	}
  	\conAct{Ordre dans $\mathbb{R}$}
  	\begin{enumerate}
  		\item Soient $a$ et $b$ deux nombres réels positifs tels que $a^2<b^2$
  		\begin{enumerate}
  			\item  Donne le signe de $a+b$
  			\item Justifie que $(a+b)(a-b)<0$
  			\item Déduis-en que $a<b$		
  		\end{enumerate}	
  		\item Soient $a$ et $b$ deux nombres réels négatifs tels que $a^2<b^2$
  		\begin{enumerate}
  			\item  Donne le signe de $a+b$
  			\item Justifie que $(a+b)(a-b)<0$
  			\item Déduis-en que $a>b$		
  		\end{enumerate}	
  		
  		\item Soient $a$ et $b$ deux nombres réels de même signe et différent de zéro. Justifie que si $a< b$ alors $\dfrac{1}{a}> \dfrac{1}{b}$.
  	\end{enumerate}
  	\dures
  	 
  	\propr{ \ding{118} Deux nombres positifs sont rangés dans le même ordre que leurs 
  		racines carrées. 
  		
  		\ding{118} 	Deux nombres négatifs sont rangés dans l’ordre contraire de 
  		celui de leurs carrés.
  		
  		\ding{118} Deux nombres positifs sont rangés dans le même ordre que 
  		celui de leurs carrés. 
  		
  		
  		\ding{118}  Deux nombres de même signe et différents de zéro sont rangés 
  		dans l’ordre contraire de celui de leurs inverses.
  	}
  	
  	\nbVo{Remarque}{Pour comparer deux nombres réels de même signe , on peut utiliser la démarche suivante:	
  		
  		$\bullet$ Préciser leur signe ;
  		
  		$\bullet$ Déterminer leur carré;
  		
  		$\bullet$ Comparer les carrés obtenus et déduire la comparaison de ces nombres. 
  	}
  	\conAct{Approfondissement}
  	
  	On considère les nombres réels suivants $a=3\sqrt{2}-1$ et  $b=3-2\sqrt{2}$
  	\begin{enumerate}
  		\item Justifie que $a>b$
  		\item Compare $\dfrac{1}{a}$ et $\dfrac{1}{b}$
  	\end{enumerate}
  	\dures
  	 
  	\nbVo{Remarque}{ Pour comparer deux nombres réels on peut étudier 
  		le signe de leur différence et déduire la comparaison 
  		de ces deux nombres réels.
  	}
  	\act{ Encadrement }
  	\conAct{Encadrement} 
  	Dowa veut connaitre un encadrement d’ordre deux de 
  	la longueur de la nouvelle parcelle de son père.
  	\begin{enumerate}
  		\item 
  		\begin{enumerate}
  			\item A l'aide de ta calculatrice, donne les 
  			approximations décimales d’ordre 3 de $\sqrt{10}$ (par défaut et par excès)
  			\item Déduis-en un encadrement de $\sqrt{10}$ 
  		\end{enumerate}
  		\item Utilise cet encadrement pour trouver un meilleur 
  		encadrement d’ordre 2 de la longueur de la 
  		nouvelle parcelle du père de Dowa.
  	\end{enumerate}
  	\dures
  	 
  	\ret{ $\bullet$ Pour déterminer un encadrement de $(a-b)$
  		connaissant un encadrement de $a$ et un 
  		encadrement de $b$, on détermine d’abord un 
  		encadrement de $(-b)$ ayant le même sens que 
  		celui de $a$ puis on détermine un encadrement de 
  		$a+(-b)$
  		
  		$\bullet$  Pour déterminer un encadrement du quotient $\dfrac{a}{b}$ on détermine d'abord un encadrement de $\dfrac{1}{b}$ ayant le même sens que celui de a puis on détermine un encadrement du produit $a\times \dfrac{1}{b}$.
  	}
  	\conAct{Approfondissement} 
  	On considère les nombres ci-dessous :
  	$x=7+3\sqrt{3}$ et $ y=-3\sqrt{5}+2\sqrt{3}$.
  	
  	Détermine un encadrement de $x$, $y$, $xy$
  	et $\dfrac{x}{y}$   par deux nombres décimaux d'ordre 2 sachant que 
  	$1,732< \sqrt{3} <1,733$ et $2,236 <\sqrt{5} <2,237$
  	
  	\dures
  	 \begin{center}
  		\Ovalbox{\textbf{ Exercice de maison}}
  	\end{center}
  \Ovalbox{\textbf{ Exercice 1 }}
  	Monsieur Benfa après le recasement, a érigé une maison sur un domaine triangulaire $ABC$ de dimensions $AB=\dfrac{3+2\sqrt{3}}{10} dam$; $AC=\dfrac{3\sqrt{3}-2}{10}dam$ et $BC=\dfrac{\sqrt{13}}{5}\left(\sqrt{5-2\sqrt{6}} \right)\left(\sqrt{5+2\sqrt{6}} \right)dam$.	
  	\begin{enumerate}
  		\item Etudie le signe du nombre réel  $5-2\sqrt{6}$ et celui de $3\sqrt{3}-2$  .
  		\item Justifie que  $BC=\dfrac{\sqrt{13}}{5}dam$.
  		\item Calcule $AB^2$, $AC^2$ et $BC^2$ .
  		\item Détermine la nature exacte du triangle $ABC$.
  	\end{enumerate} 
  	 \Ovalbox{\textbf{ Exercice 2 }}	Ecris plus simplement chacun des nombres réels suivants: \\
  	 $ A=\sqrt{125} $; $ B=\sqrt{32} $; $ C=\sqrt{108} $; $ D=\sqrt{252} $
  	 $ E=\sqrt{3} + \sqrt{27}+ \sqrt{12}$ ; $ F= \sqrt{3}\times \sqrt{5}\times \sqrt{15} $; $ G=\sqrt{\dfrac{324}{27}} $\\
  	 $ H=\sqrt{25-\sqrt{15+\sqrt{29-\sqrt{14+\sqrt{9-5}}}}} $ et \\
  	 $ I=\sqrt{158+\sqrt{113+\sqrt{69-\sqrt{21+\sqrt{9+7}}}}} $
   \sequence{VALEUR ABSOLUE}	
   \setcounter{q}{0}
  	\act{ Notion de valeur absolue}
  	
  	\hspace{0.5cm} Le lot de parcelle contenant celle du père de Dowa est traversé par une voie que l' architecte représente par la droite graduée ci-dessous. Il représente d'abord la parcelle de Joseph qui est ouvert sur la voie par le point $O$ puis il représente le Nord par le signe $+$ et le Sud par le signe $-$ . Les points $A$; $B$; $C$ et $D$  représentent les positions de certaines parcelles situées au bord de la voie. Dowa s'intéresse aux écarts entre certaines parcelles. 
  	\begin{center}
  		\begin{tikzpicture}[scale= 0.95]
  			%\clip(3.8941189772858755,-2.044711466329265) rectangle (15.353926892664857,5.531627238641024);
  			\draw [line width=0.8pt] (4.02620746092914,3.986896269535429)-- (13.02620746092914,3.986896269535429);
  			\draw [line width=0.8pt] (8,4.2)-- (8,3.8);
  			\draw [line width=0.8pt] (9.001447845911683,4.1709168129815355)-- (8.99878815854688,3.817178393462721);
  			\draw [line width=0.8pt] (7,4.2)-- (7,3.8);
  			\draw (7.956275421302749,3.637542562398451) node[anchor=north west] {0};
  			\draw [line width=0.8pt] (10,4.2)-- (10,3.8);
  			\draw [line width=0.8pt] (11,4.2)-- (11,3.8);
  			\draw [line width=0.8pt] (12,4.2)-- (12,3.8);
  			\draw [line width=0.8pt] (6.002778068052667,4.199200532530848)-- (6,3.8);
  			\draw (8.7,3.804317313765722) node[anchor=north west] {$15$};
  			\draw (9.8,3.8) node[anchor=north west] {$30$};
  			\draw (4.3,3.8) node[anchor=north west] {$-45$};
  			\draw (10.839096123508272,3.8281422782467605) node[anchor=north west] {$45$};
  			\draw (11.839744631711843,3.8400547604872797) node[anchor=north west] {$60$};
  			\draw (6.4,3.804317313765722) node[anchor=north west] {$-15$};
  			\draw (5.4,3.8281422782467605) node[anchor=north west] {$-30$};
  			\draw [line width=0.8pt] (10.799830665489349,4.134694574057762)-- (10.799830665489349,3.8759573752646697);
  			\draw [line width=0.8pt] (9.801872889270966,4.157695446409657)-- (9.799626895307005,3.8477482793832505);
  			\draw [line width=0.8pt] (7.200493575806389,4.138870920625325)-- (7.200124800859263,3.8662633959344097);
  			\draw [line width=0.8pt] (5.198659173295347,4.161392051888192)-- (5.2,3.85);
  			\draw [line width=0.8pt] (5,4.2)-- (5,3.8);
  			\draw [->,line width=0.8pt] (5.532607186935056,2.706895126086268) -- (5.226103996387775,3.7887929414499464);
  			\draw [->,line width=0.8pt] (7.614356066370923,2.8127467640236907) -- (7.200124800859263,3.8662633959344097);
  			\draw [->,line width=0.8pt] (9.484401669931955,2.8271929318486757) -- (9.799626895307005,3.8477482793832505);
  			\draw [->,line width=0.8pt] (10.635392677195362,2.7977125644418392) -- (10.799830665489349,3.8759573752646697);
  			\draw (5,4.7) node[anchor=north west] {$B$};
  			\draw (6.9,4.7) node[anchor=north west] {$C$};
  			\draw (7.7,4.7) node[anchor=north west] {$O$};
  			\draw (9.5,4.7) node[anchor=north west] {$A$};
  			\draw (10.5,4.7) node[anchor=north west] {$D$};
  			\draw (4.5,2.7) node[anchor=north west] {$-19\sqrt{5}$};
  			\draw (7,2.8) node[anchor=north west] {$-7\sqrt{3}$};
  			\draw (8.8,2.9) node[anchor=north west] {$10\sqrt{6}$};
  			\draw (10.6,2.8) node[anchor=north west] {$30\sqrt{2}$};
  		\end{tikzpicture}	
  	\end{center}
  	\conAct{Découverte} 
  	\begin{enumerate}
  		\item Détermine les distances $OA$ et $OC$.
  		\item  Quelle est la distance à zéro de chacun des nombres $-7\sqrt{3}$ et $10\sqrt{6}$
  		
  		\textbf{Information: La distance à zéro de $-7\sqrt{3}$ est appelée valeur absolue de $-7\sqrt{3}$ et est notée   $\left| -7\sqrt{3} \right | $. On a $\left| -7\sqrt{3} \right |= 7\sqrt{3} $ }
  		\item  Donne la valeur absolue de chacun des nombres $-19\sqrt{5}$ et $30\sqrt{2}$.
  	\end{enumerate}
  	\dures
  \end{multicols*}
  	\begin{multicols}{2}
  			\SetWatermarkText{\resizebox{28cm}{0.35cm}{Support d'activités 3\up{ème} }}
  		\SetWatermarkScale{1}
  		\SetWatermarkColor{black!30}
  		\SetWatermarkAngle{90}
  	\defs{ On appelle valeur absolue d’un nombre réel a et on note $\left| a \right |$ la distance de $a$ à zéro. 
  	 $a$ et $b$ sont des nombres réels. $A$ et $B$ sont les points d’abscisses respectives $a$ et $b$  dans un repère $\left(O; I \right)$ d’une droite $(\mathcal{D})$. 
  		\begin{center}
  			\begin{tikzpicture}[scale=0.6]
  				%	\clip(-1.14,-6.37) rectangle (18.1,6.35);
  				\draw [line width=0.8pt] (1,2)-- (11,2);
  				\draw (2.6,2.89) node[anchor=north west] {$A$};
  				\draw (7.6,2.93) node[anchor=north west] {$B$};
  				\draw (4.5,2.91) node[anchor=north west] {$O$};
  				\draw (5.5,2.93) node[anchor=north west] {$I$};
  				\draw (10.2,2.93) node[anchor=north west] {$(\mathcal{D})$};
  				\draw (2.6,2) node[anchor=north west] {$a$};
  				\draw (7.6,2) node[anchor=north west] {$b$};
  				\begin{scriptsize}
  					\draw [fill=black] (3,2) circle (2pt);
  					\draw [fill=black] (5,2) circle (2pt);
  					\draw [fill=black] (6,2) circle (2pt);
  					\draw [fill=black] (8,2) circle (2pt);
  				\end{scriptsize}
  			\end{tikzpicture}		
  		\end{center}
  		On appelle distance des nombres réels $a$ et $b$,  et on note $d(a, b)$ 
  		la distance des points $A$ et $B$. 
  	}
  	
  	\nbVo{Remarque}{On a $d(a, b)=\left| a-b \right |=\left| b-a \right |$ 
  	}
  	
  	\propr{ \ding{43}	Soit $x$ un nombre réel :
  		
  		$\bullet$ Si $x\geq0$, alors $\left| x \right |=x$
  		
  		$\bullet$ Si $x\leq 0$, alors $\left| x \right |=-x$
  		
  		\ding{43} Pour écrire $\left| x \right |$ sans le symbole $\left| \hspace{0.2cm}  \right |$, il faut déterminer d'abord le singe de $x$.
  	}
  
  	\conAct{Approfondissement}
  	\begin{enumerate}
  		\item  Détermine $ \left|-\sqrt{2} \right |$ ; $ \left|\sqrt{13} \right |$ ; $ \left|2-\sqrt{3} \right |$ et $ \left|7-\sqrt{5} \right |$.
  		\item  Calcule $d \left( 8+2\sqrt{2}; 5+4\sqrt{2} \right) $ et 
  		
  		$d \left(7+\sqrt{5}; 7-\sqrt{2} \right)$
  	\end{enumerate}
  	\dures	
  	
  	\act{ Notion d'intervalle de $\mathbb{R}$}
  	
  	\hspace{0.3cm} Pour situer les nouvelles parcelles, l'architecte en charge des travaux de recasement, utilise souvent une droite graduée  $\mathcal{(D)}$ muni d'un repère $\left(O; I \right)$ sur laquelle sont placés en des points $A$ et $B$ des piquets de référence d'abscisses respectives -1 et 4.  
  	
  	\conAct{Découverte}
  	\begin{enumerate}
  		\item  Représente la droite graduée $\mathcal{(D)}$ et place les points $A$ et $B$.
  		\item  Colorie en vert la portion de $\mathcal{(D)}$ qui est située à gauche de $A$ ($A$ non compris); en rouge la portion située à droite de $B$ ($B$ non compris) et en bleu la portion restante. 
  		\item 
  		\begin{enumerate}
  			\item Donne un encadrement  de l'abscisse $x$ d'un point appartenant à la portion en vert.
  			\item Donne un encadrement  de l'abscisse $y$ d'un point appartenant à la portion en rouge.
  			\item Donne un encadrement  de l'abscisse $z$ d'un point appartenant à la portion en bleu.
  		\end{enumerate}
  		\item Calcule $ \left|-1-(4) \right |$  
  	\end{enumerate}
  	\dures
   	
  	\ret{$\bullet$ L'écriture $-1\leq z\leq 4$ sera notée $z\in[-1;4]$. L'ensemble $[-1;4]$ est appelé intervalle de $\mathbb{R}$. On dira que c'est l'intervalle  fermé -1;4.
  		
  		\ding{51} Les réels -1 et 4 sont appelés les bornes de cet intervalle;
  		
  		\ding{51} $ \left|-1-(4) \right |$ est appelé amplitude de cet intervalle.
  		
  		$\bullet$ De même l'écriture $x<-1$ sera notée $x\in ]\leftarrow;-1 [$.  L'ensemble $]\leftarrow;-1 [$ est aussi appelé intervalle de $\mathbb{R}$. On dira que c'est l'intervalle des nombres inférieurs à -1.  
  	}
  \end{multicols}
  \nbVo{Vocabulaire sur les intervalles}{$a$ et $b$ sont deux nombres réels tels que $a<b$. 
  	
  	\begin{tabular}{|p{2cm}| p{4cm}| p{3cm}| p{2.5cm}|}
  		\hline
  		Écriture	& Lecture &Ensemble des $x$ tels que   &Représentation  \\
  		\hline
  		$[a;b]$& Intervalle fermé a, b.  &$a\leq x\leq b$  &
  		\begin{tikzpicture}[scale=0.5]
  			\draw [line width=0.4pt] (2,2)-- (6,2);
  			\draw (2.8349780818592176,2.9226010419026527) node[anchor=north west] {$a$};
  			\draw (4.839567003136388,2.9852444456925644) node[anchor=north west] {$b$};
  			\draw [line width=1pt] (3,2)-- (5,2);
  			\begin{scriptsize}
  				\draw [fill=black] (3,2) circle (2pt);
  				\draw [fill=black] (5,2) circle (2pt);
  			\end{scriptsize}
  		\end{tikzpicture}		
  		\\
  		\hline
  		$]a;b]$	&Intervalle a , b  ouvert en a, fermé en b.   &$a< x\leq b$  & 
  		\definecolor{uququq}{rgb}{0.25098039215686274,0.25098039215686274,0.25098039215686274}
  		\begin{tikzpicture}[scale=0.5]
  			
  			\draw [line width=0.4pt] (2,2)-- (6,2);
  			\draw (2.9339791936720565,2.978427990640255) node[anchor=north west] {$a$};
  			\draw (4.842753450570316,2.978427990640255) node[anchor=north west] {$b$};
  			\draw [line width=1pt] (3.0396030655952795,2)-- (5,2);
  			\begin{scriptsize}
  				\draw [color=uququq] (3.0396030655952795,2) circle (3pt);
  				\draw [fill=black] (5,2) circle (2pt);
  			\end{scriptsize}
  		\end{tikzpicture}		
  		\\
  		\hline
  		$[ a;b[ $	&Intervalle a , b fermé  en a, ouvert en b.   & $a\leq x< b$  &\definecolor{uququq}{rgb}{0.25098039215686274,0.25098039215686274,0.25098039215686274}
  		\begin{tikzpicture}[scale=0.5]
  			\draw [line width=0.4pt] (2,2)-- (6,2);
  			\draw (2.7,2.9830383329013648) node[anchor=north west] {$a$};
  			\draw (4.8400877491354635,2.9830383329013648) node[anchor=north west] {$b$};
  			\draw [line width=1pt] (3.0396030655952795,2)-- (5,2);
  			\begin{scriptsize}
  				\draw [fill=uququq] (3.0396030655952795,2) circle (1.5pt);
  				\draw [color=black] (5,2) circle (3pt);
  			\end{scriptsize}
  		\end{tikzpicture}  \\
  		\hline
  		$]a;b[ $	&Intervalle ouvert a, b  & $a< x< b$ & \definecolor{uququq}{rgb}{0.25098039215686274,0.25098039215686274,0.25098039215686274}
  		\begin{tikzpicture}[scale=0.5]
  			\draw [line width=0.4pt] (2,2)-- (6,2);
  			\draw (2.952061038507836,2.9830383329013648) node[anchor=north west] {$a$};
  			\draw (4.8400877491354635,2.9830383329013648) node[anchor=north west] {$b$};
  			\draw [line width=1pt] (3.0396030655952795,2)-- (5,2);
  			\begin{scriptsize}
  				\draw [color=uququq] (3.0396030655952795,2) circle (3pt);
  				\draw [color=black] (5,2) circle (3pt);
  			\end{scriptsize}
  		\end{tikzpicture} \\
  		\hline
  		$]\leftarrow;b[$	&Intervalle des nombres inférieurs à b  & $x<b$  & \definecolor{sqsqsq}{rgb}{0.12549019607843137,0.12549019607843137,0.12549019607843137}
  		\begin{tikzpicture}[scale=0.5]
  			\draw (4.844207934471067,2.9583633228509685) node[anchor=north west] {$b$};
  			\draw [line width=1.2pt] (3,2)-- (4.999985654692051,2.0053030455342364);
  			\draw [line width=0.4pt] (4.999985654692051,2.0053030455342364)-- (6,2);
  			\begin{scriptsize}
  				\draw [color=sqsqsq] (4.999985654692051,2.0053030455342364) circle (2.5pt);
  			\end{scriptsize}
  		\end{tikzpicture} \\
  		\hline
  		$]\leftarrow;b]$	& Intervalle des nombres inférieurs ou égaux  à b & $x\leq b$   & \definecolor{sqsqsq}{rgb}{0.12549019607843137,0.12549019607843137,0.12549019607843137}
  		\begin{tikzpicture}[scale=0.5]
  			\draw (4.844207934471067,2.9583633228509685) node[anchor=north west] {$b$};
  			\draw [line width=1.2pt] (3,2)-- (4.999985654692051,2.0053030455342364);
  			\draw [line width=0.4pt] (4.999985654692051,2.0053030455342364)-- (6,2);
  			\begin{scriptsize}
  				\draw [fill=sqsqsq] (4.999985654692051,2.0053030455342364) circle (3pt);
  			\end{scriptsize}
  		\end{tikzpicture} \\
  		\hline
  		$]a;\rightarrow[$	&Intervalle des nombres supérieurs à a  & $x>a$ &\definecolor{uququq}{rgb}{0.25098039215686274,0.25098039215686274,0.25098039215686274}
  		\begin{tikzpicture}[scale=0.5]
  			\draw (4.844207934471067,2.9583633228509685) node[anchor=north west] {$a$};
  			\draw [line width=0.4pt] (4,2)-- (5,2);
  			\draw [line width=1pt] (5,2)-- (7,2);
  			\begin{scriptsize}
  				\draw [color=uququq] (5,2) circle (3pt);
  			\end{scriptsize}
  		\end{tikzpicture}  \\
  		\hline
  		$[a;\rightarrow[$	&Intervalle des nombres supérieurs ou égaux  à a  &$x\geq a$  &\definecolor{uququq}{rgb}{0.25098039215686274,0.25098039215686274,0.25098039215686274}
  		\begin{tikzpicture}[scale=0.5]
  			\draw (4.847104735560534,2.9342433037421665) node[anchor=north west] {$a$};
  			\draw [line width=0.4pt] (4,2)-- (5,2);
  			\draw [line width=1pt] (5,2)-- (7,2);
  			\begin{scriptsize}
  				\draw [fill=uququq] (5,2) circle (2.5pt);
  			\end{scriptsize}
  		\end{tikzpicture}  \\
  		\hline
  \end{tabular}}
  \begin{multicols}{2}
  		\SetWatermarkText{\resizebox{28cm}{0.35cm}{Support d'activités 3\up{ème} }}
  	\SetWatermarkScale{1}
  	\SetWatermarkColor{black!30}
  	\SetWatermarkAngle{90}
  	\defs{ $a$ et $b$ sont deux nombres réels tels que $a<b$. Les intervalles de type $[a;b]$;	
  		$]a;b]$; $]a;b[$ et $[a;b[$ ont pour bornes $a$ et $b$ Ils ont pour amplitudes la distance des nombres $a$ et $b$, c'est-à-dire le nombre $ \left|a-b \right |$.
  	}
    
  	\conAct{Approfondissement}
  	
  	On donne les intervalles suivants $A=[-3;1]$; $B=]2;3\sqrt{2}[$ et $C=]\leftarrow;6]$
  	\begin{enumerate}
  		\item Donne trois nombres qui appartiennent à $A$
  		\item Calcule l'amplitude de chacun des intervalles $A$ et $B$.
  		\item Traduis par des inégalités, l'appartenance de $x$ à chacun des intervalles $A$; $B$ et $C$.
  	\end{enumerate}	
  	\dures
   	\conAct{Réunion et intersection d'intervalles}
  	On donne les intervalles $A$ et $B$ tels que $A=[-5;-2[$ et $B= ]-3;\sqrt{5}[$.
  	\begin{enumerate}
  		\item  Représente sur une même droite graduée $A$ et $B$.
  		\item  Donne l'intervalle qui contient les abscisses appartenant à la fois à $A$ et à $B$.
  		\item  Donne l'intervalle qui contient les abscisses appartenant à  $A$ ou à $B$. 
  	\end{enumerate}
  	\dures
   	
  	\ret{$\bullet$  L'ensemble qui contient les réels  appartenant à la fois aux intervalles   $A$ et  $B$ est $]-3;-2[$, on l'appelle intersection de $A$ et  $B$.
  		
  		$\bullet$ L'ensemble $[-5;\sqrt{5}[$ qui contient les réels  appartenant à l'intervalle $A$ ou à l'intervalle $B$ est appelé réunion de $A$ et de $B$.   
  	}
  	\nbVo{Notation}{ Soit $I$ et $J$ deux intervalles de $\mathbb{R}$.\\
  		\ding{43} La réunion de $I$ et $J$  est notée \og $I\cup J$ \fg{} et lit\og $ I \mbox{ }union \mbox{ } J $ \fg{}" . C'est l'ensemble des éléments appartenant à $I$ ou à $J$.
  		
  		\ding{43} L'intersection de $I$ et de $J$  est notée \og $I\cap J$ \fg{}   et on lit \og$ I \mbox{ }inter\mbox{ } J $\fg{}.. C'est l'ensemble des éléments appartenant à la fois à $I$ et à $J$.
  	}
  	
  	\conAct{Approfondissement}
  	
  	Représente sur une droite graduée et écris plus simplement :
  	
  	$]-8;\sqrt{3}[\cup ]1;5[$ 
  	
  	$[-3;\rightarrow[\cap ]-5;2]$ et $]-1;3]\cap[0;7]$ 
  	
  	\dures
  	(\`A faire)

  	\act{Racine carrée et valeur absolue}
  	\conAct{Propriété}
  	
  	On donne le tableau suivant: 
  	
  	\begin{tabular}{|c|c|c|c|c|c|}
  		\hline
  		$a$	& $-\sqrt{5}$  & -3  &0  & $\sqrt{3}$ & $\sqrt{11}$  \\
  		\hline
  		$\left| a \right |$	&  &  &  &  &  \\
  		\hline
  		$a^2$	&  &  &  &  &  \\
  		\hline
  		$\sqrt{a^2}$	&  &  &  &  &  \\
  		\hline
  	\end{tabular}
  	\begin{enumerate}
  		\item Complète le tableau ci-dessus.
  		\item Dis ce que tu constate dans chaque cas.
  	\end{enumerate}
  	\dures
   	
  	\propr{ La racine carrée du carré d’un nombre réel est égale à la valeur 
  		absolue de ce nombre.
  		
  		Autrement dit : pour tout $x\in \mathbb{R}$, $\sqrt{x^2}=\left| x \right |$
  	}
  	\conAct{Application}
  	1. Calcule  $\left( \sqrt{2}-3 \right)^2 $
  	2. Écrivons plus simplement $x=\sqrt{11-6\sqrt{2}}$
  	\dures 
   	\begin{center}
  		\Ovalbox{\textbf{ Exercice de maison}}
  	\end{center}
  	\begin{center}
  		\Ovalbox{\textbf{ Exercice 1}}
  	\end{center}
  	\begin{enumerate}
  		\item Calcule $A=\left|4-3\sqrt{2} \right |-\left|3\sqrt{2}-1 \right |$.
  		\item On considère les intervalles suivants:
  		
  		$I=\left[-2;5\right[$  $J=]-4;2]$	 et $K=]1;\rightarrow[$
  		\begin{enumerate}
  			\item  A quel intervalle appartient $A$ ?
  			\item   Traduis à l'aide  d'inégalités l'appartenance de $x$ à chacun des intervalles $I$; $J$ et $K$
  			\item  Représente sur une  droite graduée et détermine $I\cap J$;  $I\cap K$; $J\cap K$; $J\cup K$ et $I\cup J$
  		\end{enumerate} 
  		\item Calcule l'amplitude de chacun des intervalles $I$ et  $J$
  	\end{enumerate} 
  	
  	\begin{center}
  		\Ovalbox{\textbf{ Exercice 2}} 
  	\end{center}
  	
  	\hspace{0.4cm} Toto possède un foulard rectangulaire sur lequel sont inscrites les dimensions suivantes: $a=\sqrt{5+2\sqrt{6}}$ et $b=\sqrt{5-2\sqrt{6}}$. Son professeur de mathématiques l'informe que  l'aire de son foulard est l'unité et que si on pose $x=(a+b)^2$ et $y=(a-b)^2$ alors les nombres réels $x$ et $y$ sont des entiers naturels. Toto s'étonne.
  	
  	Tu es invité(e) à aider Toto à vérifier les dires du professeur.
  	\begin{enumerate}
  		\item Justifie que $b$ est réellement une longueur d'un côté du foulard.
  		\item Démontre que $a$ est l'inverse de $b$.
  		\item Le professeur de Toto a-t-il raison ?
  	\end{enumerate} 
  	\sequence{TRIGONOMETRIE}
  	\setcounter{q}{0}
   	\act{ Les rapports trigonométriques d'un angle aigu dans un triangle rectangle }
  	
  	\hspace{0.4cm}Pour déterminer autrement la hauteur $BC$ du pylône, l'ingénieur fait déplacer un bâton le long d'une corde fixée au sommet du pylône et au sol. Il observe les différentes position du point de contact  $M$  du bâton et de la corde à partir du point $A$ sous l'angle $\widehat{BAC}$ comme l'indique la figure ci-dessous où le point $ N $ est le projeté orthogonal de $ M $ sur $ (AB) $:
  	
  	\begin{tikzpicture}[scale=0.7]
  		\draw [shift={(1,1)},line width=1.2pt] (0,0) -- (0:0.936790841535849) arc (0:26.56505117707799:0.936790841535849) -- cycle;
  		\draw [line width=0.8pt] (1,1)-- (9,1);
  		\draw [line width=1.6pt] (9,1)-- (9,5);
  		\draw [line width=0.8pt] (9,5)-- (1,1);
  		\draw [line width=0.8pt] (6,4.75)-- (6,1);
  		\draw (0.2,1.5) node[anchor=north west] {$A$};
  		\draw (5.1,4.2) node[anchor=north west] {$M$};
  		\draw (5.7,1) node[anchor=north west] {$N$};
  		\draw (9.1,1.3) node[anchor=north west] {$B$};
  		\draw (9,5.7) node[anchor=north west] {$C$};
  		\draw [line width=0.8pt] (6,1.2)-- (6.2,1.2);
  		\draw [line width=0.8pt] (6.2,1.2)-- (6.2,1);
  		\draw [line width=0.8pt] (9,1.2)-- (8.8,1.2);
  		\draw [line width=0.8pt] (8.8,1.2)-- (8.8,1);
  		\draw (2,1.6) node[anchor=north west] {$\alpha$};
  		\draw [->,line width=0.8pt] (4.237449365778303,4.439432727504744) -- (5.786589154064753,4.415960912530706);
  		\draw [->,line width=0.8pt] (10,4) -- (9.190002325300135,2.9607083841403807);
  		\draw [->,line width=0.8pt] (2.8291404673360754,2.819877494296156) -- (4.26092118075234,2.796405679322118);
  		\draw (2.878553686118181,5.187803496270612) node[anchor=north west] {Bâton};
  		\draw (1.5670465079679927,3.584850278531467) node[anchor=north west] {Corde};
  		\draw (9.248731408561953,4.854722308168972) node[anchor=north west] {Pylône};
  	\end{tikzpicture}		
  	
  	\conAct{Découverte}
  	On donne $ AN=4 ul $, $ AB=6ul $, $ NM=3ul $, $ BC=4,5ul $, $ AM=5ul $ et $ AC=7,5ul $
  	\begin{enumerate}
  		\item 
  		\begin{enumerate}	
  			\item  Quelle est la nature des triangles  $ANM$ et $ABC$ ?
  			\item Comment appelle-t-on les côtés $[AM]$ et $ [AC] $ respectivement  pour les triangles $ANM$ et $ ABC $?
  			\item  Dans le triangle  $ANM$, que représente les côtés $[MN]$ et $ [AN] $ pour l'angle $\alpha$ ? 
  			
  			\bcinfo{\textbf{\underline{Information}: Le côté $[MN]$ est appelé côté opposé et le coté $ [AN] $ est appelé côté adjacent à l'angle $\alpha$ dans le triangle  $ANM$.}}
  		\end{enumerate}
  		\item Calcule les rapports $\dfrac{MN}{AM}$, $\dfrac{BC}{AC}$, $\dfrac{AN}{AM}$, $\dfrac{AB}{AC}$, $\dfrac{NM}{AN}$ et $\dfrac{BC}{AB}$ . 
  		\item  Compare les rapports $\dfrac{MN}{AM}$ et $\dfrac{BC}{AC}$, $\dfrac{AN}{AM}$ et $\dfrac{AB}{AC}$ et  $\dfrac{NM}{AN}$ et $\dfrac{BC}{AB}$. 
  	\end{enumerate}
  	\dures 
  	 \defs{ \hspace{0.4cm} Soit un angle aigu $\widehat{XOY}$, $M$ un point de $[OX)$ distinct de $O$, $N$ le projeté orthogonal de $M$ sur le support de la demi-droite $[OY)$, les rapports $\dfrac{ON}{OM}$et $\dfrac{MN}{OM}$ ne dépendent pas de la position du point $M$ sur $[OX)$.\\
  		\hspace*{-0.2cm} 	\begin{minipage}{4cm}
  			\begin{tikzpicture}[scale=0.3]
  				\draw [shift={(1,1)},line width=1.2pt] (0,0) -- (0:1.2918623697742297) arc (0:26.56505117707799:1.2918623697742297) -- cycle;
  				\draw [line width=0.8pt] (1,1)-- (9,1);
  				\draw [line width=0.8pt] (9,5)-- (1,1);
  				\draw (-0.1,1.4) node[anchor=north west] {$O$};
  				\draw (5.3,5) node[anchor=north west] {$M$};
  				\draw (5.6,1) node[anchor=north west] {$N$};
  				\draw (8.4,1.8) node[anchor=north west] {Y};
  				\draw (7.8,6.8) node[anchor=north west] {$X$};
  				\draw [line width=0.8pt] (6,1.2)-- (6.2,1.2);
  				\draw [line width=0.8pt] (6.2,1.2)-- (6.2,1);
  				\draw [line width=0.8pt] (6,1)-- (6,3.5);
  			\end{tikzpicture}
  		\end{minipage}
  		\hspace*{-1.2cm}  \begin{minipage}{5cm}
  			Ces rapports sont aussi les mêmes si on prend $M$
  			distinct de $O$ sur $[OY)$, $N$ étant le projeté orthogonal de 
  			$M$ sur le support de $[OX)$.  Ils ne dépendent que de l' angle  $\widehat{XOY}$ choisi.
  		\end{minipage}\\ 
  		
  		\par  Dans le triangle $OMN$ rectangle en $N$, on a 
  		
  		\ding{43} le rapport $\dfrac{ON}{OM}$ est appelé \og le cosinus de l'angle$\widehat{XOY}$\fg{} et noté $cos\widehat{XOY}$;
  		
  		\ding{43} le rapport $\dfrac{MN}{OM}$ est appelé \og le sinus de l'angle$\widehat{XOY}$\fg{} et noté $sin\widehat{XOY}$;
  		
  		\ding{43} Si l'angle $\widehat{XOY}$ n'est pas droit  le rapport $\dfrac{MN}{ON}$ est appelé \og le tangente de l'angle$\widehat{XOY}$\fg{} et noté $tan\widehat{XOY}$;
  		
  		\ding{43} Si l'angle $\widehat{XOY}$ n'est pas nul  le rapport $\dfrac{ON}{MN}$ est appelé \og le cotangente de l'angle $\widehat{XOY}$\fg{} et noté $cotan\widehat{XOY}$;
  		
  		\ding{43} Les rapports $\dfrac{ON}{OM}$; $\dfrac{MN}{OM}$;  $\dfrac{MN}{ON}$ et $\dfrac{ON}{MN}$ sont appelés \og les rapports trigonométriques de l'angle $\widehat{XOY}$ \fg{}
  	}
  	\nbVo{Notation}{Si $\alpha$ est la mesure en degré d’un angle $\widehat{XOY}$ les valeurs $sin\widehat{XOY}$; $cos\widehat{XOY}$; $tan\widehat{XOY}$ et $cotan\widehat{XOY}$ pourront être notées respectivement $sin\alpha$; $cos\alpha$; $tan\alpha$ et $cotan\alpha$.
  	}
  
  	\conAct{Approfondissement}
  	
  	On considère un triangle $EFG$ rectangle en $F$ tels que $EG=10 cm$ ; $EF=5\sqrt{3} cm$ et $FG=5cm$.
  	Calcule les rapports trigonométriques de l'angle $\widehat{EGF}$.
  	\\
  	\dures 
  	
   	\act{ Rapports trigonométriques des angles particuliers}\vspace{0.2cm}
  	\conAct{Rapports trigonométrique d'un angle de $ 0^\circ  $ et de $ 90^\circ $}
  	\begin{enumerate}
  		\item $ [OX)  $ et $ [OY) $ sont deux demi-droites confondues. $ M $ et $ N $ sont deux points tels que $ N\in [OX)  $ et $ M\in [OY) $, $ N $ est comme le projeté ortogonale de $ M $ sur $ [OX) $ donc   $ OM=ON $ et $ MN=0 $
  		\begin{enumerate}
  			\item Fais une figure.
  			\item Quelle est la mesure de l'angle $ \widehat{MON} $?
  			\item Détermine si possible les rapports trigonométriques  de $ 0^\circ $.
  		\end{enumerate}
  		\item $ [OX)  $ et $ [OY) $ sont deux demi-droites à supports perpendiculaires. $ I $ et $ J $ sont deux points tels que $I\in [OX)  $ et J et O sont confondus. $ IJ=OI$ et $ OJ=0 $
  		\begin{enumerate}
  			\item Fais une figure.
  			\item Quelle est la mesure de l'angle $ \widehat{IOY} $?
  			\item Détermine si possible les rapports trigonométriques  de $ 90^\circ $.
  		\end{enumerate}
  	\end{enumerate}
  	\dures
  	
  	 \conAct{Rapports trigonométriques des angles de $ 30^\circ $, $ 45^\circ $ et $ 60^\circ $}
  	\begin{center}
  		\definecolor{sqsqsq}{rgb}{0.12549019607843137,0.12549019607843137,0.12549019607843137}
  		\begin{tikzpicture}[scale=0.6]
  			\draw [shift={(1,1)},line width=0.8pt] (0,0) -- (0:0.5790738842721277) arc (0:60:0.5790738842721277) -- cycle;
  			\draw [shift={(6,1)},line width=0.8pt,color=sqsqsq] (0,0) -- (119.11494603851884:0.5790738842721277) arc (119.11494603851884:179.11494603851884:0.5790738842721277) -- cycle;
  			\draw [shift={(3.5,5.278808241166945)},line width=0.8pt,color=sqsqsq] (0,0) -- (-90:0.9265182148354043) arc (-90:-59.70333069373544:0.9265182148354043) -- cycle;
  			\draw [shift={(12,1)},line width=0.8pt] (0,0) -- (135:0.5790738842721277) arc (135:180:0.5790738842721277) -- cycle;
  			\draw [shift={(8,5)},line width=0.8pt] (0,0) -- (-90:0.5790738842721277) arc (-90:-45:0.5790738842721277) -- cycle;
  			\draw [line width=0.8pt] (6,1)-- (3.5,5.278808241166945);
  			\draw [line width=0.8pt] (3.5,5.278808241166945)-- (1,1);
  			\draw [line width=0.8pt] (3.5,5.278808241166945)-- (3.5,1);
  			\draw [shift={(3.5,5.278808241166945)},line width=0.8pt,color=sqsqsq] (-90:0.9265182148354043) arc (-90:-59.70333069373544:0.9265182148354043);
  			\draw [shift={(3.5,5.278808241166945)},line width=0.8pt,color=sqsqsq] (-90:0.8107034379809788) arc (-90:-59.70333069373544:0.8107034379809788);
  			\draw [line width=0.8pt] (8,1)-- (12,1);
  			\draw [line width=0.8pt] (12,1)-- (8,5);
  			\draw [line width=0.8pt] (8,5)-- (8,1);
  			\draw [line width=0.8pt] (3.5,1.2015639250616212)-- (3.714092709721938,1.2015639250616217);
  			\draw [line width=0.8pt] (3.714092709721938,1.2015639250616217)-- (3.7174290322863373,1);
  			\draw [line width=0.8pt] (8,1.2)-- (8.2,1.2);
  			\draw [line width=0.8pt] (8.2,1.2)-- (8.2,1);
  			\draw (3,6.2) node[anchor=north west] {$A$};
  			\draw (0.1,1.5) node[anchor=north west] {$B$};
  			\draw (3,1) node[anchor=north west] {$H$};
  			\draw (6,1.5) node[anchor=north west] {$C$};
  			\draw (7.4,1) node[anchor=north west] {$E$};
  			\draw (12,1.5) node[anchor=north west] {$F$};
  			\draw (7.6,5.9) node[anchor=north west] {$G$};
  			\draw (1.5,2) node[anchor=north west] {$60\textrm{\degre}$};
  			\draw (4.4,2) node[anchor=north west] {$60\textrm{\degre}$};
  			\draw (3.5,3.8) node[anchor=north west] {$30\textrm{\degre}$};
  			\draw (8,4) node[anchor=north west] {$45\textrm{\degre}$};
  			\draw (10,1.8) node[anchor=north west] {$45\textrm{\degre}$};
  			\draw (2.9646134709480982,0.1564965061239444) node[anchor=north west] {Figure 1};
  			\draw (9.311263242570618,0.13333355075305908) node[anchor=north west] {Figure 2};
  			\draw [line width=0.8pt] (1,1)-- (3.5,1);
  			\draw [line width=0.8pt] (2.209464828100951,1.1042332991689836) -- (2.209464828100951,0.8957667008310157);
  			\draw [line width=0.8pt] (2.2905351718990485,1.1042332991689836) -- (2.2905351718990485,0.8957667008310157);
  			\draw [line width=0.8pt] (3.5,1)-- (6,1);
  			\draw [line width=0.8pt] (4.709464828100951,1.1042332991689836) -- (4.709464828100951,0.8957667008310157);
  			\draw [line width=0.8pt] (4.790535171899049,1.1042332991689836) -- (4.790535171899049,0.8957667008310157);
  		\end{tikzpicture}
  	\end{center}
  	Soit $ x $ un nombre réel tel que $x>0$, $ABC$ est un triangle équilatéral tels que $AB=x$   et $AH=\dfrac{x\sqrt{3}}{2}$ et
  	$EFG$ est triangle rectangle isocèle en $E$ tels que $EF=x$ et $FG=x\sqrt{2}$
  	
  	\begin{enumerate}
  		\item En considérant la figure 1, détermine:
  		\begin{enumerate}
  			\item   les rapports trigonométriques de $30\textrm{\degre}$.
  			\item    les rapports trigonométriques de $60\textrm{\degre}$.
  		\end{enumerate}	 
  		\item En considérant la figure 2, détermine les rapports trigonométriques de $45\textrm{\degre}$.
  		\item  Reproduis puis complète le tableau suivant:
  		
  		\begin{tabular}{|c|c|c|c|c|c|}
  			\hline
  			Mesures $\alpha$	&$0\textrm{\degre}$  &$30\textrm{\degre}$  & $45\textrm{\degre}$ &$60\textrm{\degre}$  & $90\textrm{\degre}$ \\
  			\hline
  			$sin\alpha$	&  &  &  &  &  \\
  			\hline
  			$cos\alpha$	&  &  &  &  &  \\
  			\hline
  			$tan\alpha$	&  &  &  &  &   \\
  			\hline
  			$cotan\alpha$	&   &  &  &  &  \\
  			\hline
  		\end{tabular}
  	\end{enumerate}
  	\dures
  	 
  	\ret{ 
  		\begin{tabular}{|c|p{1cm}|c|c|c|p{1cm}|}
  			\hline
  			Mesures $\alpha$	&$0\textrm{\degre}$  &$30\textrm{\degre}$  & $45\textrm{\degre}$ &$60\textrm{\degre}$  & $90\textrm{\degre}$ \\
  			\hline
  			$sin\alpha$	& 0 &$\dfrac{1}{2}$  & $\dfrac{\sqrt{2}}{2}$ & $\dfrac{\sqrt{3}}{2}$ &1  \\
  			\hline
  			$cos\alpha$	& 1 &$\dfrac{\sqrt{3}}{2}$  &$\dfrac{\sqrt{2}}{2}$  & $\dfrac{1}{2}$ & 0 \\
  			\hline
  			$tan\alpha$	&0  & $\dfrac{\sqrt{3}}{3}$ & 1 &$\sqrt{3}$  &n'existe pas \\
  			\hline
  			$cotan\alpha$	&n'existe pas  &$\sqrt{3}$  & 1 & $\dfrac{\sqrt{3}}{3}$ & 0 \\
  			\hline
  		\end{tabular}
  		{\textbf{Remarque:}}{La tangente d'un angle droit et la cotangente d'un angle nul ne sont pas définies} 
  	}
  
  	\conAct{Approfondissement}
  	
  	Pour déterminer la hauteur $BC$ du pylône, l'ingénieur fixe les valeurs de $\alpha$ et de $AB$ dans le  triangle $ABC$ rectangle en $B$ de l'activité 1 tels que  $\widehat{BAC}=30\textrm{\degre}$  et $AB=5\sqrt{3}m$. 
  	\begin{enumerate}
  		\item Déterminer la hauteur $BC$
  		\item Calcule $AC$
  	\end{enumerate}
  	\dures
  	 
  	\act{ Propriétés} 
  	On considère le tableau des rapports trigonométriques de la consigne précédente. \\
  	\conAct{Propriétés} 
  	\begin{enumerate}
  		\item Calcule $ (cos(30^\circ))^2+ (sin(30^\circ))^2 $
  		\item \begin{enumerate} \item Calcule $ \dfrac{sin(30^\circ)}{cos(30^\circ)}  $ et donne la valeur de $ tan(30^\circ ) $ 
  			\item Compare $ \dfrac{sin(30^\circ)}{cos(30^\circ)}   $ et $ tan(30^\circ ) $.
  			\item Complète : \og Soit $ \alpha $ de mesure inférieure à $ 90^\circ $, si $ cos(\alpha)\neq \dots  $ alors $ \dfrac{sin(\alpha)}{cos(\alpha)} =tan(\dotsc ) $ \fg{}
  		\end{enumerate}
  		\item \begin{enumerate} \item Calcule $ \dfrac{cos(30^\circ)}{sin(30^\circ)}  $ et donne la valeur de $ cotan(30^\circ ) $ 
  			\item Compare $ \dfrac{cos(30^\circ)}{sin(30^\circ)}   $ et $ cotan(30^\circ ) $.
  			\item Complète : \og Soit $ \alpha $ de mesure inférieure à $ 90^\circ $, si $ sin(\alpha)\neq \dots  $ alors $ \dfrac{cos(\alpha)}{sin(\alpha)} =cotan(\dotsc ) $ \fg{}
  		\end{enumerate}
  	\end{enumerate}
  	\dures
  	 
  	\propr{ Pour tout angle non nul $\alpha$ de mesure inférieure à $90\textrm{\degre}$  on a: 
  		
  		$\bullet$  $0<sin\alpha<1$ et $0<cos\alpha<1$
  		
  		$\bullet$ $sin^2\alpha+ cos^2\alpha=1$\vspace{0.2cm}
  		
  		$\bullet$ $tan\alpha=\dfrac{sin\alpha}{cos\alpha}$ si $cos\alpha\neq 0$ \vspace{0.2cm}
  		
  		$\bullet$ $cotan\alpha=\dfrac{cos\alpha}{sin\alpha}$ si $sin\alpha\neq 0$
  	}
  	\conAct{Propriétés}
  	\begin{enumerate}
  		\item Que peut-on dire d'un angle de mesure de $ 30^\circ  $ et de mesure $ 60^\circ $ ?
  		\item Compare $cos(30^\circ) $ et $sin(60^\circ)  $ d'une part et $cos(60^\circ) $ et $sin(30^\circ)  $ d'autre part
  		\item Complète la phrase suivante pour en faire une propriété : \og Pour deux angles aigus $ a $ et $ b $ complémentaires on a :$cos(a) =sin(\dots)  $  et $cos( \dots ) =sin(b)  $. \fg{}
  		\item Compare $cotan(30^\circ) $ et $tan(60^\circ)  $ d'une part et $cotan(60^\circ) $ et $tan(30^\circ)  $ d'autre part
  		\item Complète la phrase suivante pour en faire une propriété : \og Pour deux angles aigus $ a $ et $ b $ complémentaires on a :$cotan(a) =tan(\dots)  $  et $cotan( \dots ) =tan(b)    \fg{}$
  		\item Soient $ \alpha $ et $ \beta $ deux angles aigus tel que $ sin(\alpha) =0,5$ et $ sin(\beta)=0,5 $. A l'aide de ta calculatrice, détermine les mesure en degré de $ \alpha
  		$ et $ \beta $ puis déduisant que $ \alpha=\beta $
  	\end{enumerate}
  	\dures
  	 
  	\propr{ Pour deux angles aigus $\alpha$ et $\beta$ complémentaires on a : 
  		
  		$\bullet$  $sin\alpha=cos\beta$ et $cos\alpha=sin\beta$
  		
  		$\bullet$ $tan\alpha=cotan\beta$ et $cotan\alpha=tan\beta$
  		
  		$\bullet$	Lorsque deux angles aigus ont un même rapport trigonométrique, 
  		ils ont la même mesure.
  	}
   
  	\conAct{Approfondissement}
  	
  	ABC est un triangle rectangle tel que $cos\widehat{B}=\dfrac{1}{4}$. Calcule  $sin\widehat{B}$ ; $tan\widehat{B}$ et $cotan\widehat{B}$ puis déduis-en une valeur approchée de la mesure de l'angle $\widehat{B}$ à $10^{-1}$ près.
  	\\
  	\dures

  	\begin{center}
  		\Ovalbox{\textbf{ Exercice de maison}}
  	\end{center}
  	\begin{center}
  		\Ovalbox{\textbf{ Exercice 1}}
  	\end{center}
  	
  	On considère la figure ci-dessous réalisé par l'architecte où $sin\widehat{BCA}=\dfrac{\sqrt{5}}{3}$ et $AC=\sqrt{5}$.
  	
  	\begin{tikzpicture}[scale=0.6]
  		\draw [line width=0.8pt] (1,1)-- (8,1);
  		\draw [line width=0.8pt] (8,1)-- (8,5);
  		\draw [line width=0.8pt] (8,5)-- (1,1);
  		\draw [line width=0.8pt] (8,1.5)-- (7.5,1.5);
  		\draw [line width=0.8pt] (7.5,1.5)-- (7.5,1);
  		\draw (7.7,6) node[anchor=north west] {$A$};
  		\draw (0.1,1.7) node[anchor=north west] {$C$};
  		\draw (8.1,1.7) node[anchor=north west] {$B$};
  		\draw [line width=0.8pt] (5.831116100010794,3.760637771434739)-- (6.049128785848834,3.3562416022560373);
  		\draw [line width=0.8pt] (6.049128785848834,3.3562416022560373)-- (6.485806199417089,3.6406271246182906);
  		\draw [line width=0.8pt] (6.272954819396147,4.013117039654941)-- (7.994735984913256,1);
  		\draw (5.3,5) node[anchor=north west] {$H$};
  	\end{tikzpicture}	
  	
  	\begin{enumerate}
  		\item Calcule $cos\widehat{BCA}$.
  		\item Détermine $AB$ et $CH$.
  		\item  Calcule $tan\widehat{BCA}$ puis déduis-en $cotan\widehat{BAC}$. 
  	\end{enumerate} 
  	\begin{center}
  		\Ovalbox{\textbf{ Exercice 2}} 
  	\end{center}
  	
  	$ABC$ est un triangle rectangle en $A$ tel que $AB = 3\ \mathrm{cm}$ et $BC = 5\ \mathrm{cm}$.
  	\begin{enumerate}
  		\item Calcule $AC$.
  		\item Détermine les rapports trigonométriques de l'angle $\widehat{ABH}$ (ou l'angle en $B$ selon la figure).
  	\end{enumerate}
  	\begin{center}
  		\Ovalbox{\textbf{ Exercice 3}} 
  	\end{center}
  	
  	\begin{enumerate}
  		\item $ABC$ est un triangle rectangle en $B$ tel que $\widehat{BAC} = 60^\circ$ et $BC = 12\ \mathrm{cm}$. Calcule $AC$ et $ BA $( de deux différentes manières).
  		\item $ABC$ est un triangle rectangle en $B$ tel que $\tan\widehat{BAC} = \dfrac{4}{3}$ et $BC = 2\ \mathrm{cm}$. Calcule $AB$ et $AC$.
  		\item $ABH$ est un triangle rectangle en $B$ tel que $\sin\widehat{BAH} = \dfrac{1}{3}$. Calcule $\tan\widehat{BAH}$ et $cotan\widehat{BAH}$.
  	\end{enumerate}
  	\begin{center}
  		\Ovalbox{\textbf{ Exercice 4}} 
  	\end{center}
  	\hspace{0.4cm} Le cercle $\mathcal{(C)}$ de centre $O$ a pour diamètre $2\sqrt{5} m$ tels que    $NJ=\dfrac{\sqrt{10}}{2} \left(\sqrt{3}-1 \right )m $ et $sin15\textrm{\degre}=\dfrac{\sqrt{6}-\sqrt{2}}{4}$
  	
  	\begin{tikzpicture}[scale=0.5]
  		\draw [shift={(5.67277363599308,6.025022598276512)},line width=0.8pt] (0,0) -- (-0.1675815050214416:0.9130349294112965) arc (-0.1675815050214416:59.83241849497854:0.9130349294112965) -- cycle;
  		\draw (9.561850568124978,1.6748785607458994) node[anchor=north west] {J};
  		\draw (9.5,5.8) node[anchor=north west] {$O$};
  		\draw (4,6.6) node[anchor=north west] {$M$};
  		\draw (14.4,6.5) node[anchor=north west] {$N$};
  		\draw [line width=0.8pt] (10,6) circle (4.327298711180112cm);
  		\draw [line width=0.8pt] (5.67277363599308,6.025022598276512)-- (14.327298701418915,5.999709346867662);
  		\draw [line width=0.8pt] (5.67277363599308,6.025022598276512)-- (7.739687440299886,9.690054371990062);
  		\draw [line width=0.8pt] (7.739687440299886,9.690054371990062)-- (14.327298701418915,5.999709346867662);
  		\draw [line width=0.8pt] (7.739687440299886,9.690054371990062)-- (9.629595423533942,1.6885832275788903);
  		\draw [line width=0.8pt] (5.67277363599308,6.025022598276512)-- (9.629595423533942,1.6885832275788903);
  		\draw [line width=0.8pt] (9.629595423533942,1.6885832275788903)-- (14.327298701418915,5.999709346867662);
  		\draw (6.6,7.3) node[anchor=north west] {$60\textrm{\degre}$};
  		\draw (7,10.8) node[anchor=north west] {$Q$};
  		\draw (7.7,6) node[anchor=north west] {$R$};
  		\begin{scriptsize}
  			\draw [fill=black] (10,6) circle (3pt);
  		\end{scriptsize}
  	\end{tikzpicture}
  	
  	\begin{enumerate}
  		\item Justifie que $MNQ$ et $MNJ$ sont des triangles rectangles.
  		\item Prouve que $MQ=\sqrt{5}m$ et $NQ=\sqrt{15}m$.
  		\item Justifie que $cos75\textrm{\degre}=\dfrac{\sqrt{6}-\sqrt{2}}{4}$.
  		\item Calcule $cos\widehat{MNJ}$ puis déduis-en la mesure de l'angle $\widehat{MNJ}$
  	\end{enumerate} 
  	\sequence{PROPRIETES DE THALES RELATIVES AU TRIANGLE}
  	\setcounter{q}{0}
   	\act{ Propriété de Thalès et sa réciproque}
  	
  	
  	\hspace{0.5cm}  Dans le jardin public, le conseil municipal décide de mettre deux types de gazon sur un espace triangulaire $ABC$. Cet espace sera scindé en deux parties: la première est un triangle $AMN$ réservé pour le  premier type de gazon et la seconde est un trapèze réservé pour le deuxième type de gazon. On souhaite que le rapport $\dfrac{AM}{AB}$ soit égal à $\dfrac{AN}{AC}$. L'ingénieur réalise la figure suivante:
  	\begin{center}
  		\begin{tikzpicture}[scale=0.8]
  			\draw [line width=0.8pt] (2,1)-- (8,1);
  			\draw [line width=0.8pt] (8,1)-- (4.5,4.5);
  			\draw [line width=0.8pt] (4.5,4.5)-- (2,1);
  			\draw [line width=0.8pt] (2.5,2.5)-- (7,2.5);
  			\draw [line width=0.8pt] (4.5,1.2)-- (4.7,1.2);
  			\draw [line width=0.8pt] (4.7,1.2)-- (4.7,1);
  			\draw [line width=0.8pt] (4.5,2.7)-- (4.7,2.7);
  			\draw [line width=0.8pt] (4.7,2.7)-- (4.7,2.5);
  			\draw [line width=0.8pt] (4.5,4.5)-- (4.5,1);
  			\draw (4.2,5.2) node[anchor=north west] {$A$};
  			\draw (2.4,3.1) node[anchor=north west] {$M$};
  			\draw (1.3,1.5) node[anchor=north west] {$B$};
  			\draw (8,1.4) node[anchor=north west] {$C$};
  			\draw (6.3,3.1) node[anchor=north west] {$N$};
  			\draw (3.9,3.1) node[anchor=north west] {$I$};
  			\draw (4.2,1) node[anchor=north west] {$J$};
  		\end{tikzpicture}
  	\end{center}
  	\conAct{Découverte}
  	\begin{enumerate}
  		\item
  		\begin{enumerate}
  			\item Justifie que  $mes\widehat{AMI}=mes\widehat{ABJ}$.  
  			\item 	En considérant les triangles rectangles $AMI$ et $ABJ$, justifie que $\dfrac{AI}{AM}=\dfrac{AJ}{AB}$ puis déduis-en que $\dfrac{AM}{AB}=\dfrac{AI}{AJ}$
  		\end{enumerate}	
  		\item 
  		\begin{enumerate}
  			\item Justifie que  $mes\widehat{ANI}=mes\widehat{ACJ}$. 
  			\item En considérant les triangles rectangles $ANI$ et $ACJ$, justifie que $\dfrac{AI}{AN}=\dfrac{AJ}{AC}$ puis déduis-en que $\dfrac{AN}{AC}=\dfrac{AI}{AJ}$
  		\end{enumerate}
  		\item  En déduis que $\dfrac{AM}{AB}=\dfrac{AN}{AC}$
  		
  	\end{enumerate}	
  	\dures
   	
  	\propr{{Propriété de Thalès}\\ $ABC$ est un triangle. Si un point $M$ de la droite $(AB)$ et un point $N$ de la droite $(AC)$ sont tels que les droites $(MN)$ et $(BC)$ sont parallèles, alors 	on a $\dfrac{AM}{AB}=\dfrac{AN}{AC}$.
  	}
  	\fbox{Prise de notes: ........min}
  	
  	\conAct{Approfondissement}
  	
  	$IJK$ est un triangle tel que  $IJ = 6m$ et $IK = 8m$. $M$ est un point de $(IK)$ tel que  $IM = 5m$. La parallèle à $(JK)$ passant  par $M$, coupe $(IJ)$ en un point $N$.  Calcule $IN$. (Voir figure) 
  	\begin{tikzpicture}[scale=0.7]
  		\draw [line width=0.8pt] (3,3)-- (8,3);
  		\draw [line width=0.8pt] (8,3)-- (6.61,6.23);
  		\draw [line width=0.8pt] (6.61,6.23)-- (3,3);
  		\draw [line width=0.8pt] (4.805,4.615)-- (7.305,4.615);
  		\draw (8.09,3.55) node[anchor=north west] {$K$};
  		\draw (6.45,7.11) node[anchor=north west] {$I$};
  		\draw (2.47,3.53) node[anchor=north west] {$J$};
  		\draw (7.4,5.15) node[anchor=north west] {$M$};
  		\draw (4,5.27) node[anchor=north west] {$N$};
  	\end{tikzpicture}
  	
  	\dures
  	 
  	\conAct{Réciproque de la propriété de Thalès}\\
  	On considère un triangle $ABC$ tels que $AB = 3cm$ et $AC= 6cm$. $M$ et $N$ sont deux points tels que $M\in[AB] $et $N\in [AC]$ avec $AM=1cm$ et $AN=2cm$
  	\begin{enumerate}
  		\item  Fais une figure.
  		\item  Justifie que $\dfrac{AM}{AB}=\dfrac{AN}{AC}$  
  		\item  La droite parallèle à $(BC)$ passe par $M$, coupe $(AC)$ en un point $N’$
  		\begin{enumerate}
  			\item  Marque le point $N’$ sur la figure.
  			\item  Justifie que $N’$  est confondu à $N$ et déduis-en que les droites $(MN)$ et $(BC)$ sont parallèles.
  		\end{enumerate}
  	\end{enumerate}
  	\dures
   	
  	\begin{bclogo}[couleur = lightgray!15, arrondi = 0.3, ombre = true, epOmbre =0.25, couleurOmbre= lightgray!90, logo=\bcrecyclage, epBarre=3  ]{\textbf{Réciproque de la propriété de Thalès}}
  		$ABC$ est un triangle, $M $ un point de la droite $(AB)$ et $N$ un point de la droite $(AC)$ tels que la position qu’occupe $ M$ par rapport à $A$ et $B$ est la même que 
  		celle qu’occupe $N$ par rapport à $A$ et $C$. 
  		
  		Si l’on a $\dfrac{AM}{AB}=\dfrac{AN}{AC}$ alors la droite $(MN)$ si elle existe est parallèle à la droite $(BC)$.	
  	\end{bclogo}
  	
  	\fcolorbox{yellow}{pink!10}{\begin{minipage}{1\linewidth}
  			\begin{tikzpicture}[scale=0.4]
  				\draw [line width=0.8pt] (3,3)-- (8,3);
  				\draw [line width=0.8pt] (8,3)-- (6.61,6.23);
  				\draw [line width=0.8pt] (6.61,6.23)-- (3,3);
  				\draw [line width=0.8pt] (4.805,4.615)-- (7.305,4.615);
  				\draw (7.6,3) node[anchor=north west] {$C$};
  				\draw (6,7.4) node[anchor=north west] {$A$};
  				\draw (2.1,3) node[anchor=north west] {$B$};
  				\draw (3.3,5.5) node[anchor=north west] {$M$};
  				\draw (7.3,5.5) node[anchor=north west] {$N$};
  				\draw [line width=0.8pt] (10.020235811891064,2.8333138102731756)-- (14.802031257293034,2.782979121374207);
  				\draw [line width=0.8pt] (14.802031257293034,2.782979121374207)-- (12.008456023400303,6.834921577741136);
  				\draw [line width=0.8pt] (12.008456023400303,6.834921577741136)-- (10.020235811891064,2.8333138102731756);
  				\draw [line width=0.8pt] (10.41589932901553,3.6296492434730525)-- (14.239181584756464,3.599364682440762);
  				\draw [line width=0.8pt] (17.20814131399237,2.84162956872794)-- (22.312452235155607,3.6097540277379414);
  				\draw [line width=0.8pt] (22.312452235155607,3.6097540277379414)-- (18.25499437960636,8.06887132321914);
  				\draw [line width=0.8pt] (17.20814131399237,2.84162956872794)-- (21.296228359468206,8.545078779140571);
  				\draw [line width=0.8pt] (18.25499437960636,8.06887132321914)-- (21.296228359468206,8.545078779140571);
  				\draw (11.4,8.2) node[anchor=north west] {$A$};
  				\draw (9.3,4.7) node[anchor=north west] {$B$};
  				\draw (14.4,4.5) node[anchor=north west] {$C$};
  				\draw (9.3,2.8) node[anchor=north west] {$M$};
  				\draw (14.4,3) node[anchor=north west] {$N$};
  				\draw (20,7.4) node[anchor=north west] {$A$};
  				\draw (21.7,3.6) node[anchor=north west] {$C$};
  				\draw (16.6,2.8) node[anchor=north west] {$B$};
  				\draw (21.4,9.3) node[anchor=north west] {$M$};
  				\draw (16.9,8.7) node[anchor=north west] {$N$};
  			\end{tikzpicture}
  	\end{minipage} } 
  
  	\act{Conséquence de la propriété de Thalès}
  	\conAct{Conséquence de la propriété de Thalès}
  	On considère la figure suivante:
  	\begin{center}
  		\begin{tikzpicture}[scale=0.8]
  			\draw (0.4,1.3) node[anchor=north west] {$B$};
  			\draw (5,1.1) node[anchor=north west] {$C$};
  			\draw (2.6,3.6) node[anchor=north west] {$A$};
  			\draw (1.3,2.5) node[anchor=north west] {$M$};
  			\draw (3.4,2) node[anchor=north west] {$N$};
  			\draw [line width=0.8pt] (1,1)-- (5,1);
  			\draw [line width=0.8pt,dash pattern=on 1pt off 1pt] (5,1)-- (7,3);
  			\draw [line width=1.2pt,dash pattern=on 1pt off 1pt] (7,3)-- (3,3);
  			\draw [line width=0.4pt] (3,3)-- (1,1);
  			\draw [line width=0.8pt] (3,3)-- (5,1);
  			\draw [line width=0.8pt] (1.9958966326793492,1.9958966326793492)-- (4.004103367320651,1.9958966326793492);
  			\draw [line width=1.2pt,dash pattern=on 1pt off 1pt] (5.010914580490043,3)-- (4.004103367320651,1.9958966326793492);
  			\draw (7,3.4) node[anchor=north west] {$D$};
  			\draw (4.7,3.7) node[anchor=north west] {$E$};
  		\end{tikzpicture}	
  	\end{center}
  	$ABCD$ et $AMNE$ sont des parallélogrammes tels que $(MN) \parallel (BC)$ et $(NE) \parallel (CD)$.
  	\begin{enumerate}
  		\item  Justifie que $\dfrac{AM}{AB}=\dfrac{AN}{AC}$ puis que  $\dfrac{AN}{AC}=\dfrac{AE}{AD}$.
  		\item Justifie que $MN=AE$ et $BC=AD$.
  		\item En déduis que $\dfrac{AM}{AB}=\dfrac{AN}{AC}=\dfrac{MN}{BC}$
  	\end{enumerate}
  	\dures
  	 
  	\nbVo{Conséquence de la propriété de Thalès}{Dans un triangle, si une droite parallèle au support d'un côté détermine deux triangles, alors les longueurs des côtés de l'un sont 	proportionnelles à celles des côtés de l'autre.
  		\\$ABC$ est un triangle. $M \in(AB)$ et $N \in(AC)$ tels que $(MN)\parallel(BC)$ alors $\dfrac{AM}{AB}=\dfrac{AN}{AC}=\dfrac{MN}{BC}$
  		
  		\begin{tikzpicture}[scale=0.7]
  			\draw [line width=0.8pt] (3,3)-- (8,3);
  			\draw [line width=0.8pt] (8,3)-- (6.61,6.23);
  			\draw [line width=0.8pt] (6.61,6.23)-- (3,3);
  			\draw [line width=0.8pt] (4.805,4.615)-- (7.305,4.615);
  			\draw (8.09,3.55) node[anchor=north west] {$C$};
  			\draw (6.45,7.11) node[anchor=north west] {$A$};
  			\draw (2.4,3.53) node[anchor=north west] {$B$};
  			\draw (7.4,5.15) node[anchor=north west] {$N$};
  			\draw (4,5.27) node[anchor=north west] {$M$};
  		\end{tikzpicture}
  	}
   
  	
  	\conAct{Approfondissement} 
  	\begin{enumerate}
  		\item  Construis un  triangle $AMN$ tels que $AM=4,5 cm$; $AN=6cm$ et $MN=9cm$ puis place le point $B$ de $[AM]$ tel que $AB=3 cm$ et le point $C$ de $[AN]$ tel que $AC=4 cm$
  		\item Démontre que $(BC) \parallel (MN)$ 
  		\item Calcule $BC$
  		\item La parallèle à la droite $(CM)$  passe par $N$ coupe la droite $(AM)$ en $D$.
  		\begin{enumerate} 
  			\item Démontre que $\dfrac{AM}{AD}=\dfrac{AB}{AM}$
  			\item  Calcule $AD$	
  		\end{enumerate}
  	\end{enumerate}
  	\dures
   	
  	\act{ Construction de la quatrième proportionnelle à partir de trois longueurs  $x$; $y$ et $z$}
  	2; 3  et 4 sont trois nombres pris dans cet ordre. On veut construire un segment de longueur $t$ tel que $\dfrac{2}{3}=\dfrac{4}{t}$.
  	\conAct{Construction d'une quatrième proportionnelle }
  	\begin{enumerate} 
  		\item Trace deux demi-droites de même origine $A$.  Sur 
  		l’une, marque les points $B$ et $M$ tels que $AB=2$ et $AM=3$. Sur l’autre, marque le point $C$  tel que $AC=4$. 
  		\item Construis la parallèle à $(BC)$ passant par $M$
  		; elle coupe $(AC)$ en $N$. 
  		\item  Précise le segment recherché.	
  	\end{enumerate}
  	\dures
     	\propr{{Construction d'une quatrième proportionnelle}\\La quatrième proportionnelle de trois nombres réels strictement positifs $x$, $y$ et $z$ pris dans cet ordre est le nombre réel strictement positif $t$ tel que $\dfrac{x}{y}=\dfrac{z}{t}$
  	} 
  	\section*{Programme de construction}
  	%\begin{tcolorbox}[colback=violet!10,colframe=blue!40,title=Retenons]
  	On donne trois segments de longueurs $a$, $b$ et $c$. On veut construire la quatrième proportionnelle des nombres $a$, $b$ et $c$ pris dans cet ordre. Cela revient à construire un segment $[MN]$ tel que	$\dfrac{a}{b}=\dfrac{c}{MN}$.
  	
  	Pour cela :
  	
  	$\bullet$ On trace deux demi-droites de même origine $M$;
  	
  	$\bullet$ Sur l'une des demi-droites, on marque les points $A$ et $B$ tel que $MA=a$ et $MB = b$;
  	
  	$\bullet$ Sur l'autre demi-droite, on marque le point $C$ tel que $MC=c$;
  	
  	$\bullet$ On trace la droite parallèle à $(AC)$ qui passe par $B$ ; elle coupe $(MC)$ en un point $N$.
  	
  	D’après la propriété de Thalès on a:
  	
  	$\dfrac{MA}{MB}=\dfrac{MC}{MN}$;  $\dfrac{a}{b}=\dfrac{c}{MN}$. Donc $MN$ est la quatrième proportionnelle des nombres $a$, $b$ et $c$  pris dans cet ordre. 
  	
  	\begin{tikzpicture}[scale=0.8]
  		\draw (1,1.3) node[anchor=north west] {$M$};
  		\draw (3.8,3.5) node[anchor=north west] {$N$};
  		\draw (5.4,0.6) node[anchor=north west] {$A$};
  		\draw (3,2.8) node[anchor=north west] {$C$};
  		\draw (7.5,0.5) node[anchor=north west] {$B$};
  		\draw [line width=0.8pt] (2,1)-- (5.987649522018693,0.6859119716638544);
  		\draw [line width=0.8pt] (3.6704250654847383,2.099854581569904)-- (5.987649522018693,0.6859119716638544);
  		\draw [line width=0.8pt] (4.505667303223294,2.6497367564425858)-- (7.980636208267881,0.5289048382554324);
  		\draw [line width=0.8pt] (4.505667303223294,2.6497367564425858)-- (5.5576110228806925,3.3424481549368483);
  		\draw [line width=0.8pt] (7.980636208267881,0.5289048382554324)-- (8.888330948855394,0.45740565560315527);
  		\draw [line width=0.8pt] (5.987649522018693,0.6859119716638544)-- (7.980636208267881,0.5289048382554324);
  		\draw [line width=0.8pt] (2,1)-- (3.6704250654847383,2.099854581569904);
  		\draw [line width=0.8pt] (3.6704250654847383,2.099854581569904)-- (4.5056673032232935,2.649736756442586);
  	\end{tikzpicture}
   
  	\begin{center}
  		\Ovalbox{\textbf{ Exercice de maison}}
  	\end{center}
  	\begin{center}
  		\Ovalbox{\textbf{ Exercice 1}}
  	\end{center}
  	On considère un triangle $ABH$ rectangle en $H$ tels que $AH=3cm$  et $BH=4cm$. Soit $M$ un point du segment $[BH]$ tel que $BM=2,4cm$. La perpendiculaire en $M$ à $(HB)$ coupe la droite $(AB)$ en $N$.
  	
  	
  	\begin{enumerate}
  		\item Fais un figure.
  		\item  Démontre que $(MN) \parallel (AH)$.
  		\item  Calcule $AB$ et $BN$. Déduis-en la valeur de $AN$.
  		\item Calcule alors $MN$ (en précisant la propriété utilisée).
  	\end{enumerate}
  	\begin{center}
  		\Ovalbox{\textbf{ Exercice 2}}
  	\end{center}
  	Le levé topographique de la nouvelle parcelle de Estelle donne le dessin ci-dessous:
  	
  	\begin{tikzpicture}[scale=0.6]
  		\draw [line width=0.8pt] (1,1)-- (7,1);
  		\draw [line width=0.8pt] (7,1)-- (7,4);
  		\draw [line width=0.8pt] (7,4)-- (1,4);
  		\draw [line width=0.8pt] (1,4)-- (1,1);
  		\draw [line width=0.8pt] (1,4)-- (7,1);
  		\draw [line width=0.8pt] (3.87,4)-- (9.858,1.006);
  		\draw [line width=0.8pt] (7,1)-- (9.858,1.006);
  		\draw [line width=0.8pt] (1,1.4)-- (1.4,1.4);
  		\draw [line width=0.8pt] (1.4,1.4)-- (1.4,1);
  		\draw [line width=0.8pt] (1.4,4)-- (1.4,3.6);
  		\draw [line width=0.8pt] (1.4,3.6)-- (1,3.6);
  		\draw [line width=0.8pt] (6.6,1)-- (6.6,1.4);
  		\draw [line width=0.8pt] (6.6,1.4)-- (7,1.4);
  		\draw (0.7432211730462593,0.8) node[anchor=north west] {$H$};
  		\draw (6.783095315511885,0.8477079672989145) node[anchor=north west] {G};
  		\draw (9.588852452614713,1.2332318487328873) node[anchor=north west] {$K$};
  		\draw (7.04011123646787,3.4607031636847294) node[anchor=north west] {$J$};
  		\draw (6.8259313023378825,5.152724643311609) node[anchor=north west] {$F$};
  		\draw (3.827412224518069,5.174142636724608) node[anchor=north west] {$I$};
  		\draw (0.8074751532852553,5.0027986894206204) node[anchor=north west] {$E$};
  		\draw [line width=0.8pt] (6.6,4)-- (6.6,3.6);
  		\draw [line width=0.8pt] (6.6,3.6)-- (7,3.6);
  	\end{tikzpicture}
  	
  	\begin{enumerate}
  		\item 
  		\begin{enumerate}
  			\item  Justifie que $FJ=2,1 dam$.
  			\item  Justifie que $(IJ) \parallel (EG)$.
  			\item  Calcule $EG$ et $IJ$.
  		\end{enumerate}
  		\item 
  		\begin{enumerate}
  			\item  Démontre que $(IF) \parallel (GK)$.
  			\item  Calcule $JK$ et $GK$.
  		\end{enumerate}
  	\end{enumerate}
  	\begin{center}
  		\Ovalbox{\textbf{ Exercice 3}}
  	\end{center}
  	On donne les nombres 7; 3 et 12. Construis la quatrième proportionnelle de ces nombres pris dans cet ordre.
  	 
 
  	\sequence{TRIANGLES SEMBLABLES   }
  	\setcounter{q}{0}
  	\act{ Notion de triangles semblables}
  	\hspace*{0.3cm}Dans le jardin public, il y a deux pelouses triangulaires $ABC$ et $IJK$ comme l'indique la figure ci-dessous:
  	
  	\begin{tikzpicture}[scale=1.3]
  		\draw (1.5,1.2) node[anchor=north west] {$A$};
  		\draw (3.9,1.1) node[anchor=north west] {$B$};
  		\draw (6.4,2.6) node[anchor=north west] {$K$};
  		\draw (7.4,1.2) node[anchor=north west] {$J$};
  		\draw (2.6,2.2) node[anchor=north west] {$C$};
  		\draw (4.7,1) node[anchor=north west] {$I$};
  		\draw [line width=0.8pt] (2,1)-- (4,1);
  		\draw [line width=0.8pt] (4,1)-- (2.688138340788525,1.7273369144288742);
  		\draw [line width=0.8pt] (2.688138340788525,1.7273369144288742)-- (2.0008768547026237,1.000926803167337);
  		\draw [line width=0.8pt] (5,1)-- (7.5,1);
  		\draw [line width=0.8pt] (7.5,1)-- (6.6,2.2);
  		\draw [line width=0.8pt] (6.6,2.2)-- (4.999975631141867,1.0000324924981776);
  		\draw (1.9,1.7) node[anchor=north west] {$1$};
  		\draw (3.2,1.8) node[anchor=north west] {$1,5$};
  		\draw (2.8,1) node[anchor=north west] {$2$};
  		\draw (5.1,2) node[anchor=north west] {$3\sqrt{2}$};
  		\draw (7,2) node[anchor=north west] {$2\surd2$};
  		\draw (6,0.9) node[anchor=north west] {$4\sqrt{2}$};
  		\draw (2.4,0.46) node[anchor=north west] {Figue 1};
  		\draw (6,0.46) node[anchor=north west] {Figure 2};
  	\end{tikzpicture}
  	
  	\'A la vue de ces deux pelouses, Dowa a constaté des ressemblances entre elles et décide de comparer certains rapports. 
  	
  	\conAct{Découverte }
  	\begin{enumerate} 
  		\item Calcule les rapports $\dfrac{IK}{CB}$; $\dfrac{JK}{AB}$ et $\dfrac{IJ}{CA}$
  		\item Justifie que les longueurs des côtés des triangles $ABC$ et $IJK$  sont proportionnelles.
  		\item Que peux-tu dire des triangles $ABC$ et $IJK$.
  	\end{enumerate}
  	\dures
  	 
  	\nbVo{Vocabulaire}{ 
  		Les triangles $ABC$ et $IJK$ sont semblables et on a :  $\dfrac{IK}{CB}=\dfrac{JK}{AB}= \dfrac{IJ}{CA}=k=2\sqrt{2}$.
  		
  		$\bullet$ Les côtés $[CB]$ et $[IK]$ sont dits homologues. Il en est de même des 
  		côtés $[AB]$ et $[JK]$,  comme des côtés $[CA]$ et $[IJ]$.
  		
  		$\bullet$ Dans ces conditions les sommets $A$ et $J$  sont dits homologues et les 
   angles $\widehat{BAC}$ et $\widehat{IJK}$ sont aussi dits homologues. 
  		
  		$\bullet$ Le rapport $k$ est appelé rapport de similitude du triangle $ABC$ au triangle $IJK$.
  	}
  	\defs{$T_{1}$ et $T_{2}$ sont deux triangles semblables.
  		
  		On appelle  rapport de similitude de $T_{1}$ à $T_{2}$, le rapport de la longueur 
  		d’un côté quelconque de $T_{2}$ à la longueur de son homologue de $T_{1}$.
  	}
  	\nbVo{N.B}{ Deux triangles superposables sont semblables. Mais la réciproque n'est pas toujours vraie.
  	}
  	
  	\act{Propriétés des triangles semblables} 
  	\conAct{Propriété }
  	
  	$ABC$ est un triangle, $E \in[AB]$ et $F\in[AC]$ tel que $(EF)\parallel(BC)$
  	\begin{tikzpicture}[scale=0.5]
  		\draw [line width=0.8pt] (3,2)-- (9,2);
  		\draw [line width=0.8pt] (3,2)-- (7,5);
  		\draw [line width=0.8pt] (9,2)-- (7,5);
  		\draw [line width=0.8pt] (5.602222222222222,3.951666666666666)-- (6.903333333333332,2);
  		\draw (2,2.3) node[anchor=north west] {$A$};
  		\draw (6.5,1.8) node[anchor=north west] {$E$};
  		\draw (9,2.4) node[anchor=north west] {$B$};
  		\draw (6.5,6) node[anchor=north west] {$C$};
  		\draw (5,5) node[anchor=north west] {$F$};
  	\end{tikzpicture}
  	\begin{enumerate}
  		\item  
  		\begin{enumerate} 
  			\item Justifie que $\dfrac{AE}{AB}=\dfrac{AF}{AC}=\dfrac{EF}{BC}$
  			\item En déduis que les triangles $ABC$ et $AEF$ sont semblables.
  		\end{enumerate}
  		\item Justifie que les triangles $ABC$ et $AEF$ ont leurs angles deux à deux de même mesure.
  		\item Que peux-tu dire des triangles $ABC$ et $AEF$.
  	\end{enumerate}
  	\dures
  	 
  	\propr{Si deux triangles sont semblables, alors tout angle de l’un a même 
  		mesure qu’un angle de l’autre.
  	}
  	 
  	\conAct{Propriété }
  	$ABC$ est un triangle, $E \in(AB)$ et $F\in( AC)$ tel que $(EF)\parallel(BC)$
  	
  	\begin{tikzpicture}[scale=0.6]
  		\draw (4.2,3.2) node[anchor=north west] {$A$};
  		\draw (2.7,3.8) node[anchor=north west] {$E$};
  		\draw (7,1.3) node[anchor=north west] {$B$};
  		\draw (9,4.4) node[anchor=north west] {$C$};
  		\draw (1,1.3) node[anchor=north west] {$F$};
  		\draw [line width=0.8pt] (2,1)-- (9,4);
  		\draw [line width=0.8pt] (7,1)-- (9,4);
  		\draw [line width=0.8pt] (7,1)-- (3.2181717936374437,2.851737557104243);
  		\draw [line width=0.8pt] (2,1)-- (3.2181717936374437,2.851737557104243);
  	\end{tikzpicture}
  	\begin{enumerate}
  		\item 
  		\begin{enumerate} 
  			\item Justifie que $mes\widehat{BAC}=mes\widehat{EAF}$ et que $mes\widehat{AFE}=mes\widehat{ACB}$.
  			\item Déduis-en que $mes\widehat{ABC}=mes\widehat{AEF}$.	
  		\end{enumerate}
  		\item Que peux-tu dire des triangles $ABC$ et $AEF$.
  	\end{enumerate}
  	\dures
  	 
  	\propr{  
  		Si deux triangles sont tels que deux des angles de l’un ont les 
  		mêmes mesures que deux angles de l’autre alors ces triangles sont semblables.
  	}
  	
  	\conAct{Propriété }
  	On considère deux triangles $ABC$ et $EFG$ tels que $\dfrac{AB}{EF}=\dfrac{AC}{EG}=2$ et $mes\widehat{BAC}=mes\widehat{FEG}$
  	\begin{enumerate} 
  		\item Fais une figure.
  		\item Marque les points $I$ et $J$ avec $I\in[EF)$; $J\in[EG)$ tels que $EI=AB$ et $EJ=AC$.
  		\item Justifie que les triangles $ABC$ et $EIJ$ sont  semblables.
  		\item 
  		\begin{enumerate} 	
  			\item  Justifie que $(IJ)\parallel(FG)$.
  			\item Déduis-en que les triangles $ABC$ et $EFG$ sont  semblables.
  		\end{enumerate}
  	\end{enumerate}
  	\dures
  	 
  	\propr{Si deux triangles sont tels que deux côtés de l’un ont des longueurs proportionnelles à celles de deux côtés de l’autre et que les angles formés par ces côtés ont la même mesure, alors ces triangles sont semblables.
  	
  		Si trois triangles $T_{1}$, $T_{2}$ et $T_{3}$ sont tels que $T_{1}$ est semblable à $T_{2}$ et $T_{2}$ semblable $T_{3}$ alors $T_{1}$ et $T_{3}$ sont semblables.
  	}
 
  	
  	\conAct{Approfondissement }
  	\hspace{0.4cm}Au cours de la répartition des domaines, l'ingénieur en charge des travaux a réalisé le morcellement d'un domaine de la façon suivante:
  	
  	Il trace un triangles $ABC$ isocèle rectangle en $A$ tel que $AB=AC=3 cm$ puis il place le point $O$ milieu du segment $[BC]$ et le point $H$, symétrique du point $O$ par rapport à $C$. Ensuite, il trace la perpendiculaire à la droite $(BC)$ passant par $H$  qui coupe les droites $(AB)$ et  $(AC)$ en $E$ et en $F$.
  	\begin{enumerate} 	
  		\item 
  		\begin{enumerate}
  			\item  Fais une figure du dit domaine.
  			\item Calcule $OC$.
  		\end{enumerate}
  		\item 
  		\begin{enumerate}
  			\item  Démontre que les triangles $ABC$ et $FCH$ sont semblables.
  			\item  Calcule le rapport de similitude du triangle $FCH$ au triangle $ABC$
  			\item  Déduis-en $FC$. 
  		\end{enumerate}
  	\end{enumerate}
  	\dures
  	 
  	\begin{center}
  		\Ovalbox{\textbf{ Exercice de maison}}
  	\end{center}
  	\begin{center}
  		\Ovalbox{\textbf{ Exercice 1}}
  	\end{center}
  	En considérant la figure de la consigne 4: 
  	\begin{enumerate}
  		\item  Démontre que les triangles $BEH$ et $ABC$ sont semblables.
  		\item En déduis que  les triangles $BEH$ et $FCH$ sont semblables.
  		\item  Calcule le rapport de similitude du triangle $BEH$ au triangle $ABC$
  		\item Justifie que $BE=9 cm$ et calcule $EH$.
  		\item Donne la nature exacte du triangle $BEH$
  	\end{enumerate}
  	
  	\begin{center}
  		\Ovalbox{\textbf{ Exercice 2}}
  	\end{center}
  	$\mathcal{(C)}$ est cercle de centre $I$ et de rayon $5cm$.
  	
  	\begin{tikzpicture}[scale=0.6]
  		\draw [shift={(2,1)},line width=0.8pt] (0,0) -- (0:0.9012269389796509) arc (0:30:0.9012269389796509) -- cycle;
  		\draw [line width=0.8pt] (2,1)-- (8,1);
  		\draw [line width=0.8pt] (5,1) circle (3cm);
  		\draw [line width=0.8pt] (2,1)-- (6.400923401581161,3.652812398738793);
  		\draw [line width=0.8pt] (8,1)-- (6.400923401581161,3.652812398738793);
  		\draw [line width=0.8pt] (6.400923401581161,3.652812398738793)-- (6.399623121927212,1);
  		\draw [line width=0.8pt] (6.8,1)-- (6.8,1.4);
  		\draw [line width=0.8pt] (6.8,1.4)-- (6.399819182522024,1.400000088627936);
  		\draw (8,1.6) node[anchor=north west] {$A$};
  		\draw (6.2,4.7) node[anchor=north west] {$B$};
  		\draw (1,1.4) node[anchor=north west] {$C$};
  		\draw (6,1) node[anchor=north west] {$H$};
  		\draw (4.4,1) node[anchor=north west] {$I$};
  		\draw (3.1,1.9) node[anchor=north west] {$30\textrm{\degre}$};
  		\begin{scriptsize}
  			\draw [fill=black] (5,1) circle (3pt);
  		\end{scriptsize}
  	\end{tikzpicture}
  	\begin{enumerate}
  		\item  Justifie que le triangle $ABC$ est rectangle puis calcule $AB$ et $BC$.
  		\item Démontre que  les triangles $ABC$ et $BAH$ sont semblables.
  		\item  Calcule le rapport de similitude du triangle $ABC$ au triangle $BAH$
  		\item Calcule $CH$; $BH$ et $IH$.	
  	\end{enumerate}
   \sequence{TRIANGLES RECTANGLES}
   \setcounter{q}{0}
  	\act{ Les relations métriques}
  	
  	\hspace{0.4cm}\`A l’entrée du jardin public  une grande enseigne lumineuse comportant des indications est matérialisé par la figure ci-dessous :
  	
  	\begin{tikzpicture}[scale=0.6]
  		\draw [line width=0.8pt] (3,1)-- (11,1);
  		\draw [line width=0.8pt] (11,1)-- (6.11,4.898173178458289);
  		\draw [line width=0.8pt] (6.11,4.898173178458289)-- (3,1);
  		\draw [line width=0.4pt] (6.11,4.898173178458289)-- (6.11,1);
  		\draw [line width=0.8pt] (6.11,1.3027787421418804)-- (6.500582150793547,1.3027787421418808);
  		\draw [line width=0.8pt] (6.500582150793547,1.3027787421418808)-- (6.5,1);
  		\draw [line width=0.4pt] (5.887288688440678,4.619019718246128)-- (6.1625619927316935,4.373605940805601);
  		\draw [line width=0.8pt] (6.1625619927316935,4.373605940805601)-- (6.391313809399421,4.6739175656708145);
  		\draw (5.7,5.6) node[anchor=north west] {$B$};
  		\draw (2.2,1.3) node[anchor=north west] {$A$};
  		\draw (5.7,0.9) node[anchor=north west] {$H$};
  		\draw (11,1.5) node[anchor=north west] {$C$};
  		
  	\end{tikzpicture}
  	
  	On observant cette enseigne, Basilia se demande s'il n'aurait pas une relation entre certains triangles qui se trouvent dans cette enseigne.
  	
  	\conAct{Relations métriques}
  	\begin{enumerate}
  		\item 
  		\begin{enumerate}
  			\item Démontre que les triangles $ABC$ et $HBC$ sont semblables.
  			\item Déduis-en que $BC^2= CH \times CA$ et $AC\times HB= AB\times BC$
  		\end{enumerate}
  		\item
  		\begin{enumerate} 
  			\item  Démontre que les triangles $ABC$ et $ABH$ sont semblables
  			\item Déduis-en que $AB^2= AH \times AC$.
  		\end{enumerate}
  		\item 
  		\begin{enumerate} 
  			\item Démontre que les triangles $ABH$ et $BHC$ sont semblables.
  			\item  Déduis-en que $BH^2= AH \times HC$.
  		\end{enumerate}
  	\end{enumerate}
  	\dures
  	 
  	\propr{ \ding{118} Dans un triangle rectangle, le produit des longueurs des côtés de l’angle droit est égal au produit des longueurs de l’hypoténuse et de la hauteur relative à l’hypoténuse 
  		
  		\ding{118} Dans un triangle rectangle, la longueur de la hauteur relative à 
  		l’hypoténuse est moyenne proportionnelle entre les longueurs des 
  		segments qu’elle détermine sur l’hypoténuse.
  		
  		\ding{118} Dans un triangle rectangle, le carré de la longueur d’un côté de 
  		l’angle droit est égal au produit des longueurs de l’hypoténuse et du 
  		projeté orthogonal dudit côté sur l’hypoténuse.
  		
  		\textbf{Illustration}
  		
  		\begin{tikzpicture}[scale=0.5]
  			\draw [line width=0.8pt] (1,6)-- (1,1);
  			\draw [line width=0.8pt] (1,1)-- (9,1);
  			\draw [line width=0.8pt] (9,1)-- (1,6);
  			\draw [line width=0.4pt] (1,1)-- (3.2480898876404494,4.594943820224719);
  			\draw [line width=0.8pt] (3.1146119493728994,4.337512302241196)-- (3.372486789431546,4.179168102205185);
  			\draw [line width=0.8pt] (3.372486789431546,4.179168102205185)-- (3.516276081294909,4.427327449190682);
  			\draw (0,1.7) node[anchor=north west] {$A$};
  			\draw (9.030636089682329,1.8218544181178824) node[anchor=north west] {$B$};
  			\draw (3.254061498837947,6.247183490636105) node[anchor=north west] {$H$};
  			\draw (0.07863453077144521,6.9228062497991925) node[anchor=north west] {$C$};
  			\draw [line width=0.8pt] (1,1.4)-- (1.4,1.4);
  			\draw [line width=0.8pt] (1.4,1.4)-- (1.4,1);
  		\end{tikzpicture}
  		
  		On a
  		
  		$\bullet$ $AB\times AC= AH \times BC$
  		
  		$\bullet$ $AH^2=CH \times BH$
  		
  		$\bullet$ $AB^2=BH \times BC$
  		
  		$\bullet$ $AC^2= CH \times BC$
  		
  		Les relations ci-dessus sont appelées des relations métriques 
  	}
  	\conAct{Approfondissement}
  	
  	On considère la figure suivante où $O$ est le centre du cercle et la droite $(DC)$ est tangente au cercle au point $C$ telle que $AB=5cm$ et $AC=10cm$.
  	
  	\hspace{0.7cm} {\begin{tikzpicture}[scale=0.5]
  			\draw [line width=0.8pt] (11,5) circle (2cm);
  			\draw [line width=0.8pt] (4,7)-- (11,7);
  			\draw [line width=0.8pt] (4,7)-- (11,3);
  			\draw [line width=0.8pt] (11,7)-- (9.276923076923087,3.9846153846153793);
  			\draw [line width=0.8pt] (11,7)-- (11,3);
  			\draw (10.3,3) node[anchor=north west] {$A$};
  			\draw (8.4,4) node[anchor=north west] {$B$};
  			\draw (10.6,8.1) node[anchor=north west] {$C$};
  			\draw (11.1,5.6) node[anchor=north west] {$O$};
  			\draw (3.2,6.8) node[anchor=north west] {$D$};
  			\begin{scriptsize}
  				\draw [fill=black] (11,5) circle (2pt);
  			\end{scriptsize}
  	\end{tikzpicture}}
  	
  	\begin{enumerate}
  		\item Précise la nature des triangles $ABC$ et $ADC$.
  		\item Calcule $BC$; $AD$; $BD$ et $DC$.
  	\end{enumerate}
  	\dures
  	 
  	\conAct{ Problème de construction géométrique} 
  	On donne $ AB=3cm $ et $ CB=4cm $ et $ B \in [AC] $. On cherche à construire le segment $[DB]$ de longueur $\sqrt{12}cm$. On remarquera que $ DB=\sqrt{AB\times CB} $. Pour cela:
  	\begin{itemize}
  		\item Construis le segment [AC] puis marque le point B.
  		\item Construis le demi-cercle $ \mathcal{(C)} $ de diamètre $ [AC] $
  		\item La perpendiculaire à $ (AC) $ passant par $ B $ coupe $ \mathcal{(C)} $ en $ E $.
  		\item Justifie que $ EB=\sqrt{12} cm $.
  		\item Déduis-en $ DB $
  	\end{itemize}
  	\dures
  	 
  	\nbVo{Remarques}{$\bullet$ Les propriétés métriques servent à calculer des longueurs de segments définis dans un triangle rectangle muni d'une hauteur relative à son hypoténuse
  		
  		$\bullet$ Les propriétés métriques servent  aussi à résoudre des problèmes de construction géométrique:
  		
  		- segment $[MN]$ tel que $MN=\sqrt{AB\times CD}$ où  $[AB]$ et $[CD]$ sont donnés
  		
  		- segments $[EF]$ et $[GH]$ tels que $AB\times CD= EF\times GH$ et $AB\times EF= CD\times GH$. 
  	}
  	\begin{center}
  		\Ovalbox{\textbf{ Exercice de maison}}
  	\end{center}
  	\begin{center}
  		\Ovalbox{\textbf{ Exercice 1}}
  	\end{center}
  	
  	$ABC$ est un triangle tels que $AB=3cm$; $AC=4cm$ et $BC=5cm$. $H$ est le projeté orthogonal de $A$ sur la droite $(BC)$ 
  	\begin{enumerate}
  		\item Fais une figure.
  		\item  Justifie que le triangle $ABC$ est rectangle en $A$.
  		\item  Calcule les longueurs $AH$; $BH$ et $CH$.
  		\item  La perpendiculaire en $C$ à la droite $(BC)$ coupe la droite $(AB)$ en $D$. Calcule $BD$ et $CD$.
  	\end{enumerate}
  	\begin{center}
  		\Ovalbox{\textbf{ Exercice 2}}
  	\end{center}
  	
  	L'aire de la surface d'un triangle $ABC$ rectangle en $A$ tel que $AB=5cm$ est égale à $30 cm^2$. La hauteur issue de $A$ coupe $(BC)$ en $H$. 
  	
  	\begin{enumerate}
  		\item  Calcule $BC$; $AH$ et $BH$
  		\item  
  		\begin{enumerate}
  			\item  Démontre que les triangles $ABH$ et $ACH$ sont semblables.
  			\item  Calcule le rapport de similitude du triangle $ACH$  au triangle $ABH$
  		\end{enumerate}
  		\item  La bissectrice de l'angle $\widehat{ABC}$ coupe $(AC)$ en $M$ et la parallèle à $(BC)$ passant par $M$ coupe $(AB)$ en $N$. Justifie que $MN=NB$.
  	\end{enumerate}
  	
  	\begin{center}
  		\Ovalbox{\textbf{ Exercice 3}}
  	\end{center}
  	On donne $AB=4cm$ et $CD=3cm$. Construis un segment $MN=\sqrt{AB\times CD}$.
  	 \sequence{ANGLES ET CERCLE}
  	\setcounter{q}{0}
  	\act{ Notion d’angle inscrit dans un cercle}
  	
  	\hspace{0.4cm} Le modèle ci-dessous représente le jardin public circulaire de centre $O$, subdivisé en de petites portions de terre pour l'implantation des différentes variétés de fleures.
  	
  	\definecolor{sqsqsq}{rgb}{0.12549019607843137,0.12549019607843137,0.12549019607843137}
  	\begin{tikzpicture}[scale=0.7]
  		\draw [shift={(4.899924948982784,4.514946688603571)},line width=1.2pt] (0,0) -- (-58.35579204552263:0.4885798197886218) arc (-58.35579204552263:8.892685005662901:0.4885798197886218) -- cycle;
  		\draw [line width=0.8pt] (3,1)-- (6.998496631376921,1.1096571423456547);
  		\draw [line width=0.8pt] (6.998496631376921,1.1096571423456547)-- (4.899924948982784,4.514946688603571);
  		\draw [line width=0.8pt] (4.899924948982784,4.514946688603571)-- (2.9999029071038774,0.9950154051302258);
  		\draw [line width=0.8pt] (4.9676397630584175,2.2073908987185793) circle (2.3085491156735314cm);
  		\draw [line width=0.8pt] (4.899924948982784,4.514946688603571)-- (8,5);
  		\draw [line width=0.8pt] (7.116288742546953,3.0516104095212624)-- (9,4);
  		\draw [line width=0.8pt] (7.116288742546953,3.0516104095212624)-- (9,2);
  		\draw [line width=0.8pt] (4.899924948982784,4.514946688603571)-- (4.9676397630584175,2.2073908987185793);
  		\draw [line width=0.8pt] (4.9676397630584175,2.2073908987185793)-- (6.998496631376921,1.1096571423456547);
  		\draw [line width=0.8pt] (2.9999029071038774,0.9950154051302258)-- (4.9676397630584175,2.2073908987185793);
  		\draw [line width=0.8pt] (2.9999029071038774,0.9950154051302258)-- (5.030926987805755,-0.10029056687861655);
  		\draw [line width=0.8pt] (5.030926987805755,-0.10029056687861655)-- (6.998496631376921,1.1096571423456547);
  		\draw (1.9,1.3) node[anchor=north west] {$A$};
  		\draw (4.6,-0.1) node[anchor=north west] {$B$};
  		\draw (7.1,1.5) node[anchor=north west] {$C$};
  		\draw (4.5,5.2) node[anchor=north west] {$D$};
  		\draw (8.5,2.8) node[anchor=north west] {$E$};
  		\draw (7.9,4.4) node[anchor=north west] {$F$};
  		\draw (6.5,5.6) node[anchor=north west] {$G$};
  		\draw (4.5973732302045605,2.1) node[anchor=north west] {$O$};
  		\draw (1.6,3.4) node[anchor=north west] {$\mathcal{(C)}$};
  		\begin{scriptsize}
  			\draw [fill=black] (4.9676397630584175,2.2073908987185793) circle (1pt);
  			\draw [fill=sqsqsq] (8.40663827114248,3.7012609630003417) circle (2pt);
  			\draw [fill=black] (8.483341453723407,2.2884324777923832) circle (2pt);
  			\draw [fill=black] (7.01949385956244,4.846585085058596) circle (2pt);
  		\end{scriptsize}
  	\end{tikzpicture}
  	
  	\conAct{Découverte}
  	\begin{enumerate}
  		\item 	\begin{enumerate}
  			\item Cite deux angles dont les sommets sont situés sur le cercle $\mathcal{(C)}$ et les côtés sont sécants au cercle.
  			\item Précise pour chacun d'eux l'arc intercepté.
  			\item  Pour chacun des angles cités donne l'angle au centre qui intercepte le même arc.
  		\end{enumerate}
  		\item  Indique sur la figure un quadrilatère dont les quatre sommets appartiennent au cercle  $\mathcal{(C)}$.
  		\item 	\begin{enumerate}
  			\item  Construis le cercle $\mathcal{(C)}$ de centre $O$ et de rayon $3cm$ puis construis l'angle $\widehat{COD}$ tel que $mes\widehat{COD}=70\textmd{\degre}$.
  			\item  Trace en bleu le petit arc $\overset{\frown}{CD}$ puis colorie avec la même couleur l'angle au centre qui l'intercepte.
  			\item Trace en rouge le grand arc  $\overset{\smile}{CD}$
  		\end{enumerate}
  		\bcinfo{\textbf{Information}}: \textit{\textbf{ Dans le cercle $\mathcal{(C)}$ l'angle au centre qui intercepte le petit l'arc $\overset{\frown}{CD}$ est appelé angle saillant et est noté  $\widehat{COD}$ et celui qui intercepte le grand arc  $\overset{\smile}{CD}$ est appelé angle rentrant et est noté  $\check{COD}$.}}
  	\end{enumerate}
  	\dures
  	
  	\ret{ $A$, $B$ et $O$ sont trois points non alignés du plan. $\textbf{[OA)}$ et $\textbf{[OB)}$ deux demi-droites non opposées d'origine commune $O$. 
  		
  		\ding{118}L'angle défini dans la partie intérieure aux deux demi-droites se note $\widehat{AOB}$  et se lit 	\og angle saillant $AOB$\fg{}.
  		
  		\ding{118} L'angle défini dans la partie extérieure aux deux demi-droites se note $\check{AOB}$  et se lit 	\og angle rentrant $AOB$\fg{}.\\
  		\begin{tikzpicture}[scale=0.7]
  			%\clip(-15.2,-3.56) rectangle (4.04,9.16);
  			\draw [shift={(-10,6)},line width=1pt,fill=black,fill opacity=0.36] (0,0) -- (-18.43494882292201:0.6) arc (-18.43494882292201:52.87393813172565:0.6) -- cycle;
  			\draw [shift={(-10,6)},line width=1pt] (0,0) -- (52.87393813172565:0.6) arc (52.87393813172565:341.565051177078:0.6) -- cycle;
  			\draw [line width=1pt] (-10,6)-- (-8.38,8.14);
  			\draw [line width=1pt] (-10,6)-- (-7,5);
  			\draw [color=black](-15.4,6.5) node[anchor=north west] {Angle  rentrant};
  			\draw (-9.4,6.5) node[anchor=north west] {angle saillant};
  			\draw (-9.24,8.82) node[anchor=north west] {$A$};
  			\draw (-10.5,6) node[anchor=north west] {$O$};
  			\draw (-8.2,5) node[anchor=north west] {$B$};
  			\begin{scriptsize}
  				\draw [color=black] (-8.549106052193228,7.91661299278179)-- ++(-1.5pt,-1.5pt) -- ++(3pt,3pt) ++(-3pt,0) -- ++(3pt,-3pt);
  				\draw [color=black] (-7.774,5.258)-- ++(-1.5pt,-1.5pt) -- ++(3pt,3pt) ++(-3pt,0) -- ++(3pt,-3pt);
  			\end{scriptsize}
  	\end{tikzpicture}}
  	
  	\conAct{Approfondissement}
  	
  	En utilisant la figure de l'activité 1, complète le tableau suivant:
  	
  	\begin{tabular}{|p{1.5cm}| p{3cm}| p{2cm}|}
  		\hline
  		Angle inscrit	& Angle au centre associé &Arc intercepté  \\
  		\hline
  		&  &  \\
  		\hline
  		&  &  \\
  		\hline
  		&  &  \\
  		\hline
  		&  &  \\
  		\hline
  	\end{tabular}
  	\dures
  
   
  	\act{Propriétés}
  	
  	\conAct{ La mesure d’un angle 
  		inscrit et la mesure de son angle au centre 
  		associé.}
  	
  	On considère le cercle $\mathcal{(C)}$ de centre $O$. \\
  	\hspace*{-1cm}  \begin{minipage}{4cm}
  		\begin{tikzpicture}[line cap=round,line join=round,>=triangle 45,x=1cm,y=1cm]
  			\clip(1,1) rectangle (6,5);
  			\draw [line width=1.2pt] (4,3) circle (1.5839823231336896cm);
  			\draw (1.5,3.79) node[anchor=north west] {$ \mathcal{(C)} $};
  			\draw [line width=1.2pt] (2.879955358032547,4.1200446419674535)-- (4,3);
  			\draw [line width=1.2pt] (4,3)-- (5.11,4.13);
  			\draw [line width=1.2pt] (2.879955358032547,4.1200446419674535)-- (4.510939250093889,1.500686462839245);
  			\draw [line width=1.2pt] (4.510939250093889,1.500686462839245)-- (5.11,4.13);
  			\draw [line width=1.2pt] (4,3)-- (4.510939250093889,1.500686462839245);
  			\begin{scriptsize}
  				\draw [fill=black] (4,3) circle (0.5pt);
  				\draw[color=black] (3.93,3.48) node {$O$};
  				\draw [fill=black] (5.11,4.13) circle (0.5pt);
  				\draw[color=black] (5.15,4.46) node {$A$};
  				\draw [fill=black] (2.879955358032547,4.1200446419674535) circle (0.5pt);
  				\draw[color=black] (2.79,4.52) node {$B$};
  				\draw [fill=black] (4.510939250093889,1.500686462839245) circle (1.5pt);
  				\draw[color=black] (4.55,1.24) node {$M$};
  			\end{scriptsize}
  		\end{tikzpicture}
  	\end{minipage} \hspace*{0.7cm}
  	\begin{minipage}{5.5cm}
  		1. Justifie la nature des triangles $ AOB $, $ AOM $ et $ BOM $  \\
  		2. Déduis-en que $ mes\widehat{OMB}=mes\widehat{OBM} $, $ mes\widehat{OMA}=mes\widehat{OAM} $ et $ mes\widehat{OAB}=mes\widehat{OBA} $\\
  		3. Justifie que $mes\widehat{AMB}=mes\widehat{OBM}+mes\widehat{OBM} $.\\
  		4. Justifie que $ 2 \times mes\widehat{AMB}=mes \widehat{AOB} $ puis déduis-en que $ mes\widehat{AMB}=\dfrac{1}{2}mes\widehat{AOB}  $.
  	\end{minipage}\\
  	\dures
  	 
  	\propr{Un angle inscrit dans un cercle a pour mesure la moitié de la mesure de l'angle au cercle qui lui est associé.}
  	\conAct{Égalité des mesures de deux angles 
  		inscrits interceptant le même arc}
  	
  	On considère le cercle $\mathcal{(C)}$ de centre $O$.\\
  	\begin{minipage}{3cm}
  		
  		\begin{tikzpicture}[scale=0.3]
  			\draw (9.6,-1.5) node[anchor=north west] {$A$};
  			\draw (2.6243335751979844,-0.2) node[anchor=north west] {$B$};
  			\draw (5.658197550450718,-5.8) node[anchor=north west] {$O$};
  			\draw (-0.9,-5.1) node[anchor=north west] {$\mathcal{(C)}$};
  			\draw [line width=0.8pt] (6,-6) circle (4.650571779707891cm);
  			\draw [line width=0.8pt] (3.920201073266196,-1.8404021465323916)-- (2.33891564305259,-8.867800413126826);
  			\draw [line width=0.8pt] (2.33891564305259,-8.867800413126826)-- (9.667744145468305,-3.1407219859554774);
  			\draw [line width=0.8pt] (6,-6)-- (3.920201073266196,-1.8404021465323916);
  			\draw (0.5720138272328998,-8.570859300494472) node[anchor=north west] {$M$};
  			\draw [line width=0.8pt] (3.920201073266196,-1.8404021465323916)-- (5.897892469920344,-10.649450712774193);
  			\draw [line width=0.8pt] (5.897892469920344,-10.649450712774193)-- (9.667744145468301,-3.140721985955479);
  			\draw (5.3458880235864665,-10.6) node[anchor=north west] {$C$};
  		\end{tikzpicture}
  	\end{minipage} \hspace*{0.5cm}
  	\begin{minipage}{6cm}
  		\begin{enumerate}
  			\item Quel est l'arc intercepté par les   angles $ \widehat{AMB} $ et $ \widehat{ACB} $ incrits au cercle $ \mathcal{(C)} $.
  			\item Justifie que $mes\widehat{AMB}=\dfrac{1}{2} mes\widehat{AOB}$
  			\item Justifie que $mes\widehat{ACB}=\dfrac{1}{2} mes\widehat{AOB}$
  			\item Déduis-en que $mes\widehat{AMB}=mes\widehat{ACB}$
  		\end{enumerate}
  	\end{minipage} \\
  	\dures
  	
  	   	\propr{Dans un cercle, deux angles inscrits qui interceptent le même arc ont la même mesure.}
  	\conAct{Egalité des mesures de deux angles inscrits interceptant deux arcs de mêmes longueurs}
  	
  	On considère le cercle $\mathcal{(C)}$ de centre $O$.\\
  	\begin{minipage}{4cm}
  		\begin{tikzpicture}[scale=0.3]
  			\draw (9.74080018137954,-1.2998083277878025) node[anchor=north west] {$A$};
  			\draw (2.6287886938759626,0.057088600749059665) node[anchor=north west] {$B$};
  			\draw (4,-5) node[anchor=north west] {$O$};
  			\draw (-1.4,-5) node[anchor=north west] {$\mathcal{(C)}$};
  			\draw [line width=0.8pt] (6,-6) circle (4.650571779707891cm);
  			\draw [line width=0.8pt] (3.920201073266196,-1.8404021465323916)-- (2.33891564305259,-8.867800413126826);
  			\draw [line width=0.8pt] (2.33891564305259,-8.867800413126826)-- (9.667744145468305,-3.140721985955477);
  			\draw (0.5700485264407162,-8.50539891381114) node[anchor=north west] {$M$};
  			\draw (7.82242866172397,-9.628348096048544) node[anchor=north west] {$C$};
  			\draw [line width=0.8pt] (2.338915643052589,-8.867800413126826)-- (10.650540118826823,-5.982839539702602);
  			\draw [line width=0.8pt] (10.650540118826823,-5.982839539702602)-- (8.079798926733805,-10.15959785346761);
  			\draw [line width=0.8pt] (6.845898078868418,-1.9634842931014866)-- (7.240582250848796,-0.9373054459525104);
  			\draw [line width=0.8pt] (4.342803706760804,-10.943580322313569)-- (5.033407322372828,-9.907674898895541);
  			\draw (11.050907560656515,-4.528287226720337) node[anchor=north west] {$N$};
  			\draw [line width=0.8pt] (3.920201073266196,-1.8404021465323916)-- (8.079798926733805,-10.15959785346761);
  		\end{tikzpicture}
  	\end{minipage}
  	\begin{minipage}{5cm}
  		\begin{enumerate}
  			\item  Démontre que $mes\widehat{AOB}$ = $ mes\widehat{COM}$.
  			\item Démontre que $mes\widehat{AMB} $ = $\dfrac{1}{2} mes\widehat{AOB}$.
  			\item Démontre que $mes\widehat{CNM}$ = $\dfrac{1}{2} mes\widehat{COM}$.
  			\item Déduis-en que $mes\widehat{AMB}= mes\widehat{CNM}$.
  		\end{enumerate}
  	\end{minipage}\\ \\
  	\dures
  	 
  	\propr{	Dans un cercle, deux angles inscrits qui interceptent deux arcs de même longueur ont la même mesure.}
  	\conAct{Approfondissement }
  	
  	$\mathcal{(C)}$ est un cercle de centre $O$.
  	
  	\begin{minipage}{4cm}
  		\begin{tikzpicture}[scale=0.56]
  			\draw [shift={(8,1.9740786527075755)},line width=0.8pt] (0,0) -- (72.55768066203358:0.9288980121593072) arc (72.55768066203358:111.04233894167875:0.9288980121593072) -- cycle;
  			\draw [line width=0.8pt] (8,5) circle (3.0259213472924245cm);
  			\draw [line width=0.8pt] (5.971942258368722,7.24570296268439)-- (8,1.9740786527075755);
  			\draw [line width=0.4pt] (8,1.9740786527075755)-- (9.730602988097422,7.4821791429283016);
  			\draw [line width=0.8pt] (9.730602988097422,7.4821791429283016)-- (5.971942258368722,7.24570296268439);
  			\draw [line width=0.8pt] (5.971942258368722,7.24570296268439)-- (9,10);
  			\draw [line width=0.4pt] (5.971942258368722,7.24570296268439)-- (5.072171663870785,4.235787180061188);
  			\draw [line width=0.8pt] (5.072171663870785,4.235787180061188)-- (9.730602988097422,7.4821791429283016);
  			\draw [line width=0.8pt] (5.971942258368722,7.24570296268439)-- (8,5);
  			\draw [line width=0.8pt] (8,5)-- (9.730602988097422,7.4821791429283016);
  			\draw (7.457703571799894,2.1) node[anchor=north west] {$A$};
  			\draw (9.740138687391335,8.566226105372072) node[anchor=north west] {$B$};
  			\draw (5.0691086833902474,8.274286730122006) node[anchor=north west] {$C$};
  			\draw (4.14021067123094,4.558694681484782) node[anchor=north west] {$D$};
  			\draw (7.484243515004446,10.264782470463373) node[anchor=north west] {$E$};
  			\draw (7.2,5.2) node[anchor=north west] {$O$};
  			\draw (7.3,4.2) node[anchor=north west] {$38\textmd{\degre}$};
  			\begin{scriptsize}
  				\draw [fill=black] (8.194027724130988,9.26689408162057) circle (3pt);
  			\end{scriptsize}
  		\end{tikzpicture}
  	\end{minipage}
  	\begin{minipage}{3.5cm}
  		Calcule $mes\widehat{ECB}$; $mes\widehat{BDC}$ et $mes\widehat{BOC}$.
  	\end{minipage}
  	\dures
  	 
  	\act{Propriétés}
  	
  	\conAct{Quadrilatère inscriptible dans un cercle}
  	
  	On considère un quadrilatère $ABCD$ inscrit dans un cercle $\mathcal{(C)}$ de centre $O$.
  	
  	\begin{tikzpicture}[scale=0.6]
  		\draw (4.898229587610756,5.738942516029775) node[anchor=north west] {$A$};
  		\draw (6.461955546797498,1.7052403258548925) node[anchor=north west] {$C$};
  		\draw (6.65742129169584,4.903770696918676) node[anchor=north west] {$D$};
  		\draw (2.3038660644145716,3.4644320299399736) node[anchor=north west] {$B$};
  		\draw (5.289161077407441,3.4644320299399736) node[anchor=north west] {$O$};
  		\draw [line width=0.8pt] (5,3) circle (2cm);
  		\draw [line width=0.8pt] (3,3)-- (6.349257084775684,1.5236852235439005);
  		\draw [line width=0.8pt] (5,3)-- (6.349257084775684,1.5236852235439005);
  		\draw [line width=0.8pt] (3,3)-- (5,3);
  		\draw [line width=0.8pt] (5.030485972105728,4.999767637878153)-- (3,3);
  		\draw [line width=0.8pt] (5.030485972105728,4.999767637878153)-- (6.575574128047232,4.231895355551032);
  		\draw [line width=0.8pt] (6.575574128047232,4.231895355551032)-- (6.349257084775685,1.5236852235439011);
  		\draw [line width=0.8pt] (5.030485972105728,4.999767637878153)-- (5,3);
  	\end{tikzpicture}
  	
  	\begin{enumerate}
  		\item Calcule $mes\widehat{AOC}+mes\check{AOC}$.
  		\item  Exprime la 	mesure de l’angle $mes\widehat{ABC}$ en fonction celle de    $mes\widehat{AOC}$.
  		\item Exprime la mesure de l’angle $\widehat{ADC}$ en fonction de celle de $\check{AOC}$.
  		\item  Déduis-en que $mes\widehat{ABC}+mes\widehat{ADC}=180^\circ$.
  		\item Calcule    $mes\widehat{BAD}+mes\widehat{BCD}$ puis conclure.
  	\end{enumerate}
  	\dures
  	  
  	\propr {$P_{1}$ Si un quadrilatère est inscrit dans un cercle, alors ses angles opposés sont supplémentaires.
  		
  		$P_{2}$ Si un quadrilatère est tel que deux de ses angles opposés sont supplémentaires, alors il est inscriptible dans un cercle.
  	}
  	
  	\conAct{Mesure des angles d’un polygone 
  		régulier inscrit dans un cercle}
  	
  	On considère un polygone régulier à n côté ($n\geqslant 3$). $A$, $B$ et $C $ désignent trois sommets 
  	consécutifs de ce polygone. (Voir figure)
  	
  	\begin{minipage}{5cm}
  		\begin{tikzpicture}[scale=0.8]
  			\draw (4.815799203465134,5.8047243872488234) node[anchor=north west] {$A$};
  			\draw (7.1300380904521585,3.5081514459639997) node[anchor=north west] {$C$};
  			\draw (6.511729990875472,5.045088722054612) node[anchor=north west] {$B$};
  			\draw (4.4,3.2) node[anchor=north west] {$O$};
  			\draw [line width=0.8pt] (5,3) circle (2cm);
  			\draw [line width=0.8pt] (5.030485972105728,4.999767637878153)-- (6.414250316875022,4.420806293620001);
  			\draw [line width=0.8pt] (6.414250316875022,4.420806293620001)-- (6.998828347904919,3.039405364983709);
  			\draw [line width=0.8pt] (5.030485972105728,4.999767637878153)-- (5,3);
  			\draw [line width=0.8pt] (5,3)-- (6.998828347904919,3.039405364983709);
  			\draw [line width=0.8pt] (5,3)-- (6.414250316875022,4.420806293620001);
  			\draw (2,3.2) node[anchor=north west] {$\mathcal{(C)}$};
  		\end{tikzpicture}
  	\end{minipage}
  	\begin{minipage}{3cm}
  		$\mathcal{(C)}$ est un cercle de centre $O$ et $mes\widehat{AOB}=mes\widehat{BOC}$.
  	\end{minipage}
  	\begin{enumerate}
  		\item Exprime $mes\widehat{AOB}$ en fonction de $n$.
  		\item 
  		\begin{enumerate}
  			\item Justifie que les triangles $AOB$ et $BOC$ sont superposables et isocèles.
  			\item Déduis-en que  $mes\widehat{ABO}=mes\widehat{OBC}$ et $mes\widehat{ABC}=2\times mes\widehat{ABO}$
  		\end{enumerate}
  		\item Démontre que $mes\widehat{ABC}=108^\circ-\dfrac{360^\circ}{n}$.
  	\end{enumerate}
  	\dures
  	 
  	\propr{La mesure d'un angle au sommet d'un polygone régulier de $n$ $(n\geq3)$ côtés est $180^\circ-\dfrac{360^\circ}{n}$}
  	\begin{center}
  		\Ovalbox{\textbf{ Exercice de maison}}
  	\end{center}
  	\begin{center}
  		\Ovalbox{\textbf{ Exercice 1 }}
  	\end{center}
  	On considère la figure suivante :\\
  	$I$ est le centre du cercle $\mathcal{(C)}$ 
  	\begin{center}
  		\begin{tikzpicture}[scale=0.6]
  			%\clip(-15.526527067221897,-6.173133848899474) rectangle (9.253351576442597,10.346785246876868);
  			\draw [shift={(-10.009640368544177,6.003452756109792)},line width=0.8pt] (0,0) -- (-19.705291171552158:0.7601189767995243) arc (-19.705291171552158:43.114055398612116:0.7601189767995243) -- cycle;
  			\draw [shift={(-10,3.0033185020760045)},line width=0.8pt] (0,0) -- (35.14735441422391:0.7601189767995243) arc (35.14735441422391:99.67702680331813:0.7601189767995243) -- cycle;
  			\draw [line width=0.4pt] (-10,6) circle (2.9966814979239955cm);
  			\draw [line width=0.4pt] (-10.993114152924402,8.827335190468446)-- (-12.193777050744163,3.9585685777797375);
  			\draw [line width=0.4pt] (-12.193777050744163,3.9585685777797375)-- (-7.178805936833635,4.989572339078714);
  			\draw [line width=0.4pt] (-10.993114152924402,8.827335190468446)-- (-7.178805936833635,4.989572339078714);
  			\draw [line width=0.4pt] (-12.193777050744163,3.9585685777797375)-- (-7.818674409438614,8.054730801823446);
  			\draw [line width=0.4pt] (-10.993114152924402,8.827335190468446)-- (-7.818674409438614,8.054730801823446);
  			\draw [line width=0.8pt] (-10,6)-- (-7.178805936833635,4.989572339078714);
  			\draw [line width=0.4pt] (-10.993114152924402,8.827335190468446)-- (-10,3.0033185020760045);
  			\draw [line width=0.4pt] (-10,3.0033185020760045)-- (-7.178805936833635,4.989572339078714);
  			\draw [line width=0.4pt] (-7.818674409438614,8.054730801823446)-- (-7.178805936833635,4.989572339078714);
  			\draw (-10.332380725758481,7.179622843545514) node[anchor=north west] {$I$};
  			\draw (-12.9,4.2) node[anchor=north west] {$A$};
  			\draw (-11.599245687091022,9.7) node[anchor=north west] {$B$};
  			\draw (-7.8,8.8) node[anchor=north west] {$C$};
  			\draw (-7.063869125520526,5.5073610945865585) node[anchor=north west] {$D$};
  			\draw (-10.30704342653183,3) node[anchor=north west] {$E$};
  			\draw (-9.876309339678766,4.5) node[anchor=north west] {$70^\circ$};
  			\draw (-8.8,6.4) node[anchor=north west] {$80^\circ$};
  			\draw [line width=0.4pt] (-12.19377705074416,3.958568577779736)-- (-10,3.0033185020760023);
  			\begin{scriptsize}
  				\draw [color=black] (-10.006225730091389,6.006649689801592)-- ++(-1.5pt,-1.5pt) -- ++(3pt,3pt) ++(-3pt,0) -- ++(3pt,-3pt);
  			\end{scriptsize}
  		\end{tikzpicture}
  	\end{center}
  	\begin{enumerate}
  		\item Justifie que le triangle $ACD$ est rectangle.
  		\item Démontre que $mes\widehat{DAE}=mes\widehat{EBD}$ 
  		\item Détermine la mesure de chacun des angles suivants : $\widehat{BAD}$;  $\widehat{CBD}$ et  $\widehat{BID}$.
  	\end{enumerate}
  	
  	\begin{center}
  		\ovalbox{\textbf{Exercice 2}}\\
  	\end{center}
  	$\mathcal{(C)}$ est un cercle de centre $O$ et de rayon $3cm$. $A$ et $B$ sont deux points de ce cercle tels que $mes\widehat{BOA}=60^\circ$. Les tangentes au 
  	cercle $\mathcal{(C)}$ en $A$ et en $B$ se coupent en $T$. Les droites $(OT)$ et $(AB)$ se coupent en $I$.
  	\begin{enumerate}
  		\item Détermine la nature de chacun des triangles $AOB$ et $ATB$
  		\item  Calcule la mesure de l’angle $\widehat{BTA}$.
  		\item Démontre que la droite $(OT)$ est la médiatrice du segment $[AB]$.
  		\item  Démontre que les points $O$, $A$, $T$ et $B$ sont cocycliques. 
  	\end{enumerate}
   \situ{ CONFIGURATION DE L'ESPACE}
  	\underline{\textbf{Situation de départ}} \\
  	\textbf{Texte}: Le grenier Sossa
  	
  	\hspace{0.4cm} Anani, élève de la classe de 3ème est habitué à voir les 
  	agriculteurs de sa région construire des greniers de forme 
  	cylindrique en vue de conserver le maïs.
  	L’attention de Anani a été retenue par la façon dont 
  	l’agriculteur Sossa a construit son grenier : sur le sol il a tracé un 
  	cercle de 2 m de diamètre. Au centre de ce cercle, il a implanté 
  	verticalement un pieu de 3,5 m de haut. Une à une, plusieurs 
  	perches de même longueur sont disposées sur le cercle et 
  	s’appuient par leur extrémité sur un cerceau de 
  	30 cm de rayon dont le centre est au sommet du pieu. L’ensemble 
  	est fermé par un chapeau conique. 
  	Voici ci- contre une illustration du grenier de Sossa.
  	
  	\begin{center}
  		\begin{tikzpicture}[scale=0.5]
  			%\draw [thick,dashed] (0,0)  grid (5,5);
  			\draw [thick,dashed] (5,0) arc (0:180:2.5cm and 0.8cm);
  			\draw [thick] (0,0) arc (180:360:2.5cm and 0.8cm);
  			\draw [thick] (0,0)--(2.5,5)--(5,0);
  			\draw [thick,dashed] (3.5,3) arc (0:180:1cm and 0.4cm);
  			\draw [thick] (1.5,3) arc (180:360:1cm and 0.4cm);
  		\end{tikzpicture}
  	\end{center}	 
  	
  	
  	\Ovalbox{\textbf{Tâche}} : Tu vas te construire des connaissances nouvelles en 
  	mathématiques. 
  	
  	Pour cela, tu auras tout au long de la situation d’apprentissage à :\\ 
  	\ding{43} Exprimer ta perception de chacun des problèmes posés ;\\ 
  	\ding{43} Analyser chacun des problèmes posés ; \\
  	\ding{43} Mathématiser chacun des problèmes posés ; \\
  	\ding{43} Opérer sur l’objet mathématique que tu as identifié pour chaque problème ; \\
  	\ding{43} Améliorer au besoin ta production.\\
  	\ding{43} Réinvestir tes acquis dans d’autres situations.
  	\begin{center}
  		\section*{I-INTRODUCTION}
  	\end{center}  
  	\doublebox{\textbf{Activité 0}} \textbf{ Brainstorming}\\
  	\begin{enumerate}
  		\item  Lis le texte de la situation de départ 
  		\item  Reformule la situation-problème en tes 
  		propres termes 
  		\item  Formule toutes les idées ou questions que 
  		t’inspire la situation de départ.
  	\end{enumerate}
  	
  	\underline{Stratégie et durée}: 15 min
  	\begin{center}
  		\underline{\textbf{Éléments de réponses}}
  	\end{center}
  	(Voir cahiers de recherche.)
  	
  	\begin{center}
  		\section*{II- RÉALISATION}
  	\end{center}
  	\setcounter{sequen}{0}
  	\sequence{CÔNE}
  	\setcounter{q}{0}
  	\act{ Notion de cône}
  	\hspace*{0.4cm} Anani reconnaît la forme du chapeau du grenier de Sossa et en a fait une représentation en perspective cavalière que voici:
  	\begin{center}
  		\begin{tikzpicture}[scale=0.8]
  			%\draw [thick,dashed] (0,0)  grid (5,5);
  			\draw [thick,dashed] (5,0) arc (0:180:2.5cm and 0.8cm);
  			\draw [thick] (0,0) arc (180:360:2.5cm and 0.8cm);
  			\draw [thick] (0,0)--(2.5,5)--(5,0);
  			\draw [thick,dashed] (0,0)--(5,0);
  			\draw [thick,dashed] (2.5,0)--(2.5,5);
  			\draw [thick] (2.5,5)--(3.5,-0.75);	
  			\draw [thick] (2.5,0.3)--(2.8,0.3)--(2.8,0);		
  			\filldraw(2.5,0) node[below] {$O$}	;
  			\filldraw(2.5,5) node[above ] {$S$}	;
  			\filldraw(3.5,-0.75) node[below ] {$C$};
  			\filldraw(5,0) node[right] {$B$};
  			\filldraw(0,0) node[left] {$A$};
  			\filldraw(2.5,2) node[left] {$h$};
  			\filldraw(4,2.3) node[right] {$a$};		
  		\end{tikzpicture}
  	\end{center}	 
  	Il s'intéresse à ses aires et à son volume et désir réaliser son patron.
  	
  	\conAct{Découverte}
  	\begin{enumerate}
  		\item  Quel est le nom du solide illustré ?
  		\item  Représente-le en perspective cavalière à l'échelle de $\dfrac{1}{10}$ en prenant comme hauteur $\sqrt{13}cm$.
  		\item  Dis ce que représente :\\
  		\ding{43} le point $S$\\
  		\ding{43} les segments $[SA]$; $[SB]$ et $[SC]$;\\
  		\ding{43} la droite $(SO)$;\\
  		\ding{43} $h=SO$;\\
  		\ding{43} $a=SA=SB=SC$;\\
  		\ding{43} le disque de centre $O$ et de rayon $r=OB$.
  		\item  Justifie que $a^2=h^2+r^2$. 
  	\end{enumerate}
  	\dures
  	 
  	\ret{\ding{118}Un cône de révolution encore appelé cône circulaire droit est un solide de l'espace qui a:\\
  		\ding{43} un sommet;\\
  		\ding{43} une surface plane ayant la forme d'un disque appelée base;\\
  		\ding{43} une surface non plane appelée surface latérale.
  		
  		\ding{118} Une idée de cône de révolution nous est donnée en faisant tourner un triangle rectangle autour d'un axe supportant un côté de l'angle droit.  
  	}
  	
  	\nbVo{Rappel}{Si $a$ est l'apothème, $r$ est le rayon de base et $h$ est la hauteur d'un cône de révolution , on a :
  		
  		\ding{43} Aire de base: $A_{b} = \pi r^2$
  		
  		\ding{43}  Aire latérale: $A_{L}=\pi r a$ 
  		
  		\ding{43} Aire totale : $A_{T}=A_{L}+A_{b}= \pi r a+\pi r^2=\pi r(a+r)$
  		
  		\ding{43} Volume: $V=\dfrac{\pi r^2 \times h}{3}$
  		
  		\ding{43} $a^2=h^2+r^2$.
  		
  		\ding{43}Le développement de la surface latérale d'un cône de révolution est un secteur angulaire dont le rayon est l'apothème $\textbf{a} $ du cône et dont l'angle au sommet $\alpha$ est tel que $\alpha=\dfrac{360\textmd{\degre}\times r}{ a}$.}
  	
  	\conAct{Approfondissement 1}
  	Les longueurs des génératrices d'un cône de révolution est $9 cm$. La hauteur et le rayon de ce cône sont liés par la relation $\dfrac{h}{r}=2\sqrt{2}$. Détermine pour ce cône
  	\begin{enumerate}
  		\item la hauteur  $h$ et le rayon $r$
  		\item  l'aire totale et son volume
  		\item l'angle au sommet du du secteur angulaire  de ce cône puis réalise son patron à l'échelle de $ \dfrac{1}{2} $
  	\end{enumerate}. 
  	\dures
  	\conAct{Approfondissement 2}
  	Un cône de révolution a pour rayon $r=2\sqrt{3}cm$ et a pour hauteur $h=6cm$.
  	\begin{enumerate}
  		\item Détermine la capacité de ce cône.
  		\item 
  		\begin{enumerate}
  			\item  Calcule l'aire de la surface latérale de ce cône.
  			\item  Calcule l'aire de base de ce cône.
  			\item  Déduis-en l'aire totale de ce cône.
  		\end{enumerate}
  	\end{enumerate}
  	\dures
  	 
  	\begin{center}
  		\Ovalbox{\textbf{ Exercice de maison}}
  	\end{center}
  	\begin{center}
  		\Ovalbox{\textbf{ Exercice 1 }}
  	\end{center}
  	L'aire de la surface latérale d'un cône de révolution est le triple de celle de sa base.
  	\begin{enumerate}
  		\item Détermine la mesure $\alpha$  de l'angle  du secteur circulaire représentant la surface latérale du cône.
  		\item  Sachant que la hauteur $h$ de ce cône est égale à $12 cm$ :
  		\begin{enumerate}
  			\item  Calcule le rayon $r$ et l'apothème $a$ de ce cône.
  			\item  Calcule l'aire totale et le volume de ce cône.
  		\end{enumerate}
  	\end{enumerate}
  
  	\begin{center}
  		\Ovalbox{\textbf{ Exercice de maison}}
  	\end{center}
  	\begin{center}
  		\Ovalbox{\textbf{ Exercice 1 }}
  	\end{center}
  	L'aire de la surface latérale d'un cône de révolution est le triple de celle de sa base.
  	\begin{enumerate}
  		\item Détermine la mesure $\alpha$  de l'angle  du secteur circulaire représentant la surface latérale du cône.
  		\item  Sachant que la hauteur $h$ de ce cône est égale à $12 cm$ :
  		\begin{enumerate}
  			\item  Calcule le rayon $r$ et l'apothème $a$ de ce cône.
  			\item  Calcule l'aire totale et le volume de ce cône.
  		\end{enumerate}
  	\end{enumerate}
  	\begin{center}
  		\Ovalbox{\textbf{ Exercice 2 }}
  	\end{center}
  	Une pyramide régulière $SABCD$ de sommet $S$  et de base $ABCD$ a pour hauteur $21 cm$  et pour volume $847 cm^{3}$.
  	\begin{enumerate}
  		\item  Calcule la longueur du côté du carré de sa base .
  		\item  Calcule son apothème $a$.
  		\item Calcule son aire latérale.
  		\item Calcule son aire totale.
  	\end{enumerate}
  	
  	\begin{center}
  		\Ovalbox{\textbf{ Exercice 3 }}
  	\end{center}
  	Les questions sont indépendantes.
  	\begin{enumerate}
  		\item  Par quelle fraction doit-on multiplier le rayon de base d'un cône révolution de façon à ce qu'en multipliant la hauteur par 9, le volume reste le même. 
  		
  		\begin{minipage}{5cm}
  			\item Calcule la hauteur,  le rayon  et le volume du cône de révolution obtenu en faisant tourner le triangle $ABC$ autour de la droite $(BC)$ sachant que $BC=30cm$.
  		\end{minipage}
  		\begin{minipage}{5cm}
  			\begin{tikzpicture}[scale=0.6]
  				\draw [shift={(4,6)},line width=1.2pt] (0,0) -- (-126.86989764584402:0.9) arc (-126.86989764584402:-90:0.9) -- cycle;
  				\draw [line width=0.8pt] (4,1)-- (4,6);
  				\draw [line width=0.8pt] (4,6)-- (1,2);
  				\draw [line width=0.8pt] (1,2)-- (4,2);
  				\draw [line width=0.8pt] (4,2.5)-- (3.5,2.5);
  				\draw [line width=0.8pt] (3.5,2.5)-- (3.5,2);
  				\draw (3.9,2.5) node[anchor=north west] {$B$};
  				\draw (0.3,2.4) node[anchor=north west] {$A$};
  				\draw (3.5,6.9) node[anchor=north west] {$C$};
  				\draw (2.8,4.6) node[anchor=north west] {$30\textmd{\degre}$};
  				\draw (3.6,1.8) node[anchor=north west] {$\circlearrowleft$};
  			\end{tikzpicture}
  		\end{minipage}
  		
  		\begin{minipage}{5cm}
  			\item Calcule le volume du solide de révolution engendré par un carré de $5 cm$ de côté tournant autour d'une de ses diagonales (Tu prendras $\pi=3,14$ )
  		\end{minipage}
  		\begin{minipage}{4cm}
  			\begin{tikzpicture}[scale=0.6]
  				\draw [line width=1.2pt] (4,5)-- (6,7);
  				\draw [line width=1.2pt] (4.9350353948531795,6.06496460514682) -- (5.06496460514682,5.9350353948531795);
  				\draw [line width=1.2pt] (6,7)-- (8,5);
  				\draw [line width=1.2pt] (7.0649646051468205,6.06496460514682) -- (6.9350353948531795,5.9350353948531795);
  				\draw [line width=1.2pt] (8,5)-- (6,3);
  				\draw [line width=1.2pt] (7.0649646051468205,3.93503539485318) -- (6.9350353948531795,4.06496460514682);
  				\draw [line width=1.2pt] (6,3)-- (4,5);
  				\draw [line width=1.2pt] (4.9350353948531795,3.93503539485318) -- (5.06496460514682,4.06496460514682);
  				\draw [line width=0.8pt] (6,8)-- (6,2);
  				\draw [line width=0.8pt] (4.251407405999135,5.251407405999135)-- (4.484283228924636,4.990476542236579);
  				\draw [line width=0.8pt] (4.484283228924636,4.990476542236579)-- (4.236883685902223,4.763116314097777);
  				\draw (5.5,2.8) node[anchor=north west] {$\circlearrowleft$};
  			\end{tikzpicture}
  		\end{minipage}	
  	\end{enumerate}
   
  	\sequence{SECTION PLANE}
  	\setcounter{q}{0}
  	\act{Notion de section plane}
  	\hspace*{0.5cm} Anani décide de mettre en œuvre certaines caractéristiques du 
  	grenier de Sossa. Ainsi, il réalise les figures ci-dessous dont l’une 
  	est la représentation en perspective du chapeau du grenier 
  	coupé par un plan $\mathcal{(P)}$ parallèle au plan de sa base et l’autre est 
  	la représentation en perspective d’une pyramide régulière coupé 
  	par un plan $\mathcal{(P)}$ parallèle au plan de sa base.
  \end{multicols}
  
  \definecolor{ffffff}{rgb}{1,1,1}
  \definecolor{zzttqq}{rgb}{0.6,0.2,0}
  \begin{tikzpicture}[scale=0.6]
  	\fill[line width=2pt,color=zzttqq,fill=zzttqq,fill opacity=0.03] (-11.052201811652381,0.5367821746980063) -- (-6.138647431701272,0.550659108986587) -- (-4.34465327034393,2.344653270343928) -- (-9.235387015725221,2.3535969706251674) -- cycle;
  	\fill[line width=1.2pt,color=ffffff,fill=ffffff,fill opacity=0.14] (-9.003438016933359,1) -- (-6.858008450642274,1.006059154495914) -- (-5.999429059667906,1.864638545470281) -- (-8.138799471463077,1.864638545470281) -- cycle;
  	\fill[line width=0.8pt,color=ffffff,fill=ffffff,fill opacity=0.1] (-13,-6) -- (-4,-6) -- (-2,-2) -- (-11,-2) -- cycle;
  	\fill[line width=2pt,color=zzttqq,fill=zzttqq,fill opacity=0.02] (1.5943737545226166,0.6215983936426223) -- (7.549683454967029,0.5868735557391562) -- (9.25120051223686,2.323115450912458) -- (3.313253230744182,2.323115450912458) -- cycle;
  	\fill[line width=2pt,color=zzttqq,fill=zzttqq,fill opacity=0.04] (-0.009368600793884364,-6.054526279032764) -- (8.722019707498214,-6.070493396470353) -- (10.768285267741785,-1.9779622759832103) -- (2.0165550165331974,-2.0026790443786004) -- cycle;
  	\fill[line width=0.8pt,color=zzttqq,fill=zzttqq,fill opacity=0.04] (-13,-6) -- (-3.1838877119220834,-6.048721025885665) -- (-2,-2) -- (-11,-2) -- cycle;
  	\draw [line width=0.8pt,dash pattern=on 1pt off 1pt] (-9,-3)-- (-6,-5);
  	\draw [line width=0.8pt,dash pattern=on 1pt off 1pt] (-11,-5)-- (-4,-3);
  	\draw [line width=0.8pt,dash pattern=on 1pt off 1pt] (-7.5,-4)-- (-7.5,5.5);
  	\draw [line width=0.8pt] (-7.5,5.5)-- (-4,-3);
  	\draw [line width=0.8pt] (-7.5,5.5)-- (-6,-5);
  	\draw [line width=0.8pt,dash pattern=on 1pt off 1pt] (-7.5,5.5)-- (-9,-3);
  	\draw [line width=0.8pt] (-7.5,5.5)-- (-11,-5);
  	\draw [line width=0.8pt,color=zzttqq] (-11.052201811652381,0.5367821746980063)-- (-6.138647431701272,0.550659108986587);
  	\draw [line width=0.8pt,color=zzttqq] (-6.138647431701272,0.550659108986587)-- (-4.34465327034393,2.344653270343928);
  	\draw [line width=0.4pt,color=zzttqq] (-4.34465327034393,2.344653270343928)-- (-9.235387015725221,2.3535969706251674);
  	\draw [line width=0.8pt,color=zzttqq] (-9.235387015725221,2.3535969706251674)-- (-11.052201811652381,0.5367821746980063);
  	\draw [line width=0.8pt,color=ffffff] (-9.003438016933359,1)-- (-6.858008450642274,1.006059154495914);
  	\draw [line width=0.8pt,color=ffffff] (-6.858008450642274,1.006059154495914)-- (-5.999429059667906,1.864638545470281);
  	\draw [line width=0.8pt,color=ffffff] (-13,-6)-- (-4,-6);
  	\draw [line width=0.8pt,color=ffffff] (-4,-6)-- (-2,-2);
  	\draw [line width=0.8pt,color=ffffff] (-11,-2)-- (-13,-6);
  	\draw [rotate around={0:(5.370407472006344,-4.010355913823057)},line width=0.8pt] (5.370407472006344,-4.010355913823057) ellipse (2.6180020600442737cm and 1.6893592828040032cm);
  	\draw [line width=0.4pt] (5.402895777180931,5.608754933471161)-- (2.752761970585546,-3.9824752754810837);
  	\draw [line width=0.8pt] (5.402895777180931,5.608754933471161)-- (7.9883613758387195,-4.000109398599252);
  	\draw [line width=0.4pt] (5.402895777180931,5.608754933471161)-- (6.940846645602726,-5.362018533797325);
  	\draw [rotate around={-0.5583183624421316:(5.38985656254413,1.1031124838499535)},line width=0.8pt] (5.38985656254413,1.1031124838499535) ellipse (1.239085586525175cm and 0.16042788298019578cm);
  	\draw [line width=0.8pt,dash pattern=on 1pt off 1pt] (4.161258725288073,1.115084926684446)-- (6.618454399800344,1.0911400410154597);
  	\draw [line width=0.8pt,dash pattern=on 1pt off 1pt] (5.387720306248826,1.1031333012448732)-- (6.054399470343201,0.9613376605730773);
  	\draw [line width=0.8pt,color=zzttqq] (1.5943737545226166,0.6215983936426223)-- (7.549683454967029,0.5868735557391562);
  	\draw [line width=0.8pt,color=zzttqq] (7.549683454967029,0.5868735557391562)-- (9.25120051223686,2.323115450912458);
  	\draw [line width=0.8pt,color=zzttqq] (9.25120051223686,2.323115450912458)-- (3.313253230744182,2.323115450912458);
  	\draw [line width=0.8pt,color=zzttqq] (3.313253230744182,2.323115450912458)-- (1.5943737545226166,0.6215983936426223);
  	\draw [line width=0.8pt,color=zzttqq] (-0.009368600793884364,-6.054526279032764)-- (8.722019707498214,-6.070493396470353);
  	\draw [line width=0.8pt,color=zzttqq] (8.722019707498214,-6.070493396470353)-- (10.768285267741785,-1.9779622759832103);
  	\draw [line width=0.8pt,color=zzttqq] (10.768285267741785,-1.9779622759832103)-- (2.0165550165331974,-2.0026790443786004);
  	\draw [line width=0.8pt,color=zzttqq] (2.0165550165331974,-2.0026790443786004)-- (-0.009368600793884364,-6.054526279032764);
  	\draw [line width=0.8pt,dash pattern=on 1pt off 1pt] (2.752761970585546,-3.9824752754810837)-- (7.9883613758387195,-4.000109398599252);
  	\draw [line width=0.8pt,dash pattern=on 1pt off 1pt] (5.402895777180931,5.608754933471161)-- (5.45254713434296,-3.9915684738496475);
  	\draw [line width=0.8pt,dash pattern=on 1pt off 1pt] (5.45254713434296,-3.9915684738496475)-- (6.940846645602726,-5.362018533797325);
  	\draw [line width=2pt] (-6.858008450642273,1.0060591544959137)-- (-9.003438016933359,1);
  	\draw [line width=2pt] (-6.858008450642273,1.0060591544959137)-- (-5.999429059667906,1.864638545470281);
  	\draw [line width=1.2pt,dash pattern=on 1pt off 1pt] (-5.999429059667906,1.864638545470281)-- (-8.138799471463077,1.864638545470281);
  	\draw [line width=1.2pt,dash pattern=on 1pt off 1pt] (-8.138799471463077,1.864638545470281)-- (-9.003438016933359,1);
  	\draw [line width=0.8pt] (-11,-5)-- (-6,-5);
  	\draw [line width=0.8pt] (-6,-5)-- (-4,-3);
  	\draw [line width=1.2pt,dash pattern=on 1pt off 1pt] (-4,-3)-- (-9,-3);
  	\draw [line width=1.2pt,dash pattern=on 1pt off 1pt] (-9,-3)-- (-11,-5);
  	\draw [line width=0.8pt,color=zzttqq] (-13,-6)-- (-3.1838877119220834,-6.048721025885665);
  	\draw [line width=0.8pt,color=zzttqq] (-3.1838877119220834,-6.048721025885665)-- (-2,-2);
  	\draw [line width=0.8pt,color=zzttqq] (-2,-2)-- (-11,-2);
  	\draw [line width=0.8pt,color=zzttqq] (-11,-2)-- (-13,-6);
  	\draw [line width=1.2pt,dash pattern=on 1pt off 1pt] (-8.140742300143643,1.8691269658526934)-- (-6.858008450642273,1.0060591544959137);
  	\draw [line width=1.2pt,dash pattern=on 1pt off 1pt] (-8.999996760383807,1.0000097188485797)-- (-6.002019718165304,1.8620478869728814);
  	\draw [line width=0.8pt] (-8.985529408392216,-6.394533579553936)-- (-6.820716071714553,-6.399792508027989);
  	\draw [line width=0.8pt] (-8.982492134065824,-8.15598445226318)-- (-6.811717667091587,-8.186775579454302);
  	\draw [line width=0.8pt] (-8.985529408392216,-6.394533579553936)-- (-8.982492134065824,-8.15598445226318);
  	\draw [line width=0.8pt] (-6.820716071714553,-6.399792508027989)-- (-6.811717667091587,-8.186775579454302);
  	\draw [line width=1.2pt,dash pattern=on 1pt off 1pt] (-8.998288674427224,1.0051339767183267)-- (-8.985529408392216,-6.394533579553936);
  	\draw [line width=1.2pt,dash pattern=on 1pt off 1pt] (-6.858008450642275,1.0060591544959125)-- (-6.820716071714553,-6.399792508027989);
  	\draw [line width=0.8pt] (-11,-9.94246507628567)-- (-6,-10.001714871312494);
  	\draw [line width=0.8pt] (-11,-15)-- (-6,-15);
  	\draw [line width=1.2pt,dash pattern=on 1pt off 1pt] (-11,-5)-- (-11,-9.94246507628567);
  	\draw [line width=1.2pt,dash pattern=on 1pt off 1pt] (-6,-5)-- (-6,-10.001714871312494);
  	\draw [line width=0.8pt] (-11,-9.94246507628567)-- (-11,-15);
  	\draw [line width=0.8pt] (-6,-10.001714871312494)-- (-6,-15);
  	\draw [line width=0.8pt] (5.37505079538736,-7.619736237782518) circle (1.2711599258284343cm);
  	\draw [line width=0.8pt] (5.483234618862127,-12.650284029358872) circle (2.861508144451029cm);
  	\draw [line width=0.8pt,dash pattern=on 1pt off 1pt] (2.752761970585546,-3.9824752754810837)-- (2.8138759613587063,-11.619384801320034);
  	\draw [line width=0.8pt,dash pattern=on 1pt off 1pt] (7.9883613758387195,-4.000109398599252)-- (7.996837511301049,-11.282793861247159);
  	\draw [line width=0.8pt,dash pattern=on 1pt off 1pt] (2.6217722352735624,-12.666466985430665)-- (8.344741630066874,-12.652830710660519);
  	\draw [line width=0.8pt,dash pattern=on 1pt off 1pt] (5.483234618862127,-12.650284029358872)-- (7.2649232034247575,-14.889438925385225);
  	\draw [line width=0.8pt,dash pattern=on 1pt off 1pt] (4.1038923926895325,-7.617768426089483)-- (6.646146134990301,-7.606922390903581);
  	\draw [line width=0.8pt,dash pattern=on 1pt off 1pt] (6.054399470343201,0.9613376605730773)-- (6.122024642307467,-8.648268020981037);
  	\draw [line width=0.8pt,dash pattern=on 1pt off 1pt] (6.622714145889668,1.0753087252622615)-- (6.515496103204996,-7.058280935101274);
  	\draw [line width=0.8pt,dash pattern=on 1pt off 1pt] (4.156849980560731,1.0991290207950106)-- (4.274094414426187,-6.9843288060876825);
  	\draw [line width=0.8pt,dash pattern=on 1pt off 1pt] (5.452547134342961,-3.9915684738496484)-- (5.483234618862127,-12.650284029358872);
  	\draw [line width=0.8pt,dash pattern=on 1pt off 1pt] (5.465378218286185,-7.611959909447801)-- (6.122024642307467,-8.648268020981037);
  	\draw [->,line width=0.8pt] (-1.2006136673959062,0.20814161325148525) -- (-4.7040039462370125,1.8067760123342225);
  	\draw [->,line width=0.8pt] (-1.2006136673959062,0.20814161325148525) -- (2.064682126475416,1.3305870423946837);
  	\draw [->,line width=0.8pt] (-1.6427891394826477,-7.274827914369838) -- (-4.057747487033313,-6.288436476637936);
  	\draw [->,line width=0.8pt] (-1.6427891394826477,-7.274827914369838) -- (0.6701287145095387,-6.186395983079464);
  	\draw (-1.5,0.2) node[anchor=north west] {$(\mathcal{P})$};
  	\draw (-3.1,-7.4) node[anchor=north west] {Plan de la base};
  	\draw (-11.9,-4.4) node[anchor=north west] {$A$};
  	\draw (-5.6,-4.6) node[anchor=north west] {$B$};
  	\draw (-3.9,-2.4) node[anchor=north west] {$C$};
  	\draw (-10,-2.4) node[anchor=north west] {$D$};
  	\draw (-10,1.7) node[anchor=north west] {$A'$};
  	\draw (-6.7,1.45) node[anchor=north west] {$B'$};
  	\draw (-6,2.4) node[anchor=north west] {$C'$};
  	\draw (-9.2,2.4) node[anchor=north west] {$D'$};
  	\draw (-8.2,6) node[anchor=north west] {$S$};
  	\draw (-11,-15.5) node[anchor=north west] {Figure 1};
  	\draw (4,-16) node[anchor=north west] {Figure 2};
  	\draw [line width=0.8pt,dash pattern=on 1pt off 1pt] (6.940846645602726,-5.362018533797325)-- (7.2649232034247575,-14.889438925385225);
  \end{tikzpicture}
  
  \begin{multicols*}{2}
  	\conAct{Section d’un cône de révolution et d’une pyramide régulière}	
  	\begin{enumerate}
  		\item  
  		\begin{enumerate}
  			\item Précise la nature de la section obtenue pour la pyramide régulière.
  			\item  Pour le cas de la pyramide, les côtés $[AB]$ et $[A’B’]$ par exemple sont appelés côtés correspondants.
  			
  			Donne tous les autres côtés correspondants de cette figure.
  		\end{enumerate}
  		\item Reproduis en perspective cavalière la figure 1 sachant que $AB=6cm$, $A'B'=4cm$, $\alpha=45\up{°}$ et $c=0,5$.
  	\end{enumerate}
  	\dures
  	 \newpage
  	 	\SetWatermarkText{\resizebox{28cm}{0.35cm}{Support d'activités 3\up{ème} }}
  	 \SetWatermarkScale{1}
  	 \SetWatermarkColor{black!30}
  	 \SetWatermarkAngle{90}
  	\propr{La section d’une pyramide régulière par un plan parallèle à la base est un polygone de même nature que cette base ; les supports des côtés correspondants de ces polygones sont parallèles.
  	}
  	\nbVo{Remarque}{La section d’une pyramide régulière par un plan parallèle à la base permet d'obtenir d'un côté une petite  pyramide régulière appelée \textbf{pyramide réduite} et de l'autre un \textbf{tronc de pyramide}.
  	}
  	\conAct{Section d’un cône de révolution}	
  	\begin{enumerate}
  		\item  Précise la nature de la section obtenue pour le cône de révolution.
  		\item Reproduis en perspective cavalière un cône  sachant que le rayon du disque de base est $r=3cm$ et coupé par un plan parallèle à sa base.
  	\end{enumerate}
  	\dures
  	\propr{La section d’un cône de révolution par un plan parallèle à celui de la base est un cercle.
  	}
  	\nbVo{Remarque}{La section d’un cône de révolution par un plan parallèle à la base permet d'obtenir d'un côté un cône de révolution  appelé \textbf{cône réduit} et de l'autre un \textbf{tronc de cône }.
  	}
  	\act{Échelle de réduction, rapport d'aires et rapport de volume}
  	
  	On considère le cône de révolution ci-après coupé par un plan parallèle à sa base.\\
  	\begin{tikzpicture}[scale=0.7]
  		%\draw [thick,dashed] (0,0)  grid (6,7);
  		\draw [thick,dashed] (6,0) arc (0:180:3cm and 0.8cm);
  		\draw [thick] (0,0) arc (180:360:3cm and 0.8cm);
  		\draw [thick] (0,0)--(3,7)--(6,0);
  		\draw [thick,dashed] (3,0)--(6,0);
  		\draw [thick,dashed] (3,0)--(3,7);
  		\draw [thick,dashed] (3,4)--(4.28,4);
  		\draw [thick] (3,4.2)--(3.2,4.2)--(3.2,4);
  		\draw [thick] (3,0.3)--(3.3,0.3)--(3.3,0);
  		\draw [thick,dashed] (4.27,4) arc (0:180:1.28cm and 0.6cm);
  		\draw [thick] (1.73,4) arc (180:360:1.28cm and 0.6cm);	
  		\filldraw(3,4) node[left] {$O'$};		
  		\filldraw(3,0) node[left] {$O$}	;
  		\filldraw(3,7) node[above ] {$S$};
  		\filldraw(6,0) node[right] {$A$};
  		\filldraw(4.3,4) node[right] {$A'$};
  		\filldraw(4,0) node[above] {$r$};
  		\filldraw(3.5,3.89) node[above] {$r'$};
  		\filldraw(1.3,4) node[left] {$a$};
  		\filldraw(4,5.2) node[right] {$a'$};
  		\filldraw(6.7,2) node[left] {$h_{t}$};
  		\filldraw(6.7,5) node[left] {$h'$};
  		\filldraw(7,4) node[right] {$h$};
  		\filldraw(-0.5,1.7) node[left] {\rotatebox{90}{Tronc de cône}};	
  		\filldraw(-0.5,5.4) node[left] {\rotatebox{90}{ cône réduit}};	
  		\draw [<->,line width=1pt] (-0.3,0)--(2.7,7);
  		\draw [<->,line width=1pt] (3.2,7)--(4.5,4);
  		\draw [<->,line width=1pt] (6.7,0)--(6.7,4);
  		\draw [<->,line width=1pt] (6.7,4)--(6.7,7);
  		\draw [<->,line width=1pt] (7,0)--(7,7);
  		\draw [<->,line width=1pt] (-0.5,0)--(-0.5,4);
  		\draw [<->,line width=1pt] (-0.5,4)--(-0.5,7);	
  	\end{tikzpicture}
  	
  	\conAct{Échelle de réduction}
  	\begin{enumerate}
  		\item Justifie que $(O'A') \parallel (OA)$.
  		\item  Démontre que $\dfrac{SO'}{SO}=\dfrac{SA'}{SA}=\dfrac{O'A'}{OA}$ puis déduis-en qu'il existe un nombre réel positif $k$ tel que  $k=\dfrac{r'}{r}=\dfrac{a'}{a}=\dfrac{h'}{h}$.
  		\item Démontre que $h_{t}=h\left(1-k\right)$.
  	\end{enumerate}
  	\dures
  	 
  	\ret{L’échelle $k$ de réduction est le rapport d’une longueur du solide réduit à la longueur correspondante du solide initial.
  		
  		$k=\dfrac{r'}{r}=\dfrac{a'}{a}=\dfrac{h'}{h}$ où $r$; $a$ et $h$ désignent respectivement le rayon du disque base, l'apothème et la hauteur du solide initial et $r'$; $a'$ et $h'$ désignent respectivement le rayon du disque base, l'apothème et la hauteur du solide réduit.
  		
  		$\bullet$ La hauteur du tronc de cône est donnée par la formule $h_{t}=h\left(1-k\right)$.
  	}
  	
  	\conAct{Approfondissement}
  	\hspace*{0.5cm} La section d'un cône de révolution de hauteur $h=21m$ par un plan parallèle à la base a permis d'obtenir un cône réduit de hauteur $h'=3m$.
  	\begin{enumerate}
  		\item Détermine l'échelle de réduction $k$ issue de cette section.
  		\item Détermine la hauteur du tronc de cône obtenu.
  	\end{enumerate}
  	\dures
  	 
  	\conAct{Rapport d'aires}
  	
  	On considérant le cône de révolution de l'activité 2 et en désignant  respectivement par $A_{b}$ et $A_{l}$ l'aire de base et l'aire latérale du cône initial et par  $A'_{b}$ et $A'_{l}$ l'aire de base et l'aire latérale du cône réduit:
  	\begin{enumerate}
  		\item Exprime en fonction de $k$ les rapport $\dfrac{A'_{b}}{A_{b}}$ et $\dfrac{A'_{l}}{A_{l}}$  
  		\item On désigne par $AL_{t}$ l'aire latérale du tronc de cône. Exprime $AL_{t}$  en fonction de $k$ et $A_{l}$. 
  	\end{enumerate}
  	\dures
  	 
  	 \propr{Soit $k$ l'échelle de réduction d'un cône
  		
  		$\bullet$ Si $A_{b}$; $A_{l}$ et $A_{t}$ désignent  respectivement  l'aire de base;  l'aire latérale et l'aire totale du cône initial et $A'_{b}$; $A'_{l}$ et $A'_{t}$ désignent  respectivement  l'aire de base;  l'aire latérale et l'aire totale du cône réduit,  alors: 
  		
  		$k^2=\dfrac{A'_{b}}{A_{b}}=\dfrac{A'_{l}}{A_{l}}=\dfrac{A'_{t}}{A_{t}}$
  		
  		$\bullet$ Si $AL_{t}$ désigne l'aire latérale du tronc de cône, alors:  $AL_{t}=A_{l}\left(1-k^2\right)$.
  	}
  	
  	\nbVo{Remarque}{ Ces formules sont valables pour une pyramide régulière. 
  		
  	} 
  	\conAct{Approfondissement } 
  	Une pyramide régulière  d'aire de base $A_{b}=95dm^2$ a été sectionnée par un plan parallèle au plan de sa base selon l'échelle de réduction $k=\dfrac{2}{3}$. L'aire latérale de la pyramide réduite est $A'_{l}=250 dm^2$ 
  	\begin{enumerate}
  		\item Calcule l'aire latérale du tronc de pyramide obtenu 
  		\item Détermine l'aire de base de la pyramide réduite.  
  	\end{enumerate}
  	\dures
  	 
  	\conAct{Rapport de volume } 
  	On considérant le cône de révolution de l'activité 2 et en désignant  respectivement par $V$ et $V'$ le volume du cône initial et le volume du cône réduit. 
  	\begin{enumerate}
  		\item Exprime en fonction de $k$, le rapport $\dfrac{V'}{V}$. 
  		\item Détermine l'expression du volume  $V_{t}$ du tronc de cône en fonction de $V$ et $K$. 
  	\end{enumerate}
  	\dures
  	 
  	\propr{Soit $k$ l'échelle de réduction d'un cône
  		
  		Si $V$; $V'$ et $V_{t}$ désignent  respectivement  le volume du cône  initial ; le volume du  cône réduit et le volume tronc de cône, alors : 
  		
  		$\bullet$ $\dfrac{V'}{V}=k^3$ 
  		
  		$\bullet$  $V_{t}=V\left(1-k^3\right)$
  	}
  	\nbVo{Remarque}{ Ces formules sont valables pour une pyramide régulière. 
  		
  	}
  	
  	\conAct{Approfondissement} 
  	
  	Un cône circulaire droit de volume  $V=42 m^3$ a été sectionnée par un plan parallèle au plan de sa base selon l'échelle de réduction $k=\dfrac{1}{2}$.
  	
  	\begin{enumerate}
  		\item Calcule le volume du cône réduit. 
  		\item Calcule le volume du tronc de cône de deux manières différentes. 
  	\end{enumerate}
  	\dures
  
  	\act{Autre formule du volume du tronc d’un cône ou d'une pyramide régulière}
  	
  	Anani désire déterminer le volume du tronc d'une pyramide régulière et d'un tronc cône de révolution. 
  	
  	\begin{tikzpicture}[scale=0.6]
  		%\draw [thick,dashed] (0,0)  grid (14,7);
  		\draw [thick,dashed] (6,0) arc (0:180:3cm and 0.8cm);
  		\draw [thick] (0,0) arc (180:360:3cm and 0.8cm);
  		\draw [thick] (0,0)--(3,7)--(6,0);
  		\draw [thick,dashed] (3,0)--(6,0);
  		\draw [thick,dashed] (3,0)--(3,7);
  		\draw [thick,dashed] (3,4)--(4.28,4);
  		\draw [thick] (3,4.2)--(3.2,4.2)--(3.2,4);
  		\draw [thick] (3,0.3)--(3.3,0.3)--(3.3,0);
  		\draw [thick,dashed] (4.27,4) arc (0:180:1.28cm and 0.6cm);
  		\draw [thick] (1.73,4) arc (180:360:1.28cm and 0.6cm);	
  		\filldraw(3,4) node[left] {$O'$};		
  		\filldraw(3,0) node[left] {$O$}	;
  		\filldraw(3,7) node[above ] {$S$};
  		\filldraw(5.8,0) node[right] {$A$};
  		\filldraw(4.3,4) node[right] {$A'$};
  		\filldraw(4,0) node[above] {$r$};
  		\filldraw(3.5,3.89) node[above] {$r'$};
  		\filldraw(6.7,2) node[left] {$h_{t}$};
  		\filldraw(6.7,5) node[left] {$h'$};
  		\filldraw(7,4) node[right] {$h$};
  		\draw [<->,line width=1pt] (6.7,0)--(6.7,4);
  		\draw [<->,line width=1pt] (6.7,4)--(6.7,7);
  		\draw [<->,line width=1pt] (7,0)--(7,7);
  		
  		\draw [thick,dashed] (13,1)--(8,1);
  		\draw [thick] (8,1)--(9,-2)--(13,1);
  		\draw [thick,dashed] (10,0)--(10,4);
  		\draw [thick] (8.5,5)--(9.5,3)--(12,5);	
  		\draw [thick,dashed] (8.5,5)--(12,5);
  		\draw [thick] (8.5,5)--(8,1);
  		\draw [thick] (12,5)--(13,1);
  		\draw [thick] (9.5,3)--(9,-2);
  		\filldraw(10,4) node[left] {$O'$};		
  		\filldraw(10,0) node[left] {$O$}	;
  	\end{tikzpicture}
  	On désigne par $B$ et $b$ les aires de bases parallèles respectivement du solide initial et du solide réduit; par $h'$ la hauteur du solide réduit, $h$ celui du solide initial et $h_{t}$ la hauteur du tronc du solide.
  	
  	\conAct{Autre formule}
  	Soit $V_{t}$ d’un tronc de pyramide régulière ou d'un tronc de cône de révolution.
  	\begin{enumerate}
  		\item Exprime $h$ en fonction de $h'$  et de $h_{t}$.
  		\item  Démontre que $V_{t}=\dfrac{1}{3}B\left(h'+h_{t}\right)-\dfrac{1}{3}bh'$
  	\end{enumerate}
  	\dures
  	 
  	\nbVo{Remarque}{Le volume $V_{t}$ d’un tronc d'une pyramide régulière ou d'un tronc de cône de révolution  connaissant les aires $B$ et $b$ des bases parallèles et la hauteur $h_{t}$ du tronc est donné par la formule: 
  		
  		$V_{t}=\dfrac{1}{3}B\left(h'+h_{t}\right)-\dfrac{1}{3}bh'$
  	}
  	\conAct{Autre formule}
  	
  	Soit $V_{t}$ d’un tronc de pyramide régulière ou d'un tronc de cône de révolution.
  	\begin{enumerate}
  		\item Démontre que
  		\begin{enumerate}
  			\item  $h'=\dfrac{kh_{t}}{1-k}$
  			\item  $h'=\dfrac{h_{t}\left(\sqrt{Bb}+b\right)}{B-b}$
  		\end{enumerate}
  		\item Déduis-en que $V_{t}=\dfrac{h_{t}}{3}\left(B+\sqrt{Bb}+b\right)$
  		\item soit r et r' les rayons respectifs de la base du cône initial et de la base du cône réduit. Démontre que $V_{t}=\dfrac{\pi h_{t}}{3}\left(r^2+rr'+r'^2\right)$
  	\end{enumerate}
  	\dures
  	 
  	\propr{$\bullet$ le volume $V_{t}$ d'un tronc de cône de révolution ou d'un tronc de pyramide régulière de hauteur $h_{t}$ est donné par:  $V_{t}=\dfrac{h_{t}}{3}\left(B+\sqrt{Bb}+b\right)$, où où $B$ et $b$ sont les aires des bases 
  		
  		$\bullet$ Dans le cas d'un tronc de cône de révolution de hauteur $h_{t}$, ce volume est	
  		$V_{t}=\dfrac{\pi h_{t}}{3}\left(r^2+rr'+r'^2\right)$, où $r$ et $r'$ sont les rayons des deux bases.
  	}
  	\reto
  	\begin{center}
  		\Ovalbox{\textbf{ Exercice de maison}}
  	\end{center}
  	\begin{center}
  		\Ovalbox{\textbf{ Exercice 1 }}
  	\end{center}
  	
  	Un cône de révolution a pour volume $9400 cm^3$. L'aire de sa base est égale $705 cm^2$. La section à mis hauteur dudit cône donne deux solides.
  	\begin{enumerate}
  		\item Donne le nom  de chaque solide.
  		\item Précise l'échelle de réduction.
  		\item Calcule l'aire de base du petit cône.
  		\item  Calcule le volume $V_{t}$ du grand cône.
  		\item   Calcule le volume $V'$ du petit cône puis déduis-en sa hauteur $h'$.  
  	\end{enumerate}
  	\begin{center}
  		\Ovalbox{\textbf{ Exercice 2 }}
  	\end{center}
  	Dans son programme d’extension en eau potable dans certaines zone reculées de la ville de TOFFO, le DG de la SONEB a prévu la construction d’un château d’eau dans la dite commune dont le plan est le suivant:
  	
  	\begin{tikzpicture}[scale=0.6]
  		%	\draw[dashed] (0,0) grid (18,15);
  		\draw [thick,dashed] (8,4) arc (0:180:2cm and 0.4cm);
  		\draw [thick] (8,4) arc (0:-180:2cm and 0.4cm);
  		
  		\draw [thick,dashed] (8,0) arc (0:180:2cm and 0.4cm);
  		\draw [thick] (8,0) arc (0:-180:2cm and 0.4cm);
  		\draw [thick] (2,9)--(4,4);
  		\draw [thick] (8,4)--(10,9);
  		\draw [thick] (10,9) arc (0:180:4cm and 0.6cm);
  		\draw [thick] (10,9) arc (0:-180:4cm and 0.6cm);
  		\draw [thick] (4,4)--(4,0);	
  		\draw [thick] (8,4)--(8,0);		
  		\draw [thick,dashed] (6,9)--(6,0)--(8,0);
  		\draw [thick] (10,9)--(2,9);
  		\draw [thick,dashed] (6,4)--(8,4);
  		\draw [<->,line width=1.5pt] (11,9)--(11,4);
  		\draw [<->,line width=1.5pt] (11,0)--(11,4);
  		\filldraw (10,9) node[right] {\sf $A$};
  		\filldraw (6,8.9) node[above] {\sf$O$ };
  		\filldraw (2,9) node[left] {\sf $B$};
  		\filldraw (6,4) node[left] {\sf $O'$};
  		\filldraw (8,4) node[right] {\sf $A'$};
  		\filldraw (6,0) node[left] {\sf $S$};
  		\filldraw(11,6.5) node[right] {\sf\textbf{ Partie $I$}};
  		\filldraw(11,2) node[right] {\sf\textbf{ Partie $II$}};
  	\end{tikzpicture}
  	
  	$OA=9\hspace{0.1cm}m$; $O'A'=3\hspace{0.1cm}m$;\\ $O'S=3,5\hspace{0.1cm}m$; $O'O= 12\hspace{0.1cm}m$\\ et $(O'A')\lVert(OA)$.\\
  	On prendra  $\pi=3,14$
  	\begin{enumerate}
  		\item 
  		\begin{enumerate}
  			\item  Propose un nom à chacun des solides représentés par les parties \textbf{I} et \textbf{II} du château.
  			\item   Détermine le coefficient de réduction $k$ du cône de révolution qui a permis d’obtenir le solide 
  			représentant la partie I de ce château. 
  		\end{enumerate}	
  		\item En désignant par $H$ la hauteur du cône initial dont la section a permis d’obtenir la partie $I$ du château d’eau et par $h$ celle du cône réduit, établis la formule permettant de calculer la hauteur $H$ puis justifie que $H=18\hspace{0.1cm}m$.
  		\item Calcule l’aire de la superficie latérale de ce château pour un prévision de badigeonnement.
  		\item 
  		\begin{enumerate}
  			\item  Détermine le volume $V_{_{i}}$ du cône initial dont la section a permis d’obtenir la partie \textbf{I} du château.
  			\item  Déduis-en la capacité d’eau en litre lorsque le château d’eau est plein
  			d’eau.
  		\end{enumerate}	
  	\end{enumerate}	
   
  	\section*{III- RETOUR ET PROJECTION}
  	\doublebox{\textbf{Activité 1 }} \textbf{Objectivation/auto-évaluation}
  	
  	\begin{enumerate}
  		\item  Fais le point de tout ce que tu as appris sur chaque séquence de cette SA. 
  		\item  Fais le point de tes difficultés et de tes réussites au cours de l'apprentissage.
  		\item  Identifie des situations de la vie courante où tu peux utiliser tes acquis
  	\end{enumerate}
  	\doublebox{\textbf{Activité 2}} \textbf{Réinvestissement}
  	
  	(Séance d'exercice)
  	\begin{center}
  		\it\textit{\textbf{Fin de la $SA_{2} $!!!}}
  	\end{center}
   \situ{ CALCUL LITTERAL}
  	\underline{\textbf{Situation de départ}} \\
  	\textbf{Texte}:	
  	
  	Le conseil communal de Zodji a décidé de faire aménager les abords de la rivière Zo afin de créer sur ses deux rives un espace de jeu. Pour ce faire, il lui faut après nettoyage des deux rives procéder à l’épandage d’un herbicide à raison de 4 litres sur $100 m^2$ pour une superficie de plus de $20 ha$. Cet herbicide est vendu dans un magasin de la place à $3.500 F$ le tube de un litre. Comlan, fils du comptable de la commune est un élève de la classe de troisième au CEG de Zodji. Des informations qu’il a reçues de son professeur de Sciences Physiques et des agents commerciaux qu’il a rencontrés il ressort que :
  	
  	\ding{43} la fabrication de l’herbicide nécessite le mélange de 2 produits $A$ et $B$ qui sont vendus respectivement à $900F$ la bouteille de un litre et à $5000 F$ le bidon de 5 litres.
  	
  	\ding{43} pour être efficace ce mélange doit être constitué d’au moins $40\%$ du produit $A$ et d’au moins $50 \%$ du produit $B$. les deux produits, après achat, doivent passer au contrôle de qualité dans une machine qui ne peut travailler plus de 6 heures par jour.
  	
  	\ding{43} le temps disponible pour le contrôle de qualité ne peut excéder 5 jours. 
  	
  	\ding{43} le produit A nécessite 4 min de contrôle par litre et le produit B 15 min de contrôle par bidon.
  	
  	\ding{43} le contrôle d’une bouteille du produit A coûte 90 F et celui d’un bidon du produit B coûte 500 F. 
  	Avant l’épandage l’herbicide devra être dilué dans de l’eau. 
  	
  	Le comptable a le choix entre acheter l’herbicide dans le commerce et le faire fabriquer.
  	
  	\Ovalbox{\textbf{Tâche}} : Tu vas te construire des connaissances nouvelles en mathématiques.
  	
  	Pour cela, tu auras tout au long de la situation d’apprentissage à :\\ 
  	\ding{43} Exprimer ta perception de chacun des problèmes posés ;\\ 
  	\ding{43} Analyser chacun des problèmes posés ; \\
  	\ding{43} Mathématiser chacun des problèmes posés ; \\
  	\ding{43} Opérer sur l’objet mathématique que tu as identifié pour chaque problème ; \\
  	\ding{43} Améliorer au besoin ta production.\\
  	\ding{43} Réinvestir tes acquis dans d’autres situations.
  	
  	\doublebox{\textbf{Activité 0}} \textbf{ Brainstorming}\\
  	\begin{enumerate}
  		\item  Lis le texte de la situation de départ 
  		\item  Reformule la situation-problème en tes 
  		propres termes 
  		\item  Formule toutes les idées ou questions que 
  		t’inspire la situation de départ.
  	\end{enumerate}
  	
  	\underline{Stratégie et durée}: 15 min
  	\begin{center}
  		\underline{\textbf{Éléments de réponses}}
  	\end{center}
  	(Voir cahiers de recherche.)
  		\setcounter{sequen}{0}
  		\setcounter{q}{0}
  \sequence{POLYNÔMES}
  	\act{ Notion de monômes et polynômes}
  	
  	Comlan veut aider son père à faire le choix entre acheter l'herbicide  dans le commerce et le faire fabriquer. Il désigne par $x$ le prix d'un litre du produit $B$.
  	
  	\conAct{Monômes et polynômes}
  	\begin{enumerate}
  		\item 
  		\begin{enumerate}
  			\item  Donne l'expression littérale qui traduit le prix d'un bidon du produit $B$.
  			
  			\bcinfo{\textbf{Information}: \textbf{Cette expression littérale est appelée monôme de variable $x$, de degré 1 et coefficient 5.}}
  			\item  Comment peux-tu appeler les expressions $4x$; $\sqrt{3}y^6$ et $28t^7$.
  		\end{enumerate}
  		\item Donne un exemple de monôme :
  		\begin{enumerate}
  			\item de variable $x$ et de degré 1
  			\item  de variable $x$ et de degré 3
  			\item  de variable $x$ et de degré 6
  		\end{enumerate}
  		\item  Donne l'expression littérale $P(x)$ égale à la somme des monômes trouvés à la question 2. 
  		
  		\bcinfo{\textbf{Information}: \textbf{Cette expression littérale est appelée polynôme}}
  		\item On donne le polynôme suivant $Q(x)=3-x^{2026}+x^{9}$ 
  		\begin{enumerate}
  			\item Cite tous les monômes de ce polynôme tout en précisant le coefficient et le degré de chacun d'eux.
  			\bcinfo{\textbf{Information}: \textbf{Le degré le plus élevé des monômes d'un polynôme est le degré de ce polynôme.}}
  			\item  Donne les degrés des polynômes $P(x)$ et $Q(x)$.
  		\end{enumerate}
  	\end{enumerate}
  	\dures
   	\defs{$\bullet$ On appelle monôme toute expression de la forme $ax^n$ où $a$ est un nombre réel non nul appelé coefficient, $n$ un entier naturel appelé degré et $x$ est la variable ou l'inconnue.
  		
  		$\bullet$ On appelle polynôme, toute somme algébrique de deux ou plusieurs monômes.  }
  	\nbVo{Remarques}{\ding{43} Le degré d'un polynôme est le plus grand degré des degrés de ces  monômes  
  		
  		\ding{43} Un binôme est la somme de deux monômes.
  		
  		\ding{43} Un trinôme est la somme de trois monômes. }
  	
  	\act{ Calculs sur les polynômes}
  	
  	\conAct{Somme, produit, multiplication et développement  de polynômes}
  	
  	On considère les expressions suivantes: $A(x)=2x^3+x-1$; $B(x)=3x+1$; $C(x)=x-2$; $D(x)=\left(2x-5\right)^2-3x\left(4x^2-5x\right)$; $E(x)=\left(x-1\right)^2$ et $F(x)=(x-\sqrt{2})(x+\sqrt{2})$
  	\begin{enumerate}
  		\item Calcule $A(x)+B(x)$, $A(x)-B(x)$  et $B(x)\times C(x)$
  		\item Développe, réduis et ordonne $D(x)$, $E(x)$ et $F(x)$ suivant les puissances décroissantes de $x$. 
  	\end{enumerate}
  	\dures
  	 
  	\ret{$\bullet$ Développer un polynôme, c'est passer de sa forme produit à la forme algébrique de monômes.
  		
  		$\bullet$ Ordonner un polynôme, c'est le ranger suivant les puissances croissantes ou décroissantes de la variable. }
  	
  	\conAct{Factorisations}
  	
  	On considère les expressions suivantes:
  	
  	$A(x)=(3x-2)(5x+1)-(9x^2-4)$;
  	
  	$B(x)=(x-\sqrt{5})(2x-1)+(x^2-2\sqrt{5}x+5)$;
  	
  	$C(x)= (2x+3)^2-(x-4)^2$;
  	
  	$D(x)= 9x^2(2x-3)+(x+1)^2(3-2x)$.
  	\begin{enumerate}
  		\item Écris chacune de ces expressions sous la forme d'un produit de facteurs de premier degré en $x$
  		\item Calcule $A(0)$ et $B(\sqrt{5})$ 
  	\end{enumerate}
  	
  	\dures
  	 \ret{$\bullet$ Mettre un polynôme sous forme de produit de facteurs du premier degré, c'est transformer une somme algébrique de polynôme en un produit de polynômes de degré 1
  		
  		
  		$\bullet$ Factoriser un polynôme, c'est transformer une somme algébrique de polynômes en un produit de polynômes }
 
  	\setcounter{q}{0}
  	\sequence{EQUATION DE DROITE}
  	\act{ Équations du $1\up{er}$ degré dans $\mathbb{R}$}	
  	
  	Le comptable souhaite fabriquer l'herbicide. Il veut réserver sur chaque rive un nombre $x$ d'hectares pour les premiers essais dès que l'herbicide sera fabriqué. En effet, il sait qu'en fonction du nombre $x$ d'hectares, il faut $(2x-13)(x+3)$ litres d'herbicide pour désinfecter la première rive et $\left(x^2-9\right)$ litres d'herbicide pour désinfecter la seconde rive. Comlan se demande combien d'hectare au total faut-il réserver pour les premiers essais si chaque rive reçoit la même quantité d'herbicide.
  	
  	
  	\conAct{Équations du type $ax+b=0$ et $(ax+b)(cx+d)=0$ avec $a\neq 0$ et $c\neq 0$ }
  	\begin{enumerate}
  		\item  Détermine le nombre réel solution de l'équation $2x-13=0$
  		\item Chaque rive reçoit la même quantité d'herbicide signifie que $P(x)=0$ avec $ P(x)=(2x-13)(x+3)-\left(x^2-9\right)$.
  		\begin{enumerate}
  			\item Mets sous forme du produit de facteurs du premier degré en $x$ l'équation $P(x)$.
  			\item Sachant que pour tout réels $a$ et $b$ on a $ab=0$ équivaut à $a=0$ ou $b=0$. Détermine les nombres réels solutions de l'équation $P(x)=0$.
  		\end{enumerate}
  		\item Quel est le nombre total d'hectares faut-il réserver dans ces conditions ?
  	\end{enumerate} 
  	
  	\dures
  	 
  	\defs{On appelle équation du premier degré dans $\mathbb{R}$, toute équation de la forme $ax+b=0$ avec $a\in \mathbb{R^*}$ et $b\in \mathbb{R}$ ou qui peut se mettre sous cette forme.}
  	\propr{ $\bullet$ Le produit de deux nombres réels est nuls signifie que l'un au moins de ces deux nombres est nul. Autrement dit, pour tous nombres réels $a$ et $b$, $ab=0$ équivaut à $a=0$ ou $b=0$.
  		
  		$\bullet$ Une équation de type $(ax+b)(cx+d)=0$ avec $a\neq 0$ et $c\neq 0$ se résout suivant la propriété  $(ax+b)=0$ ou $(cx+d)=0$ qui équivaut à $x=-\dfrac{b}{a}$ ou $x=-\dfrac{d}{c}$.
  	}
  	\nbVo{Remarque}{ L'équation $(ax+b)(cx+d)=0$ avec $a\neq 0$ et $c\neq 0$ n'est pas du premier degré en $x$.}
  	\conAct{Approfondissement}
  	
  	A.	On considère les équations suivantes: $(E_{1}): 2x-3=0 $; $(E_{2}): 2t-3=0 $; 
  	$(E_{3}): 3x^2+8=0 $ ; $(E_{4}): 0x+\sqrt{2}=0 $ ; $(E_{5}):\left|x+\sqrt{3} \right |=\sqrt{2} $ et $(E_{6}): -3x+17= -5x+7 $.
  	\begin{enumerate}
  		\item Précise les équations qui sont du premiers degré dans $\mathbb{R}$.
  		\item 
  		\begin{enumerate}
  			\item Résous $(E_{1})$ dans $\mathbb{R}$
  			\item Déduis-en les solutions de $(E_{2})$ dans  $\mathbb{R}$	
  		\end{enumerate} 
  		\item Résous dans $\mathbb{Z}$ et dans  $\mathbb{R}$ chacune des équations $(E_{5})$ et $(E_{6})$.
  	\end{enumerate}
  	B. On donne $ P(x)= 4x^2 - (2x-3)(x-1) + (6-2x) -9 $
  	\begin{enumerate}
  		\item Mets P(x) sous la forme de facteur du premier degré en $ x $.
  		\item Résous l'équation $ P(x)=0 $ dans  $ \mathbb{R} $
  		\item Résous l'équation $ x^2=2 $ et $ x^2=-3 $
  	\end{enumerate}
  	\dures
  	 
  	\ret{ \ding{118} Si $b$ est un réel positif, $\left|x \right |=b$  équivaut à $x=b$ ou $x=-b$
  		
  		\ding{118} Pour tout nombres réels $a$ et $b$:
  		
  		$\bullet$ Si $b<0$, alors l'équation $\left|x-a \right |=b$ n'admet pas de solution réelle;
  		
  		$\bullet$  Si $b>0$, alors l'équation $\left|x-a \right |=b$ admet  deux solutions réelles qui sont $b-a$ et $b+a$;
  		
  		$\bullet$  Si $b=0$, alors l'équation $\left|x-a \right |=b$ admet une solution réelle qui est $a$.}
  	\nbVo{Remarque}{ L'équation $\left|x-a \right |=b$ n'est pas du premier degré en $x$.}
  	
  	\act{ Équations du $1\up{er}$ degré dans $\mathbb{R}\times \mathbb{R} $}
  	
  	Le contrôle d’un litre du produit $A$ coûte $90 F$ et celui d’un bidon du produit $B$ coûte $500 F$. Le comptable désire contrôler $x$ litres du produit $A$ et $y$ bidons du produit $B$; le montant total à payer est de $86.050 F$.
  	
  	\conAct{Découverte}
  	\begin{enumerate}
  		\item Écris une équation $(E)$ traduisant cette situation.
  		
  		\bcinfo{\textbf{Information}: \textbf{Une telle équation est appelée équation du $1\up{er}$ degré dans $\mathbb{R}\times \mathbb{R} $}}
  		\item  Vérifie si le comptable peut contrôler 45 litres du produit $A$
  		et 164 bidons du produit $B$.	
  		
  		\bcinfo{\textbf{Information}: \textbf{On dit que le couple (45 ; 164) est une solution de l’équation $(E)$}}
  		\item 
  		\begin{enumerate}
  			\item  Justifie que l'équation $ (E) $ est équivalent à $ (E'): 9x + 50y - 8605=0 $
  			\item Parmi les couples $(x;y)$ des réels suivants $(50;72)$; $(945;2)$ et $(30;45)$ indique ceux qui sont solutions de $(E')$.
  			\item  Vérifie que quand un couple est solution de l'équation $(E')$ il est aussi solution de $(E)$
  		\end{enumerate}
  		\item L'équation $(E)$ admet-elle un unique couple solution ?
  		\item  On considère l'équation  $(E_{1}): 2u-3v=3$
  		\begin{enumerate}
  			\item  Détermine $u$ pour $v=3$
  			\item  Détermine $v$ pour $u=0$
  		\end{enumerate}
  	\end{enumerate}
  	\dures
  	 
  	\defs{On appelle équation du premier degré dans $\mathbb{R}\times \mathbb{R}$, toute équation de la forme $ax + by + c = 0$ où $a$, $b$ et $c$ sont des constantes réelles avec $a$ et $b$ non tous nuls. Le couple $(x, y)$ est l’inconnue de l’équation}
  	
  	\nbVo{Remarques}{\ding{118} $a$ et $b$ non tous nuls signifie $(a ;b)\neq (0;0)$.
  		
  		\ding{43} $(a ;b)\neq (0;0)$ signifie :
  		
  		$\bullet$ $a\neq 0$ et $b=0$
  		
  		$\bullet$ ou $a\neq 0$ et $b\neq 0$
  		
  		$\bullet$ ou $a= 0$ et $b\neq 0$
  		
  		\ding{118} Un couple $(x_{0};y_{0})$ est une solution $ax + by + c =0$ veut dire que $ax_{0} + by_{0} + c = 0$. 	
  
   \ding{118} Une équation du $1\up{er}$ degré dans $\mathbb{R}\times \mathbb{R}$ admet plusieurs solutions. Ses solutions sont des couples de nombres réels. 
  		
  		\ding{118} Lorsqu'on multiplie ou on divise membre à membre une équation par un nombre réel non nul, on obtient une nouvelle équation dite équivalente à l'équation initiale.	
  		
  		\ding{118} Deux équations équivalentes ont le même ensemble de solution.
  	}	
  	\conAct{Approfondissement}
  	
  	On a considère les équations suivantes  $(E_{1}): -2x-7y+14=0 $; $(E_{2}): 3a+5b-5=0 $; 
  	$(E_{3}): 3x^2+8=0 $ ; $(E_{4}): x-7y^2+\sqrt{2}=0 $ ; $(E_{5}):-2x+7=0 $ et $(E_{6}): -7y+4=0 $.
  	\begin{enumerate}
  		\item Précise les équations qui sont du premiers degré dans $\mathbb{R}\times \mathbb{R}$.
  		\item Détermine deux couples solutions de l'équation $(E_{1})$.	
  		\item Vérifie que le couple $\left(\dfrac{7}{2};1\right)$ est un couple solution des équations  $(E_{1})$ et $(E_{5})$.
  	\end{enumerate}
  	\dures
  	 
  	\act{ Équations cartésiennes d’une droite- équation réduite d'une droite}
  	
  	Le plan est muni d'un repère orthonormé $\left(O; I;J\right)$. On considère l'équation $(E): x+2y-2= 0 $. $A$, $B$, $C$ et $D$ désignent quatre points du plan dont les cordonnées sont solutions de l'équation $(E)$.
  	
  	\conAct{Équations cartésiennes d’une droite}
  	\begin{enumerate}
  		\item Complète le tableau suivant
  		
  		\begin{tabular}{|l|c|c|c|c|}
  			\cline{2-5}
  			\multicolumn{1}{c|}{} & $A$ & $B$&$C$&$D$ \\
  			\hline
  			$x$ &  & 2&$-1$&\\ 
  			\hline
  			$y$ & 1 & &&$-1$\\
  			\hline 
  		\end{tabular}	
  		
  		\item Place les points $A$, $B$, $C$ et $D$ dans le repère orthonormé $\left(O; I;J\right)$. 
  		\item Construis la droite $(AB)$ puis précise les positions des points $C$ et $D$  par rapport à la droite $(AB)$.
  		\item Place dans ce repère le point $N(2;2)$. $N$ appartient-il à $(AB)$ ?
  		\item  Détermine graphiquement les coordonnées du point $K$ de la droite $(AB)$ d'abscisse $1$.
  	\end{enumerate}
  	
  	\dures
  	 
  	\nbVo{Remarques}{Les représentations des solutions de l’équation sont des points alignés }
   
  	\propr{ \ding{118} Toutes les solutions d’une équation du $1\up{er}$ degré dans 
  		$\mathbb{R}\times \mathbb{R}$, sont représentées sur une même droite ;
  		
  		\ding{118} Tout point d’une droite représente (par ses coordonnées) l’une des solutions 
  		d’une équation du 1er degré dans $\mathbb{R}\times \mathbb{R}$.
  		
  		\ding{118} Dans le plan muni d'un repère, toute droite a une équation de la forme $ax+by+c=0$ $ avec (a;b)\neq (0;0)$. De même, toute équation de la forme $ax+by+c=0$ détermine une équation d'une droite dans un plan muni d'un repère.
  		
  		\ding{118} L'équation  $ax+by+c=0$ qui définit la droite $\mathcal{(D)}$ est appelée équation cartésienne de cette droite.}
  	
  	\ret{Pour tracer dans un plan muni d'un repère une droite dont on connait une de ses équations, on détermine d'abord les coordonnées de deux points distincts de cette droite; puis on place ces deux points dans le repère et enfin, on trace droite qui passe par ces deux points.}
  	\conAct{Approfondissement}
  	Le plan est muni d'un repère orthonormé $\left(O; I;J\right)$ et on considère les points $ A(2;3) $, $ B(1;-2) $ et $ C(1;3) $
  	\begin{enumerate}
  		\item Justifie qu'une équation de la droite $ (AB) $ est $ -5x+y+7=0 $
  		\item Construire la droite $ (AB) $, $ (AC) $ et $ (BC) $ dans le plan   muni du repère orthonormé $\left(O; I;J\right)$
  		\item Construire la droite $ (\Delta) $ d'équation $ -2x+3y-4=0 $ dans le même repère
  	\end{enumerate}
  	\dure s
  	 
  	\conAct{Équations réduite d’une droite}
  	
  	Écris l'équation de la droite  $ (\Delta) $ de la consigne 3.2 sous la forme $y=ax+b$ où $a$ et $b$ sont des réels dont on déterminera.
  	\bcinfo{\textbf{Information}: \textbf{Une telle équation est appelée équation réduite de la droite $ (\Delta) $}}\\
  	\dures
  	 
  	\defs{ Toute droite $\mathcal{(D)}$ non parallèle à l'axe des ordonnées a une équation de la forme $y=ax+b$. \\  \ding{43} L'équation  $y=ax+b$ est alors appelée équation réduite de la droite $\mathcal{(D)}$ ;  
  		
  		\ding{43} Le réel $a$ est le coefficient directeur de cette droite; 
  		
  		\ding{43} Le réel $b$ est son ordonnée à l'origine.}
  	\conAct{Coefficient directeur d'une droite non parallèle à l'axe des ordonnées}
  	
  	Soit $y=ax+b$ l'équation réduite d'une droite $(\Delta)$ et $A(x_{A};y_{A})$ et $B(x_{B};y_{B})$ deux points distincts de cette droite tels que $x_{A}\neq x_{B}$. 
  	
  	\begin{enumerate}
  		\item 
  		\begin{enumerate}
  			\item Exprime $y_{A}$ en fonction de $x_{A}$ 
  			\item Exprime $y_{B}$ en fonction de $x_{B}$ 
  		\end{enumerate}
  		\item Calcule  $y_{B}- y_{A}$ puis justifie que $a=\dfrac{y_{B}- y_{A}}{x_{B}- x_{A}}$ avec $x_{A}\neq x_{B}$.
  		\item Exprime $b$ en fonction de $a$ et des coordonnées du point $B$.
  	\end{enumerate}
  	\dures 
  	 
  	\propr{ Soit $A(x_{A};y_{A})$ et $B(x_{B};y_{B})$ deux points distincts d'une droite $\mathcal{D}$ non parallèle à l'axe des ordonnées. On a 
  		
  		\ding{43} Le coefficient directeur $a$ de la droite $\mathcal{(D)}$ est: $a=\dfrac{y_{B}- y_{A}}{x_{B}- x_{A}}$;
  		
  		\ding{43} L'ordonnée à l'origine $b$   de la droite $\mathcal{(D)}$ est: $b=y_{B}-ax_{B}$ ou $b=y_{A}-ax_{A}$.}
  	
  	\nbVo{Remarque}{ Lorsque deux points $A$ et $B$ ont la même abscisse $(x_{A}=x_{B}=k)$ alors l'équation de la droite $(AB)$ est $x=k$.
  		
  		Lorsque $ a $ est le coefficient directeur d'une droite $ (D) $ alors une équation de la droite $ (D) $  passant par $ A(x_A;y_A) $ est $ y-y_A=a(x-x_A) $}
  	\conAct{Le coefficient directeur d'une droite et un point appartenant à cette droite}
  	On considère la droite $ (\Delta)  $ dont le coefficient directeur est $ a=2 $ et le point $ A(-1;2) $ appartenant à $ (\Delta) $
  	\begin{enumerate}
  		\item En utilisant le coefficient directeur de la droite $ (\Delta ) $ donne les coordonnées du point $ H $ de $ (\Delta) $.
  		\item Construire la droite $   (\Delta)   $
  		\item Détermine l'équation cartésienne de la droite $ (\Delta) $	
  	\end{enumerate}. 
  	\dures
  	   	\ret{Pour construire dans le plan muni d'un repère orthonormé $ (O; I ;J) $, une droite non parallèle à l'axe des ordonnées pasant par le point $ A (x_A;y_A)$  et connaissant son coefficient directeur $ a $, on choisit un point $ B(x_B;y_B) $ à partir de la relation $ \dfrac{y_B-y_A}{x_B-x_A}=a$ puis on trace la droite $ (AB) $.
  		
  		La droite recherchée est la droite $ (AB) $ 
  		
  		 Une équation cartésienne d'une droite connaissant son coefficient directeur $ a $ et un de ces points $ A(x_A;y_A) $  est : $ ax-y+y_A-ax_A=0 $
  		
  		Cette équation est obtenue à partir de la relation $ \dfrac{y -y_A}{x -x_A}=a$ où $ (x;y) $ est le couple de coordonnée d'un point quelconque de la droite.}
  	\act{Droites parallèles- droites perpendiculaires }
  	\conAct{Droites parallèles}
  	\hspace{0.5cm} Dans le plan muni d'un repère  $\left(O; I;J\right)$, on considère les droites: $(\mathcal{D}_{1}): y=2x+ 1 $ ; $(\mathcal{D}_{2}): y=2x+ 3 $ et la droite  $ (\Delta) $ parallèle à $  (\mathcal{D}_{1}) $  et passant par $ A(1;0) $    
  	\begin{enumerate}
  		\item 	\begin{enumerate} \item  Trace les droites $ (\mathcal{D}_{1}) $ et $ (\Delta) $.
  			\item Détermine graphiquement les coordonnées du point $ B \in (\Delta)  $ dont l'abscisse est 0.
  			\item Calcule le coefficient directeur de la droite $(\Delta) $, compare le avec celui de la droite $ (\mathcal{D}_{1}) $ puis conclure 
  		\end{enumerate}
  		\item  Compare le coefficient directeur des droites $ (\mathcal{D}_{1}) $ et $ (\mathcal{D}_{2}) $
  		\item Trace la droite $ (\mathcal{D}_{2}) $ dans le même repère.
  		\item A l'aide de la règle et de l'équerre donne la position relative des droites $ (\mathcal{D}_{1}) $ et $ (\mathcal{D}_{2}) $ puis conclure 
  	\end{enumerate}
  	\dures
  	 \propr{\hspace{1cm} Le plan est muni d’un repère orthonormé.
  		
  		\ding{118} Si deux droites, non parallèles à l’axe des ordonnées, sont 
  		parallèles, alors elles ont même coefficient directeur. 
  		
  		$(\mathcal{D}) \parallel (\mathcal{D'})$  alors $a=a'$ avec $a$ et $a'$ les coefficients directeurs respectifs de $(\mathcal{D})$ et $(\mathcal{D'}).$
  		
  		\ding{118}  Si deux droites, non parallèles à l’axe des ordonnées, ont 
  		même coefficient directeur alors elles sont parallèles
  		
  		$a=a'$ alors $(\mathcal{D}) \parallel (\mathcal{D'})$  avec $a$ et $a'$ les coefficients directeurs respectifs de $(\mathcal{D})$ et $(\mathcal{D'}).$
  	}
  	\conAct{Application}
  	Dans le plan muni d'un repère orthonormé $ (F; G;H) $ on considère la droite $ (\mathcal{D}) $ d'équation $ y=3x-7 $.
  	\\ Détermine une équation cartésienne de la droite $ (\Delta) $ passant par $E(0;-2)  $ et parallèle à $ (\mathcal{D}) $.\\
  	\dures
  	 
  	\conAct{Droites perpendiculaires}
  	Dans le plan muni d'un repère orthonormé $ (O;I;J) $ on considère les droites $ (\Delta): y=-2x+1 $; $   \mathcal{(D)}: y=\dfrac{1}{2}x-3$ et la droite $ (\Delta') $ perpendiculaire à $ (\Delta) $ passant par $ G(0;3) $
  	\begin{enumerate}
  		\item 	\begin{enumerate} \item  Trace les droites $ ( \Delta) $ et $ (\Delta') $.
  			\item Détermine graphiquement les coordonnées du point $ H \in (\Delta')  $ dont l'ordonnées est 2.
  			\item Calcule le coefficient directeur $ a' $ de la droite $(\Delta') $, le produit $ a\times a' $  avec $ a $ le coefficient directeur de la droite  $ (\Delta) $ puis conclure 
  		\end{enumerate}
  		\item  Calcule le produit $ a\times a'' $  avec $ a'' $ le coefficient directeur des droites $ (\mathcal{D}) $
  		\item Trace la droite $ (\mathcal{D}) $ dans le même repère.
  		\item A l'aide de la règle et de l'équerre donne la position relative des droites $ (\mathcal{D}) $ et $ (\Delta) $ puis conclure 
  	\end{enumerate}
  	\dures
  	 
  	 \propr{ Le plan est muni d’un repère orthonormé.
  		
  		\ding{118} Si deux droites non parallèles à l’axe des ordonnées sont 
  		perpendiculaires alors le produit de leurs coefficients 
  		directeurs est égal à $-1$
  		
  		Si $(\mathcal{D}) \perp (\mathcal{D'})$  alors $a\times a'=-1$ avec $a$ et $a'$ les coefficients directeurs respectifs de $(\mathcal{D})$ et $(\mathcal{D'})$.
  		
  		\ding{118} Si le produit des coefficients directeurs de deux droites 
  		non parallèles à l’axe des ordonnées est égal à $-1$, alors elles sont perpendiculaires. 
  		
  		Si	$a\times a'=-1$ alors $(\mathcal{D}) \perp (\mathcal{D'})$  avec $a$ et $a'$ les coefficients directeurs respectifs de $(\mathcal{D})$ et $(\mathcal{D'})$.
  	}
  	
  	\conAct{Approfondissement} 	
  	
  	Le plan est muni d'un repère orthonormé $\left(O; I;J\right)$. 
  	\begin{enumerate}
  		\item On considère les points $A(1;-3)$ et  $B(2;-2)$. 
  		\begin{enumerate} 
  			\item Justifie que la droite $(AB)$ n'est pas parallèle à l'axe des ordonnées.
  			\item Détermine une équation réduite de la droite $(AB)$ puis déduis-en son équation cartésienne.
  			\item  Construis la droite $(AB)$ dans ce repère.
  		\end{enumerate}
  		\item   On considère la droite $(\mathcal{D}_{1}): y=\dfrac{4}{3}x+\sqrt{3}$. Détermine l'équation cartésienne de la droite $({\Delta}_{1})$ parallèle à  $(\mathcal{D}_{1})$ passant par $C(1;2)$.
  		\item On considère la droite $(\mathcal{D}_{2}): 4y=\sqrt{2}x-1$. Détermine l'équation réduite de la droite $({\Delta}_{2})$ perpendiculaire à  $(\mathcal{D}_{2})$ passant par $D(\sqrt{2};1)$.
  		\item Construis dans le même repère la droite $(\mathcal{D}_{3})$ parallèle à l'axe des ordonnées passant par $B(2;-2)$.
  	\end{enumerate}
  	
  	\dures
  	 
  	\begin{center}
  		\Ovalbox{\textbf{ Exercice de maison}}
  	\end{center}
  	\begin{center}
  		\Ovalbox{\textbf{ Exercice 1 }}
  	\end{center}
  	Le plan est muni d'un  repère orthonormé $\left(O; I;J\right)$.
  	
  	$(\mathcal{D})$ est la droite passant par le point  $A(1;4)$ et de coefficient directeur $a=-2$. On souhaite construire la droite $(\mathcal{D})$, sans déterminer un second point ni une équation de $(\mathcal{D})$. Pour cela:
  	\begin{enumerate}
  		\item Place le point $A$ dans le repère.
  		\item Construis la droite $(\mathcal{\Delta})$ d'équation $y=-2x$ dans ce repère.
  		\item Trace la droite  $(\mathcal{D})$ passant par $A$ et parallèle à $(\mathcal{\Delta})$.
  	\end{enumerate}
  	\begin{center}
  		\Ovalbox{\textbf{ Exercice 2 }}
  	\end{center}
  	Détermine les valeurs de $x$ dans chacun des cas d'équations suivantes.
  	
  	$(E_{1}): \left(x+\sqrt{2}\right)^2=\left(x+1\right)^2 $; $(E_{2}): 4x^2-1=0$ ; $(E_{3}): \left(x+3\right)^2+(3x+9)(2x-2)=0$ ; $(E_{4}): \left|x\right |=\sqrt{5} $ et  $(E_{5}): \left|2x+5\right |=0 $.
  	
  	\begin{center}
  		\Ovalbox{\textbf{ Exercice 3 }}
  	\end{center}
  	Le plan est muni d'un repère orthonormé $\left(O; U;V\right)$, l'unité graphique est $1 cm$. On donne les droites $({\Delta}_{1})$; $({\Delta}_{2})$; $({\Delta}_{3})$ et $({\Delta}_{4})$ d'équations respectives: $2x+y-5=0$; $x=5$; $y=\dfrac{1}{2}x-1$; et $y=3$; puis les points $A(2;-2)$ et  $B(0;1)$. 
  	\begin{enumerate}
  		\item Démontre que les droites $({\Delta}_{1})$ et $({\Delta}_{3})$ sont perpendiculaires. 
  		\item Démontre que les droites $({\Delta}_{2})$ et $({\Delta}_{4})$ sont perpendiculaires.
  		\item Détermine l'équation réduite de la droite $ (D) $ passant par $A$ et parallèle à $({\Delta}_{1})$ 
  		\item Détermine l'équation cartésienne de la droite $ (D') $ passant par $B$ et perpendiculaire à $({\Delta}_{3})$.
  		\item Montre que $ (D)\perp (D') $
  		\item  Représente les droites $({\Delta}_{1})$; $({\Delta}_{2})$; $({\Delta}_{3})$ et $({\Delta}_{4})$.
  	\end{enumerate}
  	\begin{center}
  		\Ovalbox{\textbf{ Exercice 4}}
  	\end{center}
  	Le plan est muni d'un repère orthonormé et on considère les
  	$A(0;1)$, $ C(-2;2), $  $B(1;-1)$, $ D(-6;0) $, $ E(0;-2) $ , $ F(4;0) $  et $ G(0;3) $.
  	\begin{enumerate}
  		\item Place ces points dans un repère orthonormé.
  		\item Détermine l'équation cartésienne de la $ (AB) $.
  		\item Justifie que les points $ C $, $ D $ et $ G $ sont alignés.
  		\item Justifie que les droites $ (CD) $ et $ (EF) $ sont parallèles
  		\item Démontre que les droites $ (AB) $ et $ (CD) $ sont perpendiculaires 
  	\end{enumerate}
  	\begin{center}
  		\Ovalbox{\textbf{ Exercice 5}}
  	\end{center}
  	\textbf{CIAM 3\up{è}}: Exercices 1, 2 pages 167 et Exercices 12, 13, 14 pages 168  
  \sequence{EQUATIONS ET INEQUATIONS }
  \setcounter{q}{0}
  	\act{ Inéquations du $1\up{er}$ degré dans $\mathbb{R}$}
  	
  	Avant la fabrication de l'herbicide, le comptable a réservé sur chaque rive plus de $20 ha$ pour les premiers essais et l'herbicide doit être utilisé à raison de 4 litres par hectare. Comlan voudrait une estimation de la quantité d'herbicide à utiliser sur chaque rive pour l'épandage.
  	
  	\conAct{Notion d'inéquations du $1\up{er}$ degré dans $\mathbb{R}$ } 	
  	\begin{enumerate}
  		\item On désigne par $x$ la quantité en litres d'herbicide à utiliser sur une superficie de plus de $20 ha$.
  		
  		Donne une inégalité qui traduit cette phrase  
  		
  		\bcinfo{\textbf{Information}: \textbf{Cette écriture est appelée inéquation du premier degré à une inconnue dans $\mathbb{R}$ }}.
  		
  		\item Écris sous forme d'intervalle l'ensemble des nombres réels $x$ vérifiant cette inégalité.
  		
  		\bcinfo{\textbf{Information}: \textbf{Cet ensemble est appelé ensemble des solutions  dans $\mathbb{R}$ de l'inéquation }}.
  	\end{enumerate}
  	\dures
  	 
  	\defs{On appelle inéquation du premier degré dans $\mathbb{R}$, toute inégalité ayant l'une  des  formes suivantes : $ax+b<0$; $ax+b\leqslant0$; $ax+b>0$ et $ax+b\geqslant0$ où $a\in \mathbb{R}\up{*}$ et $b\in \mathbb{R}$; $x$ étant l'inconnue.
  	}
  	\nbVo{Remarques}{\ding{118} On dit qu'un nombre réel est une solution d'une inéquation si on obtient une inégalité qui est vraie lorsqu'on remplace l'inconnue par ce nombre dans l'inéquation.
  		
  		\ding{118} Lorsqu'une inéquation du premier degré dans $\mathbb{R}$ est résolue, l'ensemble des solutions est généralement un intervalle de $\mathbb{R}$.
  		
  		\ding{118} L'ensemble des solutions d'une inéquation du premier degré dans $\mathbb{R}$ sera représenté par un intervalle sur une droite graduée.}
  	
  	\conAct{Approfondissement} 
  	
  	Résous chacune des inéquations et représente graphiquement leur ensemble solutions.
  	$(I_{1}): -2x+3\leqslant-3x+1$ ; $(I_{2}): 2x+3\geqslant-2x-1$; $(I_{3}): 8-(2x-7)< 4x+3$ et $(I_{4}): x\sqrt{5}+5<-2x+3$
  	\dures
  	 
  	\act{Système d'inéquations du $1\up{er}$ degré dans $\mathbb{R}$} 
  	
  	En réalité pour les premiers essais, le comptable a réservé plus de $20 ha$ sur la première rive et moins de $28 ha$ sur la deuxième ri	ve et l'herbicide doit être utilisé à raison de 4 litres par hectare. On désigne par $x$ la quantité en litres d'herbicide à utiliser sur la première et sur la deuxième rive. 
  	
  	Soit $(I_{1})$ et $(I_{2})$ les inéquations portant sur $x$ respectivement pour la première et la deuxième rive.
  	
  	\conAct{Notion de système d'inéquations du $1\up{er}$ degré dans $\mathbb{R}$}
  	
  	\begin{enumerate}
  		\item Justifie que les inéquations $(I_{1})$ et $(I_{2})$ sont respectivement $4x>20$ et $4x<28$.
  		\bcinfo{\textbf{Information}: \textbf{Des inéquations $(I_{1})$ et $(I_{2})$, on obtient le système suivant}} $(S)$: $\systeme{4x > 20,4x < 28}$
  		
  		\item Résous dans $\mathbb{R}$ les inéquations $(I_{1})$ et $(I_{2})$  et représente graphiquement leur ensemble solutions sur une même droite graduée.
  		\item Déduis-en les valeurs de $x$ qui vérifie à la fois $(I_{1})$ et $(I_{2})$.
  		
  		\bcinfo{\textbf{Information}: \textbf{ L'ensemble des valeurs de $x$ qui vérifie à la fois $(I_{1})$ et $(I_{2})$ est l'ensemble solutions du système $(S)$ }}
  	\end{enumerate}
  	
  	\dures
  	  
  	\defs{Un système d'inéquation du premier degré dans $\mathbb{R}$, est un système qui contient plusieurs inéquations dans du premier dans $\mathbb{R}$. }
  	\nbVo{Remarque}{ L'ensemble des solutions de ce système est obtenu en faisant l'intersection des ensembles des solutions de toute les inéquations qui composent ce système.
  	}
  	
  	\conAct{Approfondissement}
  	Résous chacun des systèmes inéquations suivantes.
  	
  	$(S_{1})$: $\systeme{4x-3 \geqslant 2x+5,3x+5 < 6x-1}$ et  $(S_{2})$: $\systeme{-2x +6\geqslant 0, 5x+5>3x+1}$
  	\dures
  	 
  	\act{Système d'équations du $1\up{er}$ degré dans $\mathbb{R}\times \mathbb{R}$} 
  	
  	Pour l'épandage de l'herbicide, le comptable doit utiliser deux drones de pulvérisation, $d_{1}$ et $d_{2}$. Les deux drones sont soumis à des contraintes dans leurs opérations.
  	
  	\textbf{$1\up{ère}$  contrainte}: La quantité d'herbicide épandue par le drone $d_{2}$ est le double de celle épandue par le drone $d_{1}$, diminuée de 3 hectolitres. 
  	
  	\textbf{$2\up{ème}$  contrainte}: La quantité totale d'herbicide à utiliser est de 6 hectolitres.
  	
  	On désigne par $x$ le nombre d'hectolitres d'herbicide épandus par le drone $d_{1}$ et par $y$ celui épandu par le drone $d_{2}$.
  	
  	\conAct{Découverte}
  	
  	\begin{enumerate}
  		\item Traduis par une équation $(E_{1}) $ la $1\up{ère}$  contrainte.
  		\item Traduis par une équation $(E_{2}) $ la $2\up{ème}$  contrainte.
  		\item En utilisant les équations $(E_{1})$ et $(E_{2})$, écris un système $(S)$ de deux  équations du premier degré $\mathbb{R}\times \mathbb{R}$.
  	\end{enumerate}
  	
  	\dures
  	 
  	\defs{Un système d'équations du premier degré dans $\mathbb{R}\times \mathbb{R}$, est un système qui contient plusieurs équations du premier degré dans $\mathbb{R}\times \mathbb{R}$.}
  	
  	\nbVo{Remarque}{ Résoudre un tel système c'est chercher l'ensemble des couples de nombres réels vérifiant à la fois chacune des équations du système. Il existe plusieurs méthodes.
  	}
  	
  	\conAct{Résolution par la méthode graphique}
  	
  	On considère le système d’équations suivant: $(S)$: $\systeme{2x-y=3, x+y=6}$ 
  	\begin{enumerate}
  		\item Trace dans un repère orthonormé $\left(O; I;J\right)$ les droites $(\mathcal{D}_{1}): 2x-y=3$  et $(\mathcal{D}_{2}): x+y=6$.
  		\item Détermine graphiquement les coordonnées du point d'intersection $A$ des $(\mathcal{D}_{1})$ et $(\mathcal{D}_{2})$  
  		\item Déduis-en l'ensemble $S_{\mathbb{R}\times \mathbb{R}}$  des solutions de $(S)$.
  		\item Précise le nombre $x$ d'hectolitres d'herbicide épandus par le drone $d_{1}$ et $y$ celui épandu par le drone $d_{2}$.
  	\end{enumerate}
  	\dures
  	
  	 
  	\ret{La méthode de résolution graphique d'un système de deux équations du premier degré dans $\mathbb{R}\times \mathbb{R}$ consiste à tracer les droites qui ont pour équations, les équations du système et à détermine graphiquement les coordonnées du point qui est leur intersection. Ce couple de coordonnées est le couple solution  de ce système.}
  	
  	\conAct{Résolution par la méthode de combinaison (ou addition)}
  	
  	
  	On considère le système d’équations suivant: $(S)$: $(S)$: $\systeme{2x-y=3 \hspace{0.3cm} (1), x+y=6 \hspace{0.3cm} (2)}$
  	\begin{enumerate}
  		\item 
  		\begin{enumerate}
  			\item Multiplie l'équation (2) par l'opposé du coefficient de l'inconnue $x$ dans l'équation (1).
  			\item  Additionne membre à membre la nouvelle équation obtenue et l'équation (1) pour éliminer l'inconnue $x$ puis déduis-en la valeur de $y$. 
  		\end{enumerate}
  		\item Reprends le système $(S)$, additionne membre à membre les équations (1) et (2)   pour éliminer l'inconnue $y$ puis déduis-en la valeur de $x$. 
  		\item  Déduis-en l'ensemble $S_{\mathbb{R}\times \mathbb{R}}$  des solutions de $(S)$.
  	\end{enumerate}
  	\dures
  	 
  	\ret{Pour résoudre un système de deux équations à deux inconnues par la méthode de combinaison (ou la méthode d'addition), on procède de la façon suivante:
  		
  		\ding{182} On multiplie chacune des équations (1) et (2) du système par un coefficient convenablement choisi de manière à obtenir un nouveau système dans lequel une des inconnues a des coefficient opposés;
  		
  		\ding{183} On additionne membre à membre les deux équations afin d'éliminer l'une  des inconnues;
  		
  		\ding{184} De l'équation du premier degré obtenu, on tire l'autre inconnue.
  		
  		\ding{185} On procède de la même manière pour déterminer la deuxième inconnue.
  		
  		\ding{186}  On s'assure après vérification que le couple ainsi obtenu est solution du système.
  	}
  	
  	\conAct{Résolution par la méthode de substitution}
  	
  	On considère le système d’équations suivant: $(S)$: $\systeme{2x-y=3 \hspace{0.3cm} (1), x+y=6 \hspace{0.3cm} (2)}$
  	\begin{enumerate} 
  		\item De l'équation (1), exprime $x$ en fonction de $y$, note (3) la relation ainsi trouvée.
  		\item 
  		\begin{enumerate}
  			\item  Remplace l'expression de $x$ obtenue dans (2).
  			\item L'opération effectuée en \textbf{2.(a)} te donne une équation de premier degré d'inconnue $y$. Résous cette équation; 
  		\end{enumerate}
  		\item Détermine alors la valeur de $x$ en remplaçant la valeur de $y$ trouvée dans la question \textbf{2.(b)} dans la relation notée (3). 
  		\item  Déduis-en l'ensemble $S_{\mathbb{R}\times \mathbb{R}}$  des solutions de $(S)$.
  	\end{enumerate}
  	\dures
  	
   	
  	\ret{La résolution d'un système de deux équations du premier degré à deux inconnues par la méthode de substitution consiste à:
  		
  		\ding{182} tirer une des inconnues dans l'une des équations;
  		
  		\ding{183} remplacer l'expression trouvée dans l'autre équation;
  		
  		\ding{184} résoudre l'équation du premier degré ainsi obtenue pour déterminer la valeur de l'inconnue;
  		
  		\ding{185} remplacer la valeur de l'inconnue trouvée dans l'expression tirée au départ afin de déduire la valeur de l'autre inconnue (il s'agit de l'inconnue tirée en fonction de l'autre);
  		
  		
  		\ding{186} et enfin présenter le couple solution.
  	}
  	\conAct{Résolution par la méthode de comparaison} 
  	
  	On considère le système d’équations suivant: $(S)$: $\systeme{2x-y=3 \hspace{0.3cm} (1), x+y=6 \hspace{0.3cm} (2)}$
  	\begin{enumerate}
  		\item 
  		\begin{enumerate}
  			\item De l'équation (1), exprime $x$ en fonction de $y$, note (3) la relation ainsi trouvée.
  			\item  De l'équation (2), exprime $x$ en fonction de $y$, note (4) la relation ainsi trouvée. 
  		\end{enumerate}
  		\item Pose l'égalité entre les expressions  (3) et (4)  que tu as trouvée dans \textbf{1.(a)} et \textbf{1.(b)}; déduis-en la valeur de $y$.
  		\item Reprends le procédé des questions  \textbf{1.} et \textbf{2.} pour déterminer la valeur de $x$;
  		\item  Déduis-en l'ensemble $S_{\mathbb{R}\times \mathbb{R}}$  des solutions de $(S)$.
  	\end{enumerate}
  	\dure
  	 
  	\ret{La résolution d'un système de deux équations du premier degré à deux inconnues par la méthode de comparaison revient successivement à:
  		
  		\ding{182} Exprimer l'une des inconnues en fonction de l'autre dans chacune des deux équations;
  		
  		\ding{183} Poser l'égalité des expressions trouvées et en déduire la valeur de l'inconnue présente.
  	}
  	\nbVo{Remarque}{Toutes les méthodes de résolution d'un système dans $\mathbb{R}\times \mathbb{R}$ conduisent au même ensemble de solutions.}
  	
  	\conAct{Problèmes se ramenant à des systèmes de deux 
  		équations du premier degré dans $\mathbb{R}\times \mathbb{R}$} 
  	
  	Au concert. Ce fut un beau succès pour le groupe Mygale. Les 1500 places de la salle étaient occupées. la vente des billets, les uns à $140 F$ et les autres à $60 F$, a rapporté $114 000 F$.
  	
  	Détermine le nombre de spectateurs qui ont payé $60F$ et ceux qui ont payé $140F$.
  	
  	\dures
  	   	\ret{Pour résoudre un problème se ramenant à des systèmes de deux 
  		équations du premier degré dans $\mathbb{R}\times \mathbb{R}$, on procède successivement de la façon suivante:
  		
  		\ding{182} Faire le choix des inconnues (par exemple $x$ et $y$ si l'énoncé n'a pas encore imposé un choix).
  		
  		\ding{183} Traduire l'énoncé en un système d'équations du premier degré dans $\mathbb{R}\times \mathbb{R}$.
  		
  		\ding{184} Résoudre le système obtenu.
  		
  		\ding{185} Conclure.
  		
  	}
  	
  	\act{Inéquations du $1\up{er}$ degré dans $\mathbb{R}\times \mathbb{R}$}
  	
  	Les responsables en charge du projet d'épandage se sont aperçus qu'il faut au plus $6$ hectolitres d'herbicide pour la pulvérisation des rives par les drones $d_{1}$ et $d_{2}$. 
  	De plus, le drone $d_{1}$ doit pulvériser $x$ volume d'herbicide et le drone $d_{2}$ $y$ volume. 
  	
  	Le comptable se demande comment il va pouvoir écrire cette contrainte et laisser cette consigne aux deux pilotes.
  	
  	\conAct{Notion d'inéquations du $1\up{er}$ degré dans $\mathbb{R}\times \mathbb{R}$} 
  	
  	\begin{enumerate}
  		\item Écris l'expression de l'inégalité traduisant les nouvelles contraintes.
  		\item En remplaçant le symbole de l'inégalité par celui de l'égalité, trace la droite $(\mathcal{D}_{2})$ d'équation ainsi obtenue dans un repère orthonormé $\left(O; I;J\right)$
  		\item $(\mathcal{D}_{2})$ sépare le plan en deux partie dont elle est frontière.
  		\begin{enumerate}
  			\item Justifie que le couple $(0;0)$ de coordonnées du point $O$ vérifie l'inégalité trouvée en \textbf{1}.
  			\item Justifie que le couple $(4;4)$ ne vérifie pas  cette inégalité.
  			\item Détermine la partie du plan dont les coordonnées des points sont solutions de cette équation en hachurant la zone des points ne répondant pas à cette même condition. 
  		\end{enumerate} 
  	\end{enumerate}
  	\dures
  	
  	 
  	
  	\defs{Une inéquation du premier degré dans $\mathbb{R}\times \mathbb{R}$ est toute inéquation pouvant se ramener à l'une des  formes: $ax+by+c<0$; $ax+by+c\leqslant0$; $ax+by+c>0$ et $ax+by+c\geqslant0$ avec $a$, $b$ deux nombres réels  tous non nuls et  $c\in \mathbb{R}$.
  		
  		Un couple $(x_{0}; y_{0})$ est solution de d'une de ces équations si ce couple vérifie cette inéquation.
  	}
  	\propr{Le plan est muni d'un repère.
  		
  		La droite $(\mathcal{D})$ d'équation $ax+by+c=0$  partage le plan en trois parties:
  		
  		\ding{182} L'ensemble des points de coordonnées $(x;y)$ tels que $ax+by+c=0$
  		
  		\ding{183} L'ensemble des points de coordonnées $(x;y)$ tels que $ax+by+c>0$
  		
  		\ding{184} L'ensemble des points de coordonnées $(x;y)$ tels que $ax+by+c<0$
  	}
  	\nbVo{ Méthode de résolution}{ Soit à résoudre, graphiquement, une inéquation de l'une des forme ci-dessus:
  		
  		\ding{182} on pose $(\mathcal{D}): ax+by+c=0$ ;
  		\ding{183} on trace la droite $(\mathcal{D})$ dans le plan rapporté à un repère;
  		
  		\ding{184} La droite $(\mathcal{D})$ partage le plan en trois parties dont deux demi-plans de frontière $(\mathcal{D})$;
  		
  		\ding{185} on choisit un point de coordonnées $(x_{0}; y_{0})$; 
  		
  		\ding{185} si $(x_{0}; y_{0})$  est solution de l'inéquation, alors le demi-plan de frontière  $(\mathcal{D})$ contenant les points de coordonnées $(x_{0}; y_{0})$ est la partie du plan qui contient les points dont les coordonnées sont des solutions;
  		
  		\ding{186} on hachure le demi-plan des solution ou celui ne contenant pas les solutions suivants l'indication de la consigne;
  		
  		\ding{187} on présente l'ensemble des solutions de l'inéquation (dans le cas d'une inégalité stricte, les coordonnées des points de $(\mathcal{D})$ ne sont pas solution de l'inéquation). 
  		
  	}
  	\conAct{Système d'inéquations du $1\up{er}$ degré dans $\mathbb{R}\times \mathbb{R}$} 
  	
  	Les responsables en charge du projet d'épandage se sont aperçus également qu'il faut 
  	que la quantité d'herbicide épandue par le drone $d_{2}$ soit supérieure ou égale au double de celle épandue par le drone $d_{1}$, diminuée de 3 hectolitres. Ainsi, on a  $2x-y\leqslant3$. Pour donner une consigne assez complète aux deux pilotes le comptable écrit le système suivant $(S)$: $\systeme{2x-y\leqslant3 \hspace{0.3cm} (1), x+y\leqslant6 \hspace{0.3cm} (2)}$.
  	\begin{enumerate}
  		\item Résous graphiquement l'inéquation (1)  
  		\item Résous graphiquement l'inéquation (2) 
  		\item Déduis-en l'ensemble des solutions du système $(S)$.  
  	\end{enumerate}
  	\dures
  	 
  	\nbVo{Méthode}{Pour résoudre graphiquement un système d'inéquations dans $\mathbb{R}\times \mathbb{R}$, on résout graphiquement chacune des inéquations.
  		
  		L'ensemble des solutions de ce système est l'intersection des ensembles de solutions de chacune des inéquations.}
  	\nbVo{Remarque}{C'est toujours graphiquement qu'on résout les inéquations et systèmes d'inéquations dans $\mathbb{R}\times \mathbb{R}$. }	
  	\begin{center}
  		\Ovalbox{\textbf{ Exercice de maison}}
  	\end{center}
  	\begin{center}
  		\Ovalbox{\textbf{ Exercice 1 }}
  	\end{center}
  	Dans chacun des cas suivants résous  dans $\mathbb{R}\times \mathbb{R}$ le système d'inconnue $(x;y)$ présenté par la méthode indiquée.
  	\begin{enumerate}
  		\item  $(S_{1})$: $\systeme{3x+4y=4 , 2x+y=12 }$ (Méthode graphique)
  		\item  $(S_{2})$: $\systeme{-2x+y=1 , 2x+y=12 }$ (Méthode d'addition)
  		\item  $(S_{3})$: $\systeme{x+y=3 , -x+2y=0 }$ (Méthode de substitution)
  		\item  $(S_{4})$: $\systeme{x+y=-3 , -x+2y=10 }$ (Méthode comparaison)
  	\end{enumerate}
  	\begin{center}
  		\Ovalbox{\textbf{ Exercice 2 }}
  	\end{center}
  	\begin{enumerate}
  		\item Résous dans $\mathbb{R}$ les inéquations suivantes:   $(I_{1}): \dfrac{x+2}{4}\leqslant x-\dfrac{1}{3}$; $(I_{2}): -13x+15>-6x+6$ et  $(I_{3}): 3x-5\leqslant -5x<32-x$.
  		\item Résous dans $\mathbb{R}$ $(S_{1})$: $\systeme{x+1<-3x+5 , 2x+7<5x-2 }$ et 
  		$(S_{2})$: $\systeme{-2x+6\geqslant 0 , 5x+5>3x+1 }$ 
  	\end{enumerate}
  	\begin{center}
  		\Ovalbox{\textbf{ Exercice 3 }}
  	\end{center}
  	Dans cet exercice, les questions sont indépendantes. 
  	\begin{enumerate}
  		\item $60$ livres de deux sortes occupent 210 cm sur une étagère, ces 
  		livres sont rangés les uns contre les autres, leurs épaisseurs est 
  		soit $4 cm$ soit $2 cm$. 
  		
  		Détermine le nombre de livres de chaque sorte.
  		\item  Kodjo est un éleveur de moutons et de volailles. dans sa ferme, on compte 72 têtes et 188 pattes d'animaux. 
  		
  		Détermine le nombre de volailles et le nombre de moutons que possède Kodjo.
  		\item La longueur d'un terrain rectangulaire dépasse le double de sa largeur de $6 m$. le périmètre de ce terrain est $84 m$.
  		
  		Détermine les dimensions de ce terrain.
  		\item   Une droite $(\mathcal{D})$ a pour équation cartésienne $ax+by-3=0$. Cette droite passe par les points $A(2;-1)$	 et $B(7;2)$. Détermine $a$ et $b$.
  	\end{enumerate}
  	\begin{center}
  		\Ovalbox{\textbf{ Exercice 4 }}
  	\end{center}
  	\begin{enumerate}
  		\item Résous l'inéquation   $(E): 3x-y>5$ dans $\mathbb{R}\times \mathbb{R}$
  		\item Résous les systèmes d'inéquations 
  		
  		$(S_{1})$: $\systeme{2x-y\leqslant3, x+y\leqslant6,x>0,y>0}$    
  		
  		$(S_{2})$: $\systeme{-2x+y\leqslant2, x+y>12 }$ dans $\mathbb{R}\times \mathbb{R}$
  	\end{enumerate}
  	
  \sequence{SOMME DE DEUX VECTEURS-PRODUIT D'UN VECTEUR PAR UN NOMBRE REEL}
  	\setcounter{q}{0}
  	\act{ Notion de somme de deux vecteurs et du produit d’un vecteur par un nombre réel}
  	
  	Au cours de l’épandage de l’herbicide, les déplacements des drones $d_{1}$ et $d_{2}$ sont schématisés par la figure suivante.
  	
  	\begin{tikzpicture}[scale=0.7]
  		%\draw[dashed] (-4,-4) grid (5,8);	
  		\draw [line width=0.8pt] (1,1)-- (9,1);
  		\draw [line width=0.8pt] (1,1)-- (3,3);
  		\draw [line width=0.8pt] (3,3)-- (5,1);
  		\draw [line width=0.8pt] (5,1)-- (7,3);
  		\draw [line width=0.8pt] (3,3)-- (7,3);
  		\draw [line width=0.8pt] (7,3)-- (9,1);
  		\draw [line width=0.8pt] (0.7,0)-- (9,0);
  		\draw (2.6,3.8) node[anchor=north west] {$A$};
  		\draw (6.8,3.8) node[anchor=north west] {$B$};
  		\draw (4.7,1) node[anchor=north west] {$C$};
  		\draw (0.6,1) node[anchor=north west] {$D$};
  		\draw (8.7,1) node[anchor=north west] {$E$};
  		\draw (4.8,3.8) node[anchor=north west] {$I$};
  		\draw (6.7,1) node[anchor=north west] {$K$};
  		\draw (2.6,1) node[anchor=north west] {$J$};
  		\draw (0.7,0)  node[ below] {$(\mathcal{D})$}; 
  		\begin{scriptsize}
  			\draw [fill=black] (1,1) circle (1.5pt);
  			\draw [fill=black] (9,1) circle (1.5pt);
  			\draw [fill=black] (3,3) circle (1.5pt);
  			\draw [fill=black] (5,1) circle (1.5pt);
  			\draw [fill=black] (7,3) circle (1.5pt);
  			\draw [fill=black] (4,3) circle (1.5pt);
  			\draw [fill=black] (5,3) circle (1.5pt);
  			\draw [fill=black] (6,3) circle (1.5pt);
  			\draw [fill=black] (2,1) circle (1.5pt);
  			\draw [fill=black] (3,1) circle (1.5pt);
  			\draw [fill=black] (4,1) circle (1.5pt);
  			\draw [fill=black] (6,1) circle (1.5pt);
  			\draw [fill=black] (7,1) circle (1.5pt);
  			\draw [fill=black] (8,1) circle (1.5pt);
  		\end{scriptsize}
  	\end{tikzpicture}
  	$ABCD$ et $ABEC$ sont des parallélogrammes, $(DC) \parallel(\mathcal{D}) $ et $DC=CE$
  	\conAct{Somme de deux vecteurs }
  	
  	\begin{enumerate}
  		\item Reproduis les vecteurs $\overrightarrow{CB}$ et $\overrightarrow{CE}$ puis Construire le vecteur $\overrightarrow{CP}= \overrightarrow{CB}+\overrightarrow{CE}$ 
  		
  		\bcinfo{\textbf{Information}: \textbf{On dit que le vecteur $\overrightarrow{CP}$ est la somme des deux vecteurs $\overrightarrow{CB}$ et $\overrightarrow{CE}$ }}
  		\item Détermine la somme des vecteurs suivants: $\overrightarrow{AD}+\overrightarrow{AB}$; $\overrightarrow{AB}+\overrightarrow{DC}$ et $\overrightarrow{DJ}+\overrightarrow{JC}$
  	\end{enumerate}
  	\dures
  	 
  	\propr{$A$,$B$ et $C$ étant trois points quelconques du plan on a: $\overrightarrow{AB}+\overrightarrow{BC}=\overrightarrow{AC}$ (\textbf{ Cette 
  			égalité est appelée " Relation de Chasles"})}
  	%\nbVo{Remarque}{Pour construire la somme de deux vecteurs $\overrightarrow{AB}$ et $\overrightarrow{AC}$, on construit d'abord un parallélogramme sur les côtés $[AB]$ et $[AC]$, soit par exemple le parallélogramme $CBPE$. La somme des vecteurs $\overrightarrow{CB}$ et $\overrightarrow{CE}$ est le vecteur $\overrightarrow{CP}$.}
  	\conAct{Définition du produit d’un vecteur par un réel }
  	
  	On considère la figure de l'activité 1.
  	\begin{enumerate}
  		\item Complète chacune des écritures suivantes : $\overrightarrow{AB}= \dots \overrightarrow{AI}$; $\overrightarrow{DE}= \dots \overrightarrow{AB}$;  $\overrightarrow{DE}= \dots \overrightarrow{AI}$ et $\overrightarrow{DE}= \dots \overrightarrow{IA}$ 
  		\item Dis ce tu penses des directions, du sens puis des longueurs des 
  		vecteurs $\overrightarrow{DE}$ et  $ \overrightarrow{AB}$  d’une part, et des vecteurs $\overrightarrow{DE}$ et $ \overrightarrow{IA}$ d’autre 
  		part.
  		
  		\bcinfo{\textbf{Information}: \textbf{On dit que le vecteur $\overrightarrow{DE}$ est le produit du vecteur $\overrightarrow{AB}$ par 2 et le vecteur $\overrightarrow{DE}$ est le produit du vecteur $\overrightarrow{IA}$ par $-4$.}}
  	\end{enumerate}
  	\dures
  	 
  	\defs{On appelle produit du vecteur non nul $\overrightarrow{AB}$  par un nombre réel non nul $k$, le vecteur $\overrightarrow{MN}$ , noté $k\overrightarrow{AB}$  tel que :
  		
  		\ding{43} Les droites $(MN)$ et $(AB)$ ont la même direction, 
  		
  		\ding{43} $\overrightarrow{MN}$ et $\overrightarrow{AB}$ sont de même sens lorsque $k>0$
  		
  		\ding{43} $\overrightarrow{MN}$ et $\overrightarrow{AB}$ sont de sens contraires lorsque $k<0$
  		
  		\ding{43} $MN=\mid k\mid AB$.
  		
  		%	La norme du vecteur $\overrightarrow{MN}$ est $\Arrowvert\overrightarrow{MN}\| = MN$  et celle de $k\overrightarrow{MN}$ est $\Arrowvert k\overrightarrow{MN}\|=\mid k\mid MN$
  	}
  	\nbVo{Par conventions}{$\bullet$ Le produit du vecteur nul par un nombre réel est le vecteur 
  		nul noté $\overrightarrow{0}$;
  		
  		$\bullet$ Le produit d’un vecteur quelconque par $0$ est le vecteur nul.}
  	
  	\Ovalbox{\textbf{Consigne 3}} \textbf{Propriétés }
  	
  	On considère la figure de l'activité 1.
  	
  	Simplifie les écritures suivantes en complétant les égalités suivantes:
  	\begin{enumerate}
  		\item $2\overrightarrow{AB}+ 3\overrightarrow{AB}=(2+\dots)\dots=\dots$	
  		\item $-2\left(2\overrightarrow{JK}\right)=(-2 \times \dots)\dots=\dots$
  		\item $4\overrightarrow{AC}+ 4\overrightarrow{AB}=4\left(\dots+\dots\right)$
  	\end{enumerate}
  	\dures
  	 
  	\propr{$\overrightarrow{AB}$ et  $\overrightarrow{DC}$ sont deux vecteurs, $k$ et $k'$ sont deux nombres réels.
  		
  		\ding{118} $k\left(k'\overrightarrow{AB}\right)=(kk')\overrightarrow{AB}$ 
  		
  		\ding{118} $k\overrightarrow{AB}+k'\overrightarrow{AB}=(k+k')\overrightarrow{AB}$
  		
  		\ding{118} $k\overrightarrow{AB}+k\overrightarrow{DC}=k\left(\overrightarrow{AB}+\overrightarrow{DC}\right)$
  		
  		\ding{118} $\overrightarrow{AB}=-\overrightarrow{BA}$ (vecteurs opposées)
  		
  		\ding{118} $\overrightarrow{AB}+\overrightarrow{BA}=\overrightarrow{AA}=\overrightarrow{0}$}
  	
  	\conAct{Approfondissement }
  	
  	$ABC$ est un triangle.
  	\begin{enumerate}
  		\item Écris plus simplement les vecteurs suivants  
  		\begin{enumerate}
  			\item $\overrightarrow{V}=\overrightarrow{AC}+2\overrightarrow{CB}-\overrightarrow{CA}$
  			\item $\overrightarrow{U}=3\overrightarrow{AB}-3\overrightarrow{DB}+3\overrightarrow{DA}=\overrightarrow{0}$
  		\end{enumerate}
  		\item  Construis les points $I$ et $J$ tels que $\overrightarrow{AB}=2\overrightarrow{AI}$ , $\overrightarrow{JC}=\overrightarrow{AJ}$
  		et $\overrightarrow{W}=\overrightarrow{AI}-\overrightarrow{AJ}$
  	\end{enumerate}
  	\dures
  	 
  	\act{Vecteurs colinéaires-vecteurs directeur d'une droite}
  	
  	
  	\conAct{Vecteurs colinéaires et vecteurs directeur d'une droite }
  	
  	\begin{enumerate}
  		
  		\item Trouve un vecteurs ayant la même direction que le vecteur $\overrightarrow{DE}$
  		
  		\bcinfo{\textbf{Information}: \textbf{On dit que ces vecteurs sont colinéaires.}}
  		\item Cite sur la figure trois vecteurs colinéaires à  $\overrightarrow{AI}$.
  		\item Exprime le vecteur $\overrightarrow{DE}$ en fonction du vecteur.
  		
  		\item ) Cite trois vecteurs de même direction que la droite $(\mathcal{D})$
  		
  		\bcinfo{\textbf{Information}: \textbf{On dit que ces vecteurs sont des vecteurs directeurs de la droite $(\mathcal{D})$ }}
  	\end{enumerate}
  	\dures
  	 
  	\defs{On dit que deux vecteurs sont colinéaires lorsqu'ils ont même direction ou lorsque l’un d’eux est le vecteur nul.}
  	\nbVo{Remarque}{Le vecteur nul ($\overrightarrow{0})$ est colinéaire à tout vecteur de du plan.}
  	\defs{On dit qu’un vecteur non nul $\overrightarrow{AB}$ est un vecteur directeur d’une droite $(\mathcal{D})$ lorsque les droites $(AB)$ et $ (\mathcal{D})$ ont la même direction}
  	
  	\conAct{Caractérisation vectorielle de la colinéarité du parallélisme et de points alignés }
  	
  	On donne les points $A$, $L$ et $B$; puis les droites $(AB)$ et $(CD)$
  	\begin{enumerate}
  		\item On suppose qu'il existe un nombre réel $k$ tel que $\overrightarrow{AB}=k\overrightarrow{CD}$.
  		\begin{enumerate}
  			\item Justifie que les vecteurs $\overrightarrow{AB}$ et $\overrightarrow{CD}$ sont colinéaires.
  			\item Justifie que les droites $(AB)$ et $(CD)$ sont parallèles.
  		\end{enumerate}
  		\item On suppose que les vecteurs $\overrightarrow{AB}$ et $\overrightarrow{CD}$ sont colinéaires, soit $(AB)\parallel(CD)$.
  		
  		Justifie qu'il existe un nombre réel $k$ tel que $\overrightarrow{AB}=k\overrightarrow{CD}$.
  		\item 
  		\begin{enumerate}
  			\item On suppose qu'il existe un nombre réel $k$ tel que $\overrightarrow{AL}=k\overrightarrow{AB}$. Justifie que les points $A$, $L$ et $B$ sont alignés.
  			\item  On suppose que les points $A$, $L$ et $B$ sont alignés. Justifie qu'il existe un nombre réel $k$ tel que $\overrightarrow{AL}=k\overrightarrow{AB}$.
  		\end{enumerate}
  	\end{enumerate}
  	\dures\\

  	\propr{ $P_{1}:$ $\overrightarrow{AB}$ et $\overrightarrow{CD}$ sont deux vecteurs et $\overrightarrow{CD}$ n'est pas le vecteur nul.
  		
  		$\bullet$ Si les vecteurs $\overrightarrow{AB}$ et $\overrightarrow{CD}$ sont colinéaires alors il existe un nombre réel $k$ tel que $\overrightarrow{AB}=k\overrightarrow{CD}$.
  		
  		$\bullet$ S’il existe un nombre réel $k$ tel que : $\overrightarrow{AB}=k\overrightarrow{CD}$, alors les vecteurs  $\overrightarrow{AB}$ et $\overrightarrow{CD}$ sont colinéaires.
  		
  		$P_{2}:$  $(AB)$ et $(CD)$ sont deux droites du plan.
  		
  		$\bullet$ Si les droites  $(AB)$ et $(CD)$ sont parallèles, alors il existe 
  		un nombre réel $k$ tel que $\overrightarrow{CD}=k\overrightarrow{AB}$.
  		
  		$\bullet$ S’il existe un nombre réel $k$ tel que $\overrightarrow{CD}=k\overrightarrow{AB}$ alors les droites  $(AB)$ et $(CD)$ sont parallèles.
  		
  		$P_{2}:$ Soit $A$, $B$ et $M$ trois  points distincts du plan.
  		
  		$\bullet$ Si les points $A$, $B$ et $M$  sont alignés alors il existe un 
  		nombre réel $k$ tel que $\overrightarrow{AM}=k\overrightarrow{AB}$.
  		
  		$\bullet$ S’il existe un nombre réel $k$ tel que $\overrightarrow{AM}=k\overrightarrow{AB}$ alors les 
  		points$A$, $B$ et $M$ sont alignés.
  	}
  	\conAct{Approfondissement}
  	
  	$A$, $B$, $C$, $D$, $E$, $K$ et $T$ sont des points du plan tels que $\overrightarrow{AT}=\dfrac{3}{2}\overrightarrow{DC}$; $\overrightarrow{KD}=5\overrightarrow{AE}$ et $\overrightarrow{BT}=\dfrac{1}{2}\overrightarrow{KC}-\overrightarrow{CD}$
  	\begin{enumerate}
  		\item Démontre que les vecteurs $\overrightarrow{AB}$ et $\overrightarrow{DK}$ sont colinéaires.
  		\item Déduis-en que la position relative des droites $(AB)$ et $(DK)$
  		\item Justifie que les points $A$, $B$ et $E$ sont alignés.
  	\end{enumerate}
  	
  	\dures 
  	
  	\act{Vecteurs orthogonaux}
  	
  	Soit $ABCD$ un rectangle.
  	
  	\begin{tikzpicture}[scale=0.6]
  		\draw [line width=0.8pt] (1,1)-- (7,1);
  		\draw [line width=0.8pt] (7,1)-- (7,4);
  		\draw [line width=0.8pt] (7,4)-- (1,4);
  		\draw [line width=0.8pt] (1,4)-- (1,1);
  		\draw [line width=0.8pt] (1.4,4)-- (1.4,3.6);
  		\draw [line width=0.8pt] (1.4,3.6)-- (1,3.6);
  		\draw [line width=0.8pt] (1,1.404305646997127)-- (1.4,1.4);
  		\draw [line width=0.8pt] (1.4,1.4)-- (1.4,1);
  		\draw [line width=0.8pt] (6.6,1.4)-- (6.6,1);
  		\draw [line width=0.8pt] (6.6,1.4)-- (7,1.4);
  		\draw (0.15,1.5) node[anchor=north west] {$A$};
  		\draw (7,1.5) node[anchor=north west] {$B$};
  		\draw (7.1,4.4) node[anchor=north west] {$C$};
  		\draw (0.1,4.6) node[anchor=north west] {$D$};
  	\end{tikzpicture}
  	\begin{enumerate}
  		\item Cite deux droites perpendiculaires.
  		\item Détermine un vecteur directeur de chacune de ces deux droites.
  		
  		\bcinfo{\textbf{Information}: \textbf{On dit que ces vecteurs sont orthogonaux.}}
  	\end{enumerate} 
  	\dures
  	 
  	\defs{On dit que deux vecteurs non nuls sont orthogonaux, lorsqu'ils sont 
  		des vecteurs directeurs de deux droites perpendiculaires.}
  	\nbVo{Remarque}{Le vecteur nul est orthogonal à tout vecteur du 
  		plan. }
  	\nbVo{Notation}{Si $\overrightarrow{AB}$ et $\overrightarrow{DC}$ sont deux vecteurs orthogonaux, On note: $\overrightarrow{AB} \perp \overrightarrow{DC}$.}	
  	\begin{center}
  		\Ovalbox{\textbf{ Exercice de maison}}
  	\end{center}
  	\begin{center}
  		\Ovalbox{\textbf{ Exercice 1 }}
  	\end{center}
  	$RST$ est un triangle. On désigne par $I$ le milieu de $[RT]$ et $J$ le milieu de $[ST]$.
  	\begin{enumerate}
  		\item  Démontre  que les droites $(IJ)$ et $(RS)$ sont parallèles.
  		\item Donne quatre vecteurs directeurs de la droite $(SR)$.
  		\item Donne trois vecteurs directeurs de la droite $(RT)$.
  	\end{enumerate}
  	\begin{center}
  		\Ovalbox{\textbf{ Exercice 2 }}
  	\end{center}
  	On considère les points $A$, $B$, $C$ $I$ et $E$ tels que $\overrightarrow{AB} \perp \overrightarrow{AC}$; $\overrightarrow{AI}=\dfrac{3}{2}\overrightarrow{AB}$ et $\overrightarrow{BE}=\dfrac{1}{2}\overrightarrow{AC}$.
  	
  	\begin{enumerate}
  		\item  Construis les points $A$, $B$, $C$, $I$ et $E$ .
  		\item Justifie que les vecteurs $\overrightarrow{IC}$ et $\overrightarrow{IE}$ sont colinéaires et déduis-en que les points $C$,$I$ et $E$ sont alignés.
  		\item Justifie que les vecteurs $\overrightarrow{EB}$ et $-3\overrightarrow{AB}$ sont orthogonaux.
  	\end{enumerate}
  	\begin{center}
  		\Ovalbox{\textbf{ Exercice 3 }}
  	\end{center}
  	On donne le triangle $ABC$ tel que $AB=5cm$; $BC=6cm$ et $AC=4cm$
  	\begin{enumerate}
  		\item  Construis le triangle $ABC$ .
  		\item Place les points $R$ et $S$ tel que $\overrightarrow{AS}=\dfrac{3}{2}\overrightarrow{AB}$ et $\overrightarrow{BR}=2\overrightarrow{AC}$.
  		\item Calcule $\Arrowvert\overrightarrow{AS}\|$ et $\Arrowvert\overrightarrow{BR}\|$.
  	\end{enumerate}
   \sequence{CALCULS SUR LES COORDONNEES DE VECTEURS}
  	\setcounter{q}{0}
  	\act{ Coordonnées d'un vecteur}
  	Afin d'évaluer les distances entre les différentes positions occupées par sa machine au cours de l'épandage, le technicien se propose de les placer dans un repère qu'il s'est choisi.
  	
  	\conAct{Abscisse d'un point sur une droite} \\
  	Observe ci-dessous la droite $(\mathscr{D})$ muni du repère $(O, I)$.\\
  	\begin{minipage}{6cm}
  		\begin{tikzpicture}[scale=0.8]
  			\clip(-4.659411298389018,-0.5507863055794087) rectangle (6.696864649217748,0.5848411643063186);
  			\draw (5.317559878253363,0.473307792605) node[anchor=north west] {$(\mathscr{D})$};
  			\draw [line width=0.6pt,domain=-4.659411298389018:6.696864649217748] plot(\x,{(-0-0*\x)/8});
  			\draw (0.8,0.019231104919507497) node[anchor=north west] {$1$};
  			\draw (-0.2,-0.004274767678385951) node[anchor=north west] {0};
  			\draw (-1.3,-0.005743884715754291) node[anchor=north west] {-1};
  			\begin{scriptsize}
  				\draw [color=black] (-3,0)-- ++(-1pt,-1pt) -- ++(2pt,2pt) ++(-2pt,0) -- ++(2pt,-2pt);
  				\draw[color=black] (-2.8816407047015895,0.1389751292835) node {$A$};
  				\draw [color=black] (5,0)-- ++(-1pt,-1pt) -- ++(2pt,2pt) ++(-2pt,0) -- ++(2pt,-2pt);
  				\draw[color=black] (5.118326966891841,0.1389751292835) node {$B$};
  				\draw [color=black] (1,0)-- ++(-1pt,-1pt) -- ++(2pt,2pt) ++(-2pt,0) -- ++(2pt,-2pt);
  				\draw[color=black] (1.1183431310951257,0.1389751292835) node {$I$};
  				\draw [color=black] (-2,0)-- ++(-1pt,-1pt) -- ++(2pt,2pt) ++(-2pt,0) -- ++(2pt,-2pt);
  				\draw[color=black] (-1.8701505393277074,0.1389751292835) node {$D$};
  				\draw [color=black] (0,0)-- ++(-1pt,-1pt) -- ++(2pt,2pt) ++(-2pt,0) -- ++(2pt,-2pt);
  				\draw[color=black] (-0.061728728507736606,0.1301455029669) node {$O$};
  				\draw [color=black] (-1,0)-- ++(-1pt,0 pt) -- ++(2pt,0 pt) ++(-1pt,-1pt) -- ++(0 pt,2pt);
  				\draw [color=black] (2,0)-- ++(-1pt,0 pt) -- ++(2pt,0 pt) ++(-1pt,-1pt) -- ++(0 pt,2pt);
  				\draw [color=black] (3,0)-- ++(-1pt,0 pt) -- ++(2pt,0 pt) ++(-1pt,-1pt) -- ++(0 pt,2pt);
  				\draw [color=black] (4,0)-- ++(-1pt,0 pt) -- ++(2pt,0 pt) ++(-1pt,-1pt) -- ++(0 pt,2pt);
  			\end{scriptsize}
  		\end{tikzpicture}
  		
  	\end{minipage}
  	\begin{enumerate}
  		\item Écris le vecteur $\overrightarrow{OA}$ en fonction du vecteur $\overrightarrow{OI}$
  		\item Écris le vecteur $\overrightarrow{OB}$ en fonction du vecteur $\overrightarrow{OI}$
  	\end{enumerate}
  	\dure 
  	\resul
  	\begin{enumerate}
  		\item Écrivons le vecteur $\overrightarrow{OA}$ en fonction du vecteur $\overrightarrow{OI}$\\
  		Les vecteurs $ \overrightarrow{OA} $ et $ \overrightarrow{OI} $ sont colinéaires de sens contraire.
  		On a donc $\overrightarrow{OA}=-3\overrightarrow{OI}$
  		\item Écrivons le vecteur $\overrightarrow{OB}$ en fonction du vecteur $\overrightarrow{OI}$\\
  		Les vecteurs $ \overrightarrow{OB} $ et $ \overrightarrow{OI} $ sont colinéaires de même sens  .
  		On a $\overrightarrow{OB}=5\overrightarrow{OI}$
  	\end{enumerate}
  	\defs{ La droite $(\mathscr{D})$ est munie du repère $(O, I)$. On appelle abscisse d'un point $M$ de $(\mathscr{D})$ dans le repère $(O, I)$, le nombre réel $x$ tel que $\overrightarrow{OM} = x\overrightarrow{OI}$.}
  	\conAct{Les coordonnées d'un point}
  	Le plan est muni d'un repère $(O, I, J)$; $M$ est un point du plan. Observe la figure.\\
  	\begin{minipage}{5cm}
  		\begin{tikzpicture}[scale=0.8]
  			\clip(6,3) rectangle (11,7);
  			\draw [line width=0.6pt,domain=5.2:14.88] plot(\x,{(-33.82--5.54*\x)/1.24});
  			\draw [line width=0.6pt,domain=5.2:14.88] plot(\x,{(--15.36--0.56*\x)/4.82});
  			\draw (9.66,4.28) node[anchor=north west] {$x$};
  			\draw (10,5) node[anchor=north west] {$ M_1 $};
  			\draw [line width=0.6pt] (10.38,6.3)-- (9.889954014856738,4.335762292182525);
  			\draw [line width=0.6pt] (10.38,6.3)-- (7.4333639900955095,5.936158471878317);
  			\draw (7.5,6) node[anchor=north west] {$M_2 $};
  			\draw (6.8,6.18) node[anchor=north west] {$ y $};
  			\draw [->,line width=0.6pt](7,4)--(10.38,6.3);
  			\begin{scriptsize}
  				\draw [color=black] -- ++(-1pt,-1pt) -- ++(2pt,2pt) ++(-2pt,0) -- ++(2pt,-2pt);
  				\draw[color=black] (7.2,3.85) node {$O$};
  				\draw [color=black] -- ++(-1pt,-1pt) -- ++(2pt,2pt) ++(-2pt,0) -- ++(2pt,-2pt);
  				\draw[color=black] (10.54,6.61) node {$M$};
  				\draw [color=black] (7.963318004952247,4.111920764060842)-- ++(-1pt,-1pt) -- ++(2pt,2pt) ++(-2pt,0) -- ++(2pt,-2pt);
  				\draw[color=black] (8.12,4.23) node {$I$};
  				\draw [color=black] (7.216681995047756,4.968079235939158)-- ++(-1pt,-1pt) -- ++(2pt,2pt) ++(-2pt,0) -- ++(2pt,-2pt);
  				\draw[color=black] (7.38,5) node {$J$};
  				\draw [color=black] (9.889954014856738,4.335762292182525)-- ++(-1pt,-1pt) -- ++(2pt,2pt) ++(-2pt,0) -- ++(2pt,-2pt);
  				\draw [color=black] (7.4333639900955095,5.936158471878317)-- ++(-1pt,-1pt) -- ++(2pt,2pt) ++(-2pt,0) -- ++(2pt,-2pt);
  			\end{scriptsize}
  		\end{tikzpicture}
  	\end{minipage}
  	\hspace*{-1.5cm} \begin{minipage}{5cm}
  		\begin{enumerate}
  			\item Écris le vecteur $\overrightarrow{OM_1}$ en fonction du vecteur $\overrightarrow{OI}$
  			\item Écris le vecteur $\overrightarrow{OM_2}$ en fonction du vecteur $\overrightarrow{OJ}$
  			\item Dis ce que vaut la somme $\overrightarrow{OM_1} + \overrightarrow{M_2}$
  			\item Déduis-en que $\overrightarrow{OM} = x\overrightarrow{OI} + y\overrightarrow{OJ}$
  		\end{enumerate}
  	\end{minipage}
  	\defs{1. Le plan est muni du repère $(O, I, J)$, $M$ est un point du plan. On appelle coordonnées du point $M$ dans le repère $(O, I, J)$ le couple $(x, y)$ de nombres réels vérifiant l'égalité : $\overrightarrow{OM} = x\overrightarrow{OI} + y\overrightarrow{OJ}$. $x$ est appelé l'abscisse et $y$ l'ordonnée du point $M$ dans le repère $(O, I, J)$. On note $M(x, y)$ ou $M \begin{pmatrix} x \\ y \end{pmatrix}$. La droite $(OI)$ muni du repère $(O, I)$ est appelée \textbf{axes des abscisses} et la droite $(OJ)$ muni du repère $(O, J)$ est appelé \textbf{axe des ordonnées}.
  		
  		2. Le repère $(O, I, J)$ est dit \textbf{orthonormé} lorsque $(OI) \perp (OJ)$ et $OI = OJ = 1$.}
  	
  	\conAct{Application} 
  	Dans le plan muni d'un repère $(O, I, J)$, les points $A, B, C, D, E$ sont tels que :
  	
  	$ \overrightarrow{OA} = 2\overrightarrow{OI} + 2\overrightarrow{OJ}$,  $  \overrightarrow{OD} = 3\overrightarrow{OJ} $,  $
  	\overrightarrow{OB} = -2\overrightarrow{OI} $,  $ \overrightarrow{OE} = -3\overrightarrow{OI} - \overrightarrow{OJ} $,  $
  	\overrightarrow{OC} = \overrightarrow{OI} - \overrightarrow{OJ} $ et $ \overrightarrow{OG} = -\overrightarrow{OI}
  	$
  	1. Donnons le couple de coordonnées de chacun des points dans le repère $(O, I, J)$. \\
  	
  	2. Place ces points dans le plan muni du repère orthonormé $(O, I, J)$. \\
  	3. Construis les vecteurs $\overrightarrow{AB}$ et $\overrightarrow{DE}$.\\
  	\dures
  	 
    \conAct{Coordonnées d'un vecteur dans le plan }
  	\par Le plan est muni d'un repère $(O, I, J)$. On considère les points $ A $ et $B$ tel que les points $ A_1 $ et $ B_1 $ sont les projétés respectifs des points  $ A $ et $B$  sur l'axe des abscisses parallèlement à l'axe des ordonnées et les points $ A_2 $ et $ B_2 $ sont les projétés respectifs des points  $ A $ et $B$  sur l'axe des ordonnées   parallèlement à l'axe  des abscisses comme l'indique la figure ci-dessous.\\
  	\begin{minipage}{6cm}
  		\begin{tikzpicture}[scale=0.6]
  			\clip(-0.5,-0.5) rectangle (7,4);
  			\draw [line width=0.8pt,domain=-6.33:16.33] plot(\x,{(-0-0*\x)/5});
  			\draw [line width=0.8pt,domain=-6.33:16.33] plot(\x,{(-0--4*\x)/3});
  			\draw [->,line width=0.8pt] (4.8,2.4) -- (6.4,3.2);
  			\draw [line width=1.2pt,dash pattern=on 1pt off 1pt] (1.8,2.4)-- (4.8,2.4);
  			\draw [line width=1.2pt,dash pattern=on 1pt off 1pt] (4.8,2.4)-- (3,0);
  			\draw [line width=1.2pt,dash pattern=on 1pt off 1pt] (2.4,3.2)-- (6.4,3.2);
  			\draw [line width=1.2pt,dash pattern=on 1pt off 1pt] (6.4,3.2)-- (4,0);
  			\draw [line width=1.2pt,dash pattern=on 1pt off 1pt] (4.8,2.4)-- (5.8,2.4);
  			\begin{scriptsize}
  				\draw [color=black] (0,0)-- ++(-1pt,-1pt) -- ++(2pt,2pt) ++(-2pt,0) -- ++(2pt,-2pt);
  				\draw[color=black] (-0.29,0.29) node {$O$};
  				\draw [color=black] (1,0)-- ++(-1pt,-1pt) -- ++(2pt,2pt) ++(-2pt,0) -- ++(2pt,-2pt);
  				\draw[color=black] (1.05,-0.19) node {$I$};
  				\draw [color=black] (0.6,0.8)-- ++(-1pt,-1pt) -- ++(2pt,2pt) ++(-2pt,0) -- ++(2pt,-2pt);
  				\draw[color=black] (0.37,1.09) node {$J$};
  				\draw [color=black] (1.8,2.4)-- ++(-1pt,-1pt) -- ++(2pt,2pt) ++(-2pt,0) -- ++(2pt,-2pt);
  				\draw[color=black] (1.45,2.55) node {$A_2$};
  				\draw [color=black] (2.4,3.2)-- ++(-1pt,-1pt) -- ++(2pt,2pt) ++(-2pt,0) -- ++(2pt,-2pt);
  				\draw[color=black] (2.2,3.37) node {$B_2$};
  				\draw [color=black] (3,0)-- ++(-1pt,-1pt) -- ++(2pt,2pt) ++(-2pt,0) -- ++(2pt,-2pt);
  				\draw[color=black] (3.01,-0.25) node {$A_1$};
  				\draw [color=black] (4,0)-- ++(-1pt,-1pt) -- ++(2pt,2pt) ++(-2pt,0) -- ++(2pt,-2pt);
  				\draw[color=black] (4.15,-0.25) node {$B_1$};
  				\draw [color=black] (4.8,2.4)-- ++(-1pt,-1pt) -- ++(2pt,2pt) ++(-2pt,0) -- ++(2pt,-2pt);
  				\draw[color=black] (4.61,2.59) node {$A$};
  				\draw [color=black] (6.4,3.2)-- ++(-1pt,-1pt) -- ++(2pt,2pt) ++(-2pt,0) -- ++(2pt,-2pt);
  				\draw[color=black] (6.57,3.41) node {$B$};
  				\draw [color=black] (5.8,2.4)-- ++(-1pt,-1pt) -- ++(2pt,2pt) ++(-2pt,0) -- ++(2pt,-2pt);
  				\draw[color=black] (5.99,2.39) node {$C$};
  			\end{scriptsize}
  		\end{tikzpicture}
  	\end{minipage}
  	\hspace*{-1.5cm}\begin{minipage}{4cm}
  		1. Dis ce que vaut la somme $\overrightarrow{AC} + \overrightarrow{CB}$. \\
  		2. Cite un vecteur égal à chacun des vecteurs $\overrightarrow{AC}$ et $\overrightarrow{CB}$ sur la figure. 
  	\end{minipage}\\
  	3. Justifie que $\overrightarrow{A_1B_1}$ et $\overrightarrow{OI}$ d'une part puis $\overrightarrow{A_2B_2}$ et $\overrightarrow{OJ}$ d'autre part sont colinéaires. \\
  	4. Justifie qu'il existe un couple $ (x;y) $ tel que  $\overrightarrow{AB}=x\overrightarrow{OI}+ y\overrightarrow{OJ}$\\
  	\dures
  	 
  	\defs{Le plan est muni du repère $(O, I, J)$, $A$ et $B$ sont des points du plan. On appelle coordonnées du vecteur $\overrightarrow{AB}$ (ou composantes scalaires du vecteur $\overrightarrow{AB}$ dans le repère $(O, I, J)$) le couple de nombres réels $(x, y)$ vérifiant l'égalité $\overrightarrow{AB} = x\overrightarrow{OI} + y\overrightarrow{OJ}$. On note $AB(x, y)$ ou $\overrightarrow{AB} \begin{pmatrix} x \\ y \end{pmatrix}$ pour exprimer que les coordonnées du vecteur $\overrightarrow{AB}$ dans le repère $(O, I, J)$ sont les nombres réels $x$ et $y$ pris dans cet ordre. \nbVo{NB}{Contrairement à un point du plan, on ne parlera de l'abscisse ni de l'ordonnée pour un vecteur. On préférera les expressions "première coordonnée et deuxième coordonnée".}}
  	\conAct{Application } 
  	Dans le plan muni d'un repère $(O, I, J)$On donne
  	\\
  	$ \overrightarrow{BA} = 2\overrightarrow{OI} + 2\overrightarrow{OJ}$,  $  \overrightarrow{ED} = 3\overrightarrow{OJ} $,  $
  	\overrightarrow{FB} = -2\overrightarrow{OI} $ et $ \overrightarrow{MG} = -\overrightarrow{IO}
  	$
  	\\ Détermine les composantes de chacun de ces vecteurs.
  	\dures
  	
  	\act{Calcul dans un repère orthonormé}
  	\conAct{ Égalité de deux vecteurs}
  	
  	$A$ et $B$ sont deux points du plan muni du repère $(O, I, J)$ avec $A(x_A, y_A)$ et $B(x_B, y_B)$ \\
  	1. Écris le vecteur $\overrightarrow{AB}$ en fonction des vecteurs $\overrightarrow{OA}$ et $\overrightarrow{OB}$. \\
  	2. Écris le vecteur $\overrightarrow{AB}$ en fonction des vecteurs $\overrightarrow{OI}, \overrightarrow{OJ}$ et des coordonnées des points $A$ et $B$.
  	\ret{Soit $A(x_A, y_A)$ et $B(x_B, y_B)$ deux points du plan muni d'un repère $(O, I, J)$. $\overrightarrow{AB}$ a pour coordonnées :
  		$ \overrightarrow{AB} \begin{pmatrix} x_B - x_A \\ y_B - y_A \end{pmatrix} $}
  	
  	\conAct{Propriétés}
  	$A, B, C, D$ sont des points du plan muni du repère $(O, I, J)$ tels que $\overrightarrow{AB}(x, y)$ et $\overrightarrow{CD}(x', y')$. \\
  	1. Écris chacun des vecteurs $\overrightarrow{AB}$ et $\overrightarrow{CD}$ en fonction des vecteurs $\overrightarrow{OI}$ et $\overrightarrow{OJ}$. \\
  	2. Déduis-en que si les vecteurs $\overrightarrow{AB}$ et $\overrightarrow{CD}$ sont égaux alors on a $(x = x')$ et $(y = y')$. \\
  	3. Écris le vecteur $\overrightarrow{AB} + \overrightarrow{CD}$ en fonction des vecteurs $\overrightarrow{OI}$ et $\overrightarrow{OJ}$ puis déduis les coordonnées de la somme $\overrightarrow{AB} + \overrightarrow{CD}$ dans le repère $(O, I, J)$. \\
  	4. Écris le vecteur $k\overrightarrow{AB}$ en fonction des vecteurs $\overrightarrow{OI}$ et $\overrightarrow{OJ}$ puis déduis-en les coordonnées du vecteur $k\overrightarrow{AB}$ dans le repère $(O, I, J)$. \\
  	5. On note $(x_A, y_A)$ et $(x_B, y_B)$ les couples de coordonnées respectifs des points $A$ et $B$; $I$ milieu du segment $[AB]$. A partir de l'égalité vectorielle $\overrightarrow{AI} = \overrightarrow{IB}$, détermine les coordonnées du point $I$ en fonction des coordonnées de $A$ et $B$.\\
  	 
  	\propr{ P$_1 $. $A$ et $B$ sont deux points du plan muni d'un repère. Si $A(a, b)$ et $B(c, d)$ alors $\overrightarrow{AB}(c-a, d-b)$. \\
  		P$_2 $. $A, B, C, D$ sont des points du plan muni d'un repère. Si $\overrightarrow{AB}(x, y)$ et $\overrightarrow{CD}(x', y')$ alors $(\overrightarrow{AB} + \overrightarrow{CD})(x+x', y+y')$.\\
  		P$_3 $. $A$ et $B$ sont deux points du plan muni d'un repère et $k$ un réel. Si $\overrightarrow{AB}(x, y)$ alors $k\overrightarrow{AB}(kx, ky)$ \\
  		P$_4 $. $A$ et $B$ sont deux points du plan muni d'un repère et $I$ le milieu du segment $[AB]$. Si $A(x, y); B(x', y')$ alors $I\left(\frac{x+x'}{2}; \frac{y+y'}{2}\right)$}
  	
  	\conAct{Application} On donne $C(-2, 0); D(-2, -2); S(5, 4)$. Détermine les coordonnées de $I$ et $J$ milieu respectifs de $[CD]$ et $[DS]$\\
  	\act{Colinéarités  deux vecteurs}
  	\conAct{ Conditions de colinéarité de deux vecteurs}
  	$A, B, C, D$ sont des points du plan muni d'un repère $(O, I, J)$ tels que $\overrightarrow{AB}\begin{pmatrix} x \\ y \end{pmatrix}$ et $\overrightarrow{CD}\begin{pmatrix} x' \\ y' \end{pmatrix}$. On suppose que les vecteurs $\overrightarrow{AB}$ et $\overrightarrow{CD}$ sont colinéaires et non nuls.
  	\begin{enumerate}
  		\item Montre qu'il existe un nombre réel $k$ tel que ($x = kx'$ et $y = ky'$)
  		\item Justifie que  $xy' - yx' = 0$.
  		\item On considère les vecteurs $ \overrightarrow{EF}(2;-1) $ et $ \overrightarrow{GH}(-4;2) $
  		\begin{enumerate}
  			\item Justifie que $ x_{\overrightarrow{EF}}\times y_{\overrightarrow{GH}}-y_{\overrightarrow{EF}}\times x_{\overrightarrow{GH}}=0 $.
  			\item  Justifie que les vecteurs $ \overrightarrow{EF} $ et $ \overrightarrow{GH} $ sont colinéaires.
  		\end{enumerate}
  	\end{enumerate}
  	\dures
  	 
  	\propr{ Le plan est muni du repère $(O, I, J)$, si $\overrightarrow{AB}(x, y)$ et $\overrightarrow{CD}(x', y')$ sont colinéaires alors on a : $xy' - x'y = 0$
  		
  		Le plan est muni d'un repère $(O, I, J)$ on donne $\overrightarrow{AB}(x, y)$ et $\overrightarrow{CD}(x', y')$. Si $xy' - x'y = 0$ alors les vecteurs $\overrightarrow{AB}(x, y)$ et $\overrightarrow{CD}(x', y')$ sont colinéaires}
  	
  	\conAct{Application } Le plan est muni du repère $(O, I, J)$
  	\begin{enumerate}
  		\item On donne $A(4, -6); B(10, 8); C(0, -2); D(3, 5)$
  		\begin{enumerate}
  			\item Les vecteurs $\overrightarrow{AB}$ et $\overrightarrow{CD}$ sont ils colinéaires ?
  			\item Les vecteurs $\overrightarrow{BD}$ et $\overrightarrow{CD}$ sont ils colinéaires ?
  		\end{enumerate}
  		\item On donne $A(5, 3); B(-3, 2); C(0, -4); A'(1, -2); B'(-1, -1)$;  $C'(5, -4)$
  		\begin{enumerate}
  			\item Vérifie si les points $A, B$ et $C$ sont alignés
  			\item Vérifie si $C' \in (A'B')$
  		\end{enumerate}
  	\end{enumerate}
  	\textbf{Stratégie :} TI : ...... min \quad TG : ...... min \quad TC : ...... min
  	\act{Orthogonalité de deux  vecteurs dans un repère orthonormé }
  	\conAct{Distance entre deux point dans un repère orthonormé}
  	
  	On donne $A(x_A, y_A)$ et $B(x_B, y_B)$ dans un repère orthonormée $(O, I, J)$ \\
  	\begin{minipage}{5cm}
  		\begin{tikzpicture}[scale=0.8]
  			\clip(-0.5,-0.5) rectangle (5,6);
  			\draw[line width=0.8pt] (0.21213203435596456,0) -- (0.2121320343559646,0.21213203435596453) -- (0,0.21213203435596456) -- (0,0) -- cycle; 
  			\draw[line width=0.8pt] (4,2.2121320343559647) -- (3.7878679656440353,2.2121320343559647) -- (3.7878679656440353,2) -- (4,2) -- cycle; 
  			\draw [line width=1.2pt,dash pattern=on 1pt off 1pt] (4,5)-- (4,0);
  			\draw [line width=1.2pt,dash pattern=on 1pt off 1pt] (2,0)-- (2,2);
  			\draw [line width=1.2pt,dash pattern=on 1pt off 1pt] (2,2)-- (0,2);
  			\draw [line width=1.2pt,dash pattern=on 1pt off 1pt] (4,5)-- (0,5);
  			\draw [line width=1.2pt,dash pattern=on 1pt off 1pt] (2,2)-- (4,5);
  			\draw [line width=1.2pt,dash pattern=on 1pt off 1pt] (2,2)-- (4,2);
  			\draw [line width=0.8pt,domain=-7.56:14.36] plot(\x,{(-0-0*\x)/4});
  			\draw [line width=0.8pt] (0,-4.88) -- (0,10.58);
  			\begin{scriptsize}
  				\draw [color=black] (2,2)-- ++(-1pt,-1pt) -- ++(2pt,2pt) ++(-2pt,0) -- ++(2pt,-2pt);
  				\draw[color=black] (1.78,2.35) node {$A$};
  				\draw [color=black] (4,5)-- ++(-1pt,-1pt) -- ++(2pt,2pt) ++(-2pt,0) -- ++(2pt,-2pt);
  				\draw[color=black] (4.16,5.31) node {$B$};
  				\draw [color=black] (4,0)-- ++(-0.5pt,-0.5pt) -- ++(1pt,1pt) ++(-1pt,0) -- ++(1pt,-1pt);
  				\draw[color=black] (4,-0.2) node {$x_{B}$};
  				\draw [color=black] (2,0)-- ++(-0.5pt,-0.5pt) -- ++(1pt,1pt) ++(-1pt,0) -- ++(1pt,-1pt);
  				\draw[color=black] (2.08,-0.2) node {$x_A$};
  				\draw [color=black] (0,2)-- ++(-0.5pt,-0.5pt) -- ++(1pt,1pt) ++(-1pt,0) -- ++(1pt,-1pt);
  				\draw[color=black] (-0.34,2.23) node {$y_A$};
  				\draw [color=black] (0,5)-- ++(-0.5pt,-0.5pt) -- ++(1pt,1pt) ++(-1pt,0) -- ++(1pt,-1pt);
  				\draw[color=black] (-0.34,5.19) node {$y_B$};
  				\draw [color=black] (4,2)-- ++(-0.5pt,-0.5pt) -- ++(1pt,1pt) ++(-1pt,0) -- ++(1pt,-1pt);
  				\draw[color=black] (4.16,2.27) node {$C$};
  				\draw [color=black] (0,0)-- ++(-1pt,-1pt) -- ++(2pt,2pt) ++(-2pt,0) -- ++(2pt,-2pt);
  				\draw[color=black] (-0.44,0.27) node {$O$};
  			\end{scriptsize}
  		\end{tikzpicture}
  	\end{minipage}
  	\hspace*{-1cm} \begin{minipage}{5cm}
  		\begin{enumerate}
  			\item Exprime la distance $AC$ en fonction de $y_A$ et $y_B$
  			\item Exprime la distance $BC$ en fonction de $x_A$ et $x_B$
  			\item Sachant que $(a - b)^2 = (b - a)^2$, déduis-en que $AB = \sqrt{(x_B - x_A)^2 + (y_B - y_A)^2}$
  			\item Exprime la distance $AB$ en fonction de $x_{\overrightarrow{AB}}$ et $y_{\overrightarrow{AB}}$
  		\end{enumerate}
  	\end{minipage}
  	\dures
  	 
  	\propr{A et B sont deux points du plan muni d'un repère orthonormé.
  		
  		Si $A(x_A, y_A)$ et $B(x_B, y_B)$ alors \\ $AB = \sqrt{(x_B - x_A)^2 + (y_B - y_A)^2}$ 
  		
  		Si $ \overrightarrow{AB}(a;b) $ alors $ AB=\sqrt{a^2+b^2} $}
  	\conAct{ Conditions d'orthogonalité}
  	
  	On considère le triangle $ ABC $ rectangle en $ C $ tel que $ C(2;2) $, $ A(2;5) $ et $ B(4;2) $.
  	\begin{enumerate}
  		\item Justifie que $ x_{\overrightarrow{AC}}\times x_{\overrightarrow{BC}}+y_{\overrightarrow{AC}}\times y_{\overrightarrow{BC}}=0 $
  		\item Montre que les vecteurs $ \overrightarrow{AC} $ et $ \overrightarrow{BC} $ sont orthogonaux.
  		\item On considère les vecteurs $ \overrightarrow{EF}(x;y) $ et $ \overrightarrow{GH}(x';y') $ orthogonaux et non nuls tel que $ \overrightarrow{OM}=\overrightarrow{EF} $ et $ \overrightarrow{ON}=\overrightarrow{GH} $.
  		\begin{enumerate}
  			\item Justifie que le triangle $ OMN $ est un triangle rectangle en $ O $.
  			\item Montre que $ \overrightarrow{MN}=\overrightarrow{GH}-\overrightarrow{EF} $ puis déduis-en les composantes du vecteur $ \overrightarrow{MN} $.
  			\item Calcule de deux différente manière $ MN^2 $.
  			\item Justifie que $ x\times x'+y\times y'=0 $.
  		\end{enumerate}
  	\end{enumerate}
  	\dures
  	 
  	\propr{$A, B, C$ et $D$ sont des points du plan muni d'un repère orthonormé. Les vecteurs $\overrightarrow{AB}$ et $\overrightarrow{CD}$ ont respectivement pour coordonnées $(x, y)$ et $(x', y')$
  		\begin{enumerate}
  			\item Si les vecteurs $\overrightarrow{AB}$ et $\overrightarrow{CD}$ sont orthogonaux alors $xx' + yy' = 0$
  			\item Si $xx' + yy' = 0$ alors les vecteurs $\overrightarrow{AB}$ et $\overrightarrow{CD}$ sont orthogonaux
  		\end{enumerate}
  	}
  	\conAct{Application } Dans un repère orthonormé $(O, I, J)$ on donne
  	$A(-2; 3); B(3; 1); C(0; -2); M(-1; \frac{2}{3}); N(\frac{1}{3}, 2); P(4; -\frac{5}{3})$
  	\begin{enumerate}
  		\item Calcule $AB$ et $AC$
  		\item Donne la nature du triangle $ABC$
  		\item Montre que $\overrightarrow{MN} \perp \overrightarrow{NP}$
  		\item Détermine les coordonnées du point $ D $ pour que le quadrilatère $ABCD $ soit un parallélogramme 
  		\end{enumerate}
  		\newpage
  			\begin{enumerate}
  				\setcounter{enumi}{4}
  		\item Détermine les coordonnées du point $ E $ pour que $ B $ soit le milieu du segment $ [AE] $
  	\end{enumerate}
  	\textbf{Stratégie :} TI : ...... min \quad TG : ...... min \quad TC : ...... min	
  	\begin{center}
  		EXERCICE DE MAISON
  	\end{center}
  	\Ovalbox{\textbf{ Exercice 1 }}  	
  	Le plan est muni du repère $(O,I,J)$. Soient $A(1,-1)$, $B(2,3)$, $C(1,4)$, $D(-5,3)$ et $E(-6,-1)$.
  	\begin{enumerate}
  		\item Calcule les composantes des vecteurs $\overrightarrow{AB}$ et $\overrightarrow{DC}$.
  		\item Détermine les coordonnées des images des points $A$, $O$ et $E$ par translation de vecteur $\overrightarrow{AB}$.
  		\item Place le point $H$ défini par $\overrightarrow{EH} = \overrightarrow{BA} + \overrightarrow{CB}$.
  		Calcule les coordonnées du point $H$.
  	\end{enumerate}
  \Ovalbox{\textbf{ Exercice 2 }}
  	 Le plan est muni d'un repère orthonormé $(O,I,J)$. $A$, $B$ et $C$ sont des points du plan tels que :
  	$$\overrightarrow{OA} = 7\overrightarrow{OI}+4\overrightarrow{OJ}\,, \qquad \overrightarrow{BO} = -8\overrightarrow{OI}-4\overrightarrow{OJ}\,, \qquad \overrightarrow{OC} = 7\overrightarrow{OJ}+\overrightarrow{IO}.$$
  	\begin{enumerate}
  		\item[1.] (a) Détermine les coordonnées des points $A$, $B$ et $C$ puis place ces points dans le repère $(O,I,J)$.
  		
  		(b) Calcule les distances $AB$, $AC$ et $BC$.
  		
  		(c) Donne la nature du triangle $ABC$ en justifiant ta réponse.
  		
  		(d) Détermine les coordonnées du point $E$ tel que $ABCE$ soit un parallélogramme.
  		\item $H$ est le point du plan tel que $H$ soit le projeté orthogonal de $B$ sur la droite $(AC)$.
  		Calcule les longueurs $AH$ et $BH$.
  		\item Soit $(\mathscr{C})$ le cercle circonscrit au triangle $ABC$.
  		\begin{enumerate}
  			\item[(a)] Justifie que $\text{mes}\,\widehat{EBC} = \text{mes}\,\widehat{EAC}$.
  			\item[(b)] Justifie que le point $O$ appartient au cercle $(\mathscr{C})$ puis détermine $\text{mes}\,\widehat{EOC}$ et $\text{mes}\,\widehat{BOC}$ sachant que $\text{mes}\,\widehat{EBC} = 30°$.
  			\item[(c)] Démontre que les triangles $BEC$ et $AEC$ sont semblables.
  		\end{enumerate}
  	\end{enumerate}
  	
 \Ovalbox{\textbf{ Exercice 3 }}
  	Le plan est muni du repère orthonormé $(O,I,J)$.
  	\begin{enumerate}
  		\item Place les points $A(-3,0)$, $B(2,1)$, $C(4,3)$ et $D(-1,2)$.
  		\item Montre que les segments $[AC]$ et $[BD]$ ont le même milieu $K$.
  		\item Montre que le triangle $OBD$ est rectangle et isocèle.
  		\item Soit $E$ le point tel que $BODE$ soit un parallélogramme.
  		Quelles sont les coordonnées de $E$ ?
  		\item Quelle est la nature du quadrilatère $AOCE$ ?
  	\end{enumerate}
  	
  \Ovalbox{\textbf{ Exercice 4 }}	
  	Le plan est muni du repère $(O,I,J)$. On donne : $A(1,-1)$, $B(-1,-2)$ et $C(-2,2)$.
  	\begin{enumerate}
  		\item Détermine le couple de coordonnées du point $G$ pour que : $\overrightarrow{GA}+2\overrightarrow{GB}+\overrightarrow{GC} = \vec{0}$.
  		\item Démontre que $\overrightarrow{BG}$ et $\overrightarrow{BD}$ sont colinéaires.
  		\item Justifie que les points $B$, $G$ et $D$ sont alignés.
  	\end{enumerate}
  	
  \Ovalbox{\textbf{ Exercice 5 }}	
  	Dans le plan muni d'un repère orthonormé $(O,I,J)$, on considère les points $A(-2,-1)$ et $B(4,3)$. On note $(\mathscr{C})$ le cercle de diamètre $[AB]$ et $M$ le centre de $(\mathscr{C})$.
  	\begin{enumerate}
  		\item Dessine la figure.
  		\item Calcule les coordonnées de $M$.
  		\item Calcule le rayon du cercle $(\mathscr{C})$ (on donnera la valeur exacte).
  		\item Soit $F$ le point de coordonnées $(3,4)$.
  		Démontre que $F$ est un point du cercle $(\mathscr{C})$.
  		\item Que peut-on dire du triangle $AFB$ ?
  		\item On précise que $FA = \sqrt{50}$ et $FB = \sqrt{2}$.
  		Calcule l'arrondi au degré de l'angle $\widehat{FAB}$.
  	\end{enumerate}
  	
  \Ovalbox{\textbf{ Exercice 6 }} Soient les points $A(2,3)$ ; $B(4,-1)$ et $M(3,-2)$ dans un repère $(O,I,J)$.
  	Calculer les coordonnées des points $C$ et $D$ tels que $ABCD$ soit un parallélogramme de centre $M$.
  	
   \situ{ORGANISATION DES DONNEES}
   \setcounter{sequen}{0}
   \setcounter{q}{0}
 	\underline{\textbf{Situation de départ}} \\
 	\textbf{Texte}:	
 	
 	Pour préparer une course en pirogue, 115 élèves d'un collège de Cotonou ont entrepris de lancer une commande de T-shirts. En vue d'établir le bon de commande, les élèves ont communiqué leur taille à Sovi leur premier responsable. Selon l'information donnée par le fournisseur, les tailles S, M, L, XL et XXL sont destinées aux personnes mesurant respectivement moins de 160 cm, de 160cm à moins de 165 cm, de 165 à moins de 170cm, de 170 à moins de 175cm, de 175 et plus.
 	\par Sovi a voulu mesurer la longueur de cette course. Il a jeté dans la lagune, à un mètre devant lui un caillou. Il s’est alors formé à la 
 	surface de l’eau une succession de rides circulaires concentriques. Sovi s’est intéressé aux positions successives de la première 
 	ride.\\
 	\hspace*{-1.5cm} \begin{minipage}{13cm}
 		\definecolor{wewdxt}{rgb}{0.43137254901960786,0.42745098039215684,0.45098039215686275}
 		\definecolor{xdxdff}{rgb}{0.49019607843137253,0.49019607843137253,1}
 		\definecolor{ududff}{rgb}{0.30196078431372547,0.30196078431372547,1}
 		\begin{tikzpicture}[scale=0.8]
 			\clip(-3,2) rectangle (10,5);
 			\draw [line width=2pt] (-2,3)-- (10,3);
 			\draw [shift={(-1,3)},line width=0.8pt,dash pattern=on 1pt off 1pt]  plot[domain=-0.38131938506331853:0.4636476090008061,variable=\t]({1*2.2840534144367113*cos(\t r)+0*2.2840534144367113*sin(\t r)},{0*2.2840534144367113*cos(\t r)+1*2.2840534144367113*sin(\t r)});
 			\draw [shift={(0.6,3.03)},line width=0.8pt,dash pattern=on 1pt off 1pt]  plot[domain=-0.38131938506331853:0.4636476090008061,variable=\t]({1*2.2840534144367113*cos(\t r)+0*2.2840534144367113*sin(\t r)},{0*2.2840534144367113*cos(\t r)+1*2.2840534144367113*sin(\t r)});
 			\draw [shift={(2.2,3.06)},line width=0.8pt,dash pattern=on 1pt off 1pt]  plot[domain=-0.38131938506331853:0.4636476090008061,variable=\t]({1*2.2840534144367113*cos(\t r)+0*2.2840534144367113*sin(\t r)},{0*2.2840534144367113*cos(\t r)+1*2.2840534144367113*sin(\t r)});
 			\draw [shift={(3.8,3.09)},line width=0.8pt,dash pattern=on 1pt off 1pt]  plot[domain=-0.38131938506331853:0.4636476090008061,variable=\t]({1*2.2840534144367113*cos(\t r)+0*2.2840534144367113*sin(\t r)},{0*2.2840534144367113*cos(\t r)+1*2.2840534144367113*sin(\t r)});
 			\draw [shift={(5.4,3.12)},line width=0.8pt,dash pattern=on 1pt off 1pt]  plot[domain=-0.38131938506331853:0.4636476090008061,variable=\t]({1*2.2840534144367113*cos(\t r)+0*2.2840534144367113*sin(\t r)},{0*2.2840534144367113*cos(\t r)+1*2.2840534144367113*sin(\t r)});
 			\draw [shift={(7,3.15)},line width=0.8pt,dash pattern=on 1pt off 1pt]  plot[domain=-0.38131938506331853:0.4636476090008061,variable=\t]({1*2.2840534144367113*cos(\t r)+0*2.2840534144367113*sin(\t r)},{0*2.2840534144367113*cos(\t r)+1*2.2840534144367113*sin(\t r)});
 			\draw [line width=0.4pt] (-1,3)-- (-0.98,3.62);
 			\draw [line width=0.4pt] (-0.98,3.62)-- (0.2,3.62);
 			\draw [line width=0.4pt] (0.2,3.62)-- (0.2,3);
 			\draw (-0.7,3.56) node[anchor=north west] {1 m};
 			\draw (-2.26,2.92) node[anchor=north west] {\parbox{1.96 cm}{\begin{minipage}{1cm}
 						\hspace*{0.2cm}S \\ (Sovi)
 			\end{minipage} }};
 			\draw [->,line width=0.4pt] (0.6,4.16) -- (1.2840534144367113,3);
 			\draw (0.58,4.58) node[anchor=north west] {Point de contact de pierre avecc la surface de l'eau};
 			\begin{scriptsize}
 				\draw [fill=ududff] (-2,3) circle (0.5pt);
 				\draw [fill=ududff] (10,3) circle (0.5pt);
 				\draw [fill=xdxdff] (-1,3) circle (0.5pt);
 				\draw [color=black] (-0.98,3.62)-- ++(-0.5pt,0 pt) -- ++(1pt,0 pt) ++(-0.5pt,-0.5pt) -- ++(0 pt,1pt);
 				\draw [color=black] (0.2,3.62)-- ++(-0.5pt,0 pt) -- ++(1pt,0 pt) ++(-0.5pt,-0.5pt) -- ++(0 pt,1pt);
 				\draw [fill=wewdxt] (0.2,3) circle (0.5pt);
 				\draw [fill=black] (1.2840534144367113,3) circle (0.5pt);
 				\draw[color=black] (1,2.75) node {$A_0$};
 				\draw [fill=black] (2.883856387779232,3) circle (0.5pt);
 				\draw[color=black] (2.6,2.75) node {$A_1$};
 				\draw [fill=black] (4.483265205796297,3) circle (0.5pt);
 				\draw[color=black] (4.3,2.75) node {$A_2$};
 				\draw [fill=black] (6.082279562192152,3) circle (0.5pt);
 				\draw[color=black] (5.80,2.75) node {$A_3$};
 				\draw [fill=black] (7.683289956699467,3.060949397671566) circle (0.5pt);
 				\draw[color=black] (7.5,2.75) node {$A_4$};
 				\draw [fill=black] (9.279122638209712,3) circle (0.5pt);
 				\draw[color=black] (9,2.75) node {$A_5$};
 			\end{scriptsize}
 			
 		\end{tikzpicture}
 	\end{minipage}
 	\begin{minipage}{13cm}
 		\hspace{0.7cm}	$ \underbrace{ \textcolor{white}{ et par une estime conviction je peux y arriver  il suffit}  }_{\mbox{ \small Position successives de la première ride formée}} $
 	\end{minipage}
 	
 
 	
 \par 	Un dispositif lui a permis de consigner dans un tableau les distances $ SA_1 $, $ SA_2 $, $ SA_3 $, $ SA_4 $, $ SA_5 $ (cf figure) ainsi que les temps correspondants.
 	
 	
 	\begin{tabular}{|p{2.5cm}|c|c|c|c|c|c|}
 		\hline 
 		Distance $ SA_i $ en $ (m) $	& 2 & 3,5 & 4 & 6 & 7 & 8,5 \\ 
 		\hline 
 		Temps en $ (s) $	& 2 & 5 & 6 &  10&12  & 15 \\ 
 		\hline 
 	\end{tabular} 
 	
 	\par Kokou, posté au point d'arrivée de la course à noté que la ride observée est parvenue à son niveau au  bout de 3 min.
 	Sovi se demande quoi faire pour assurer à chaque élève un T-shirt convenable et comment calculer la longueur de la course.	
 	
 	\Ovalbox{\textbf{Tâche}} : Tu vas te construire des connaissances nouvelles en mathématiques.
 	
 	Pour cela, tu auras tout au long de la situation d’apprentissage à :\\ 
 	\ding{43} Exprimer ta perception de chacun des problèmes posés ;\\ 
 	\ding{43} Analyser chacun des problèmes posés ; \\
 	\ding{43} Mathématiser chacun des problèmes posés ; \\
 	\ding{43} Opérer sur l’objet mathématique que tu as identifié pour chaque problème ; \\
 	\ding{43} Améliorer au besoin ta production.\\
 	
 	\vspace*{-1cm}
 	\begin{center}
 		\section*{I-INTRODUCTION}
 	\end{center}  
 	\doublebox{\textbf{Activité 0}} \textbf{ Brainstorming}\\
 	\begin{enumerate}
 		\item Lis le texte de la situation de départ dégage le problème qu'il pose.
 		\item Reformule ce problème en tes propre termes.
 		\item Formule toutes les idées et question que t'inspire le texte de la situation de départ.
 	\end{enumerate}
 	
 	\underline{Stratégie et durée}: 15 min
 	\begin{center}
 		\underline{\textbf{Éléments de réponses}}
 	\end{center}
 	(Voir cahiers de recherche.)
 
 	\begin{center}
 		\section*{II- RÉALISATION}
 	\end{center}
  \sequence{APPLICATION AFFINE}	
 	\act{Découverte} 
 	
 	Pour connaître la longueur de la course, Sovi se propose de rechercher la relation entre la distance d parcourue par la
 	ride et le temps de parcours.\\
 	\conAct{Définition  d'une application affine}
 	\begin{enumerate}
 		\item Pour chaque temps (t), vérifie que la distance est $ d=\dfrac{1}{2}t+1 $.
 		\\ \textbf{\underline{Inforlation:} L'expression $ \dfrac{1}{2}t+1 $ est appelé une application affine}
 		\item Détermine alors la longueur du trajet de la
 		course.
 	\end{enumerate}
 	\dures
 	 
 	\defs{$ a $ et $b$ étant des nombres réels ; on appelle application affine $  f$  de coefficient $ a $ et de terme constant $ b $,   } 
 	\defs{l'application qui à chaque réel $ x $ associe le nombre $ ax+b $.\\
 		On dit que l'application $ f $  est définie sur $ \mathbb{R} $ par $ f(x)=ax+b $
 	}
 	\nbVo{Vocabulaire}{Pour tout nombre réel $ \alpha $, $ f(\alpha) $ est appelé l'image de $ \alpha $ par $ f $ et $ \alpha $  est appelé l'antécédent de $ f(\alpha) $ par $f$}
 	
 	\conAct{Reconnaissance d'une application affine}
 	On donne : $ f(x)=2x+3 $, $ g(x)=x^2-1 $, $ h(x)=\sqrt{3}x $, i(x)=0, $ j(x)=x^3+2 $ et $ k(x)=2 $.
 	\begin{enumerate}
 		\item Identifie parmi ces expressions celles qui sont des applications.
 		\item Pour chaque application identifiée, précise son coefficient et son terme constant.
 		\item Une application affine est définie par $ f(2)=3 $ et $ f(0)=-1 $. Détermine l'expression de $ f $.
 	\end{enumerate}
 	\dures
 	\nbVo{NB}{Soit $ e $ et $ c $ deux nombres réels distincts.\\ Si une application affine est définie par $ f(x)=ax +b $  telle que l'image de $ e $ et $ c $ par $ f $ sont respectivement $ f(e) $ et $ f(c) $ alors on a \\
 		$ a=\dfrac{f(e)-f(c)}{e-c}  $ \\ 
 		$ b=f(e)-a(e) $ ou $ b=f(c)-a(c) $}
 	\conAct{Sens de variation}
 	On considère les application affines suivantes: $ f(x)=2x-1 $, $ g(x)=-x+4 $ et $ h(x)=5 $.
 	\begin{enumerate}
 		\item Range dans l'ordre croissant les réels suivants: 2; 3 et 5.
 		\item Pour chaque application détermine l'image des réels suivants: 2; 3 et 5 puis les comparées.
 		\item Compare le coefficient de chaque application à 0.
 		\item Que peut-on dire d'une application affine dont le coefficient est inférieur à 0, supérieur à 0 ou égal à 0.
 	\end{enumerate}
 	\dures
 	 
 	\defs{$f$ étant une application affine définie sur $ \mathbb{R} $ par $ f(x)=ax+b $,\\
 		\ding{118} Pour $ a $ strictement positif, quand $ x $ augmente,  $ f(x) $ augmente. On dit que l'application $ f $ est strictement croissante sur $ \mathbb{R} $.\\
 		\ding{118} Pour $ a $ strictement négatif, quand $ x $ augmente,  $ f(x) $ diminue. On dit que l'application $ f $ est strictement décroissante sur $ \mathbb{R} $. \\
 		\ding{118} Pour $ a $ nul, quand $ x $ augmente, la valeur de  $ f(x) $ ne change pas. On dit que l'application $ f $ est constante sur $ \mathbb{R} $.
 	}
 	\conAct{Approfondissement}
 	Soit les applications affine suivante: $ f(x)=-3x-1 $, $ g(x)=7x+4 $ et $ h(x)=2$.
 	\begin{enumerate}
 		\item Donne le sens de variation de chacune de ces applications.
 		\item Détermine l'antécédent de $-12$ par $f$, l'image de $4$ par $ g $
 	\end{enumerate}
 	\dures
 	
 	\begin{center}
 		\Ovalbox{\textbf{ Exercice de maison}}
 	\end{center}
  \Ovalbox{\textbf{ Exercice 1}} $f$ est l'application affine définie par :
  $$f(x) = (4\sqrt{2}-6)x + \sqrt{27}.$$
  \begin{enumerate}
  	\item Compare $4\sqrt{2}$ et 6.
  	\item $f$ est-elle croissante ou décroissante ?
  	\item Range par ordre croissant $f\left(\dfrac{3}{2}\right)$ ; $f\left(\dfrac{2}{5}\right)$ ; $f\left(\dfrac{2}{3}\right)$ et $f(2)$.
  \end{enumerate}
  
  \Ovalbox{\textbf{ Exercice 2 }} Trouver l'expression de la fonction affine $f$ telle que :
  $$f\left(\frac{1}{2}\right) = 1 \quad \text{et} \quad f(1) = \frac{1}{4}.$$
  
 \Ovalbox{\textbf{ Exercice 3 }} Soit $f$ la fonction affine telle que $f(x) = -2x+3$.
  \begin{enumerate}
  	\item Calculer $f(5)$ et $f(-2)$.
  	\item Quelle est l'image de 1 par $f$ ?
  	\item Tracer la représentation graphique de $f$.
  \end{enumerate}
  
  \Ovalbox{\textbf{ Exercice 4 }} Représenter graphiquement les fonctions suivantes :
  \begin{enumerate}
  	\item $f(x) = 3x-2$
  	\item $f(x) = \dfrac{2}{3}x$
  	\item $f(x) = -\dfrac{4}{3}x+3$
  	\item $f(x) = 3$
  	\item $f(x) = 1-2x$
  	\item $f(x) = -\dfrac{1}{3}x+2$
  \end{enumerate}
  
 \Ovalbox{\textbf{ Exercice 5 }} Déterminer les fonctions affines définies par :
  \begin{enumerate}
  	\item $f(2)=1$ et $f(3)=5$
  	\item $f(-2)=-3$ et $f(4)=1$
  	\item $f(0)=4$ et $f(-3)=-2$
  	\item $f(2)=3$ et $f(3)=3$
  \end{enumerate}
  
  \Ovalbox{\textbf{ Exercice 6}} Des élèves décident de vendre des croissants à domicile pour financer leur voyage de fin d'année. Le montant facturé comprend le prix des croissants auquel s'ajoute une somme fixe pour la livraison. Le prix facturé est fonction affine $f$ du nombre de croissants.
  On sait que 4 croissants livrés coûtent 19 francs et 10 croissants livrés coûtent 37 francs.
  \begin{enumerate}
  	\item Dessiner la représentation graphique de la fonction $f$. En abscisse, on prendra $1\,cm$ pour 1 croissant, en ordonnée $1\,cm$ pour 5 francs.
  	\item Par lecture graphique, déterminer le prix de 12 croissants livrés.
  	\item Déterminer $a$ et $b$ tels que $f(x) = ax+b$.
  	\item Retrouver par le calcul, le prix de 12 croissants.
  	\item Quel est le montant de la livraison ?
  \end{enumerate}
  
 	\sequence{APPLICATIONS LINEAIRE}
 	\setcounter{q}{0}
 	\act{Découverte} 
 	En outre pour vérifier les résultats de Sovi, Aîcha une de ses camarades a jeté un deuxième caillou dans la lagune. Elle obtient le tableau ci-dessous:
 	
 	\begin{tabular}{|c|c|c|c|c|c|}
 		\hline 
 		d( en m )	& 1 &2,5  & 3 & 5 & 6 \\ 
 		\hline 
 		t( en s )	&  2&  5&6  &10  &12  \\ 
 		\hline 
 	\end{tabular} \\
 	
 	\conAct{Définition}
 	\begin{enumerate}
 		\item Justifie que ce tableau est un tableau de proportionnalité.
 		\item Déduis-en la forme générale de la distance $d$ en fonction du temps $t$
 	\end{enumerate}
 	\dures
 	 
 	\defs{On appelle application linéaire une application affine dont le terme constant est nul}
 	\nbVo{Traduction}{ $ a  $ étant un nombre réel, l'application linéaire $ g $ de coefficient $ a $ est l'application qui, chaque réel $ x $ associe le nombre $ ax $\\
 		\textbf{\underline{Notation}} On note $ g(x)=ax $}
 	\conAct{Reconnaissance et détermination d'une application linéaire}
 	\begin{enumerate}
 		\item  Parmi les application suivante identifie celles qui sont des application linéaire: $ f(x)=-3x $, $ g(x)=7x+4 $, $ h(x)=2$, $ f_1(x)=-x-1 $, $ g_1(x)=x $ et $ h_1(x)=2x$.
 		\item  Détermine l'application linéaire $k$ dont le coefficient est $ -2 $ puis déduis-en son sens de variation.
 	\end{enumerate}
 	\dures 
 	
 	
 	 
 	\conAct{Propriétés}
 	$ g $ étant une application linéaire, $ u $, $ v $ et $ k $ étant des nombres réel quelconques. Démontre que $ g(u+v)=g(u)+g(v) $ et $ g(ku)=kg(u) $. \\
 	\dures
 	 
 	\propr{$ g $ étant une application linéaire, $ u $, $ v $ et $k$ étant des nombres réel quelconques, on a  
 		\begin{dinglist}{45}
 			\item $ g(u+v)=g(u)+g(v) $
 			\item $ g(ku)=kg(u) $
 	\end{dinglist} }
 	\conAct{Application }
 	On désigne par $ g $ l'application linéaire telle que $g(-1)=\sqrt{3}$, $g(-\sqrt{3})=3 $ et $ g(\sqrt{2})=-\sqrt{6} $ et $ f(x)=2x $.
 	\begin{enumerate}
 		\item  Calcule $ g(\sqrt{6}) $ et $ g(1+\sqrt{3}) $ sans calculer le coefficient de $g$
 		\item Représente graphiquement $f$.
 	\end{enumerate}
 	\dures
 	 \nbVo{Remarque}{La représentation graphique d'une application linéaire est une droite qui passe par l'origine du repère.}	
 	
 	\reto 
 	\begin{center}
 		\Ovalbox{\textbf{ Exercice de maison}}
 	\end{center}
  Dans le plan rapporté à un repère orthonormé $(O,I,J)$, on considère l'application linéaire $f$ telle que :
 	$$f(1)+f(2) = 5.$$
 	\begin{enumerate}
 		\item Détermine l'application $f$.
 		\item Représente graphiquement $f$.
 		\item Détermine le nombre réel $f(x^2+5)-f(x^2+10)$.
 	\end{enumerate}
 
 	\sequence{STATISTIQUES}
 	\setcounter{q}{0}
 	\act{Regroupement}
 	Le professeur de sport décide de séparer ses élèves en plusieurs groupes. Pour cela, il mesure la taille en centimètre de ses élèves et trouve :
 	
 	\begin{tabular}{cccccc}
 		\hline
 		156 & 162 & 158 & 171 & 167 & 172 \\
 		164 & 165 & 175 & 162 & 156 & 166 \\
 		167 & 156 & 163 & 165 & 163 & 175 \\
 		169 & 167 & 158 & 167 & 163 & 158 \\
 		158 & 170 & 170 & 175 & 175 & 175 \\
 		\hline
 	\end{tabular} \\
 	Il forme trois (3) groupes :
 	
 	\begin{itemize}
 		\item Le groupe des petits formés des élèves dont la taille appartient à l'intervalle $[150; 160[$.
 		\item Le groupe des moyens formés des élèves dont la taille appartiennent à l'intervalle $[160; 170[$.
 		\item Le groupe des grands formés des élèves dont la taille appartient à l'intervalle $[170; 180[$.
 	\end{itemize}Il obtient ainsi trois classes ($[150; 160[$, $[160; 170[$ et $[170; 180[$) et organise ces données dans un tableau des effectifs et des fréquences.\\
 	\renewcommand{\arraystretch}{1.5}
 	
 	\hspace*{-1cm}\begin{tabular}{|p{1.7cm}|c|c|c|c|}
 		\hline
 		classes & $[150; 160[$ & $[160;170[$ & $[170; 180[$ & Total \\ \hline
 		Effectifs & & & & \\ \hline
 		Fréquences en \% & & & & \\ 
 		\hline
 	\end{tabular}
 	
 	\conAct{Définition de la classe modale}
 	\begin{enumerate}
 		\item Quel est l'effectif total de cette série statistique ? 
 		\item Calcule l'amplitude de chaque classe.
 		\item Complète le tableau de l'activité.
 		\item Quelle est la classe ayant l'effectif le plus élevé ?
 	\end{enumerate}
 	\dures
 	 
 	\defs{Une classe modale d'une série statistique est une classe correspondant à l'effectif le plus élevé.}
 	
 	\nbVo{Rappel}{
 		\begin{enumerate}
 			\item[\ding{68}] La population est l'ensemble sur lequel porte une étude statistique.
 			\item[\ding{68}] Un individu est chaque élément de la population étudiée.
 			\item[\ding{68}]	Le caractère est sur quoi porte une étude statistique. La nature du    caractère peut être :
 			\begin{enumerate}
 				\item[\ding{65}]  qualitative (la couleur des cheveux, les sports pratiqués ou le type de film préféré, \dots ) 
 				\item[\ding{65}] quantitative (la taille, l'âge, le temps passé devant la télévision, \dots ).
 			\end{enumerate}
 			\item[\ding{68}] Les modalités sont les différentes réponses que le caractère permet d'avoir ou toutes valeurs prises par un caractère.
 			\item[\ding{68}]	L'effectif d’une modalité $ (n_i )$ est le nombre d'individus de la population qui représente cette modalité.
 			\item[\ding{68}] L'effectif total $ (N) $ est le nombre total des individus d’une population.
 		\end{enumerate}
 	}
 	\nbVo{Remarque}{ Si $M$ désigne la moyenne d'une série statistique regroupée en classe alors $M=\dfrac{\sum n_{i}c_{i}}{N}$ où $c_{i}$ est le centre de la classe $i$ , $n_{i}$ est l'effectif de la classe $i$ et $N$ est l'effectif total. 
 		
 		\ding{43} Si $[a;b[$ est une classe alors $c_{i}=\dfrac{a+b}{2}$.
 		
 		\textbf{\large Rappel}
 		
 		Dans le cas d'une série statistique simple, la  moyenne $M=\dfrac{\sum n_{i}x_{i}}{N}$ où $x_{i}$ est la modalité ; $n_{i}$ est l'effectif de la modalité $x_{i}$ et $N$ est l'effectif total. }
 	
 	\ret{Si $F$ désigne la fréquence en pourcentage d'une classe alors $F=\dfrac{Effectif \hspace*{0.3cm} de\hspace*{0.3cm} la\hspace*{0.3cm} classe \times100}{Effectif \hspace*{0.3cm} total}$  }
 	\conAct{Approfondissement}
 	Après une évaluation sommative, les notes sur 20 en mathématiques, obtenues par les élèves d'une classe de troisième du CEG DOUNIAN se présentent comme suit:  
 	
 	\begin{tabular}{cccccc}
 		\hline
 		5 & 14 & 20 & 14 & 8& 18 \\
 		18 & 8 & 14 & 18 & 5 & 18 \\
 		\hline
 	\end{tabular}
 	\begin{enumerate}
 		\item Précise: la population, le caractère étudié et sa nature.
 		\item 
 		\begin{enumerate}
 			\item Dresse le tableau des effectifs de cette série statistique.
 			\item Quel est le mode de cette série statistique ? 
 			\item Calcule la moyenne de cette série statistique.	
 		\end{enumerate}
 		\item 
 		\begin{enumerate}
 			\item Repartis ces notes dans les classes $[5;10[$; $[10;15[$ et $[15;20[$ puis dresse le tableau des effectifs et des fréquences de cette série statistique.
 			\item  Donne la classe modale de cette série statistique.
 			\item Calcule la moyenne de cette série statistique.
 		\end{enumerate}
 	\end{enumerate}
 	\newpage
 	\dures
 	 
 	\conAct{Les diagrammes}
 	On considère le tableau de l'activité.
 	\begin{enumerate}
 		\item Construis le diagramme circulaire de cette série statistique.
 		\item Construis le diagramme semi-circulaire de cette série statistique.
 	\end{enumerate}
 	\dures
 	\defs{\ding{70} Un diagramme circulaire d'une série statistique est un disque partagé en secteurs circulaires de mesure proportionnelles aux effectifs ou aux fréquences des modalités ou des classes.\\
 		\ding{70}Un diagramme semi-circulaire d'une série statistique est un demi-disque partagé en secteurs circulaires de mesure proportionnelles aux effectifs ou aux fréquences des modalités ou des classes }
 	
 	\ret{
 		\ding{118} \textbf{\large Angle au centre d'un diagramme circulaire}\\
 		Si $\alpha_{i}$ est la mesure de l'angle au centre associé à la classe $i$ du diagramme circulaire alors\\ $\alpha_{i}=\dfrac{(Effectif \hspace*{0.3cm} de\hspace*{0.3cm} la\hspace*{0.3cm} classe \hspace*{0.2cm} i)\times360^\circ}{Effectif \hspace*{0.3cm} total}$\\ ou bien\\
 		
 		$\alpha_{i}=\dfrac{(Fréquence \hspace*{0.3cm} de\hspace*{0.3cm} la\hspace*{0.3cm} classe \hspace*{0.2cm} i) \times360^\circ}{100}$
 		
 		\vspace*{0.5cm}
 		\ding{118} \textbf{\large Angle au centre d'un diagramme semi-circulaire}\\
 		Si $\alpha_{i}$ est la mesure de l'angle au centre associé à la classe $i$ du diagramme semi-circulaire alors\\ $\alpha_{i}=\dfrac{(Effectif \hspace*{0.3cm} de\hspace*{0.3cm} la\hspace*{0.3cm} classe \hspace*{0.2cm} i)\times180^\circ}{Effectif \hspace*{0.3cm} total}$\\ ou bien\\
 		
 		$\alpha_{i}=\dfrac{(Fréquence \hspace*{0.3cm} de\hspace*{0.3cm} la\hspace*{0.3cm} classe \hspace*{0.2cm} i) \times180^\circ}{100}$}
 	
 	\conAct{Histogramme}
 	On considère le tableau de l'activité\\
 	En  référent d'un diagramme en bâton, fais réalise un histogramme de cette série statistique.\\
 	\dures
 	\defs{Un histogramme est un graphique constitué de bandes juxtaposées dont les aires sont proportionnelles aux effectifs ( ou fréquences) des classes.} 
 	\conAct{Approfondissement}
 	Une série statistique nous a permis d'avoir l'histogramme suivant:
 	
 	
 	\definecolor{ffqqqq}{rgb}{1,0,0}
 	\definecolor{qqwuqq}{rgb}{0,0.39215686274509803,0}
 	\definecolor{qqqqff}{rgb}{0,0,1}
 	\definecolor{zzttqq}{rgb}{0.6,0.2,0}
 	\begin{tikzpicture}[scale=0.6]
 		\begin{axis}[x=1cm,y=1cm,axis lines=middle, xmin=-0.3,xmax=8.5,ymin=-0.3,ymax=7.5,
 			xtick={-1,0,1,...,8},
 			ytick={-1,0,1,...,8},]
 			%	\clip(-0.1,-0.2) rectangle (9,7.5);
 			\fill[line width=2.8pt,color=zzttqq,fill=zzttqq,fill opacity=0.35] (1,0) -- (1,4) -- (2,4) -- (2,0) -- cycle;
 			\fill[line width=0.4pt,color=qqqqff,fill=qqqqff,fill opacity=0.35] (2,3) -- (3,3) -- (3,0) -- (2,0) -- cycle;
 			\fill[line width=2.8pt,color=qqwuqq,fill=qqwuqq,fill opacity=0.35] (4,0) -- (4,6) -- (3,6) -- (3,0) -- cycle;
 			\fill[line width=0.4pt,color=ffqqqq,fill=ffqqqq,fill opacity=0.35] (4,5) -- (5,5) -- (5,0) -- (4,0) -- cycle;
 			\fill[line width=0pt,fill=black,fill opacity=0.35] (5,2) -- (6,2) -- (6,0) -- (5,0) -- cycle;
 			\draw [line width=0.4pt,color=zzttqq] (1,0)-- (1,4);
 			\draw [line width=0.4pt,color=zzttqq] (1,4)-- (2,4);
 			\draw [line width=0.4pt,color=zzttqq] (2,4)-- (2,0);
 			\draw [line width=0.4pt,color=zzttqq] (2,0)-- (1,0);
 			\draw [line width=0.4pt,color=qqqqff] (2,3)-- (3,3);
 			\draw [line width=0.4pt,color=qqqqff] (3,3)-- (3,0);
 			\draw [line width=0.4pt,color=qqqqff] (3,0)-- (2,0);
 			\draw [line width=0.4pt,color=qqqqff] (2,0)-- (2,3);
 			\draw [line width=0.4pt,color=qqwuqq] (4,0)-- (4,6);
 			\draw [line width=0.4pt,color=qqwuqq] (4,6)-- (3,6);
 			\draw [line width=0.4pt,color=qqwuqq] (3,6)-- (3,0);
 			\draw [line width=0.4pt,color=qqwuqq] (3,0)-- (4,0);
 			\draw [line width=0.4pt,color=ffqqqq] (4,5)-- (5,5);
 			\draw [line width=0.4pt,color=ffqqqq] (5,5)-- (5,0);
 			\draw [line width=0.4pt,color=ffqqqq] (5,0)-- (4,0);
 			\draw [line width=0.4pt,color=ffqqqq] (4,0)-- (4,5);
 			
 			\filldraw (0,7) node[right] {Effectif };
 			\filldraw (7,0) node[above] {Modalité };
 			
 			\begin{scriptsize}
 				\draw [color=black] (1,0)-- ++(-0.5pt,0 pt) -- ++(1pt,0 pt) ++(-0.5pt,-0.5pt) -- ++(0 pt,1pt);
 				\draw [color=black] (1,4)-- ++(-0.5pt,0 pt) -- ++(1pt,0 pt) ++(-0.5pt,-0.5pt) -- ++(0 pt,1pt);
 				\draw [color=black] (2,4)-- ++(-0.5pt,0 pt) -- ++(1pt,0 pt) ++(-0.5pt,-0.5pt) -- ++(0 pt,1pt);
 				\draw [color=black] (2,0)-- ++(-0.5pt,0 pt) -- ++(1pt,0 pt) ++(-0.5pt,-0.5pt) -- ++(0 pt,1pt);
 				\draw [color=black] (2,3)-- ++(-0.5pt,0 pt) -- ++(1pt,0 pt) ++(-0.5pt,-0.5pt) -- ++(0 pt,1pt);
 				\draw [color=black] (3,3)-- ++(-0.5pt,0 pt) -- ++(1pt,0 pt) ++(-0.5pt,-0.5pt) -- ++(0 pt,1pt);
 				\draw [color=black] (3,0)-- ++(-0.5pt,0 pt) -- ++(1pt,0 pt) ++(-0.5pt,-0.5pt) -- ++(0 pt,1pt);
 				\draw [color=black] (4,0)-- ++(-0.5pt,0 pt) -- ++(1pt,0 pt) ++(-0.5pt,-0.5pt) -- ++(0 pt,1pt);
 				\draw [color=black] (4,6)-- ++(-0.5pt,0 pt) -- ++(1pt,0 pt) ++(-0.5pt,-0.5pt) -- ++(0 pt,1pt);
 				\draw [color=black] (3,6)-- ++(-0.5pt,0 pt) -- ++(1pt,0 pt) ++(-0.5pt,-0.5pt) -- ++(0 pt,1pt);
 				\draw [color=black] (4,5)-- ++(-0.5pt,0 pt) -- ++(1pt,0 pt) ++(-0.5pt,-0.5pt) -- ++(0 pt,1pt);
 				\draw [color=black] (5,5)-- ++(-0.5pt,0 pt) -- ++(1pt,0 pt) ++(-0.5pt,-0.5pt) -- ++(0 pt,1pt);
 				\draw [color=black] (5,0)-- ++(-0.5pt,0 pt) -- ++(1pt,0 pt) ++(-0.5pt,-0.5pt) -- ++(0 pt,1pt);
 				\draw [color=black] (5,2)-- ++(-0.5pt,0 pt) -- ++(1pt,0 pt) ++(-0.5pt,-0.5pt) -- ++(0 pt,1pt);
 				\draw [color=black] (6,2)-- ++(-0.5pt,0 pt) -- ++(1pt,0 pt) ++(-0.5pt,-0.5pt) -- ++(0 pt,1pt);
 				\draw [color=black] (6,0)-- ++(-0.5pt,0 pt) -- ++(1pt,0 pt) ++(-0.5pt,-0.5pt) -- ++(0 pt,1pt);
 			\end{scriptsize}
 		\end{axis}
 		
 	\end{tikzpicture}
 	
 	\begin{enumerate}
 		\item Dresse le tableau des effectifs et des fréquences de cette série statistique.
 		\item  Réalise le diagramme semi-circulaire de cette série statistique
 	\end{enumerate}
 	\dures 
  
 	\begin{center}
 		\Ovalbox{\textbf{ Exercice de maison}}
 	\end{center}
\Ovalbox{\textbf{ Exercice 1 }}Trois candidats à une élection ont obtenu les résultats donnés dans le tableau suivant.
 
 \begin{center}
 	\begin{tabular}{|l|c|c|c|c|}
 		\hline
 		Candidats & X & Y & Z & Total \\
 		\hline
 		Suffrages & 8350 & 6221 & 831 & \\
 		\hline
 		Pourcentages & & & & \\
 		\hline
 		Angle au centre & & & & \\
 		\hline
 	\end{tabular}
 \end{center}
 
 \begin{enumerate}
 	\item Compléter ce tableau.
 	\item Tracer le diagramme circulaire des résultats.
 \end{enumerate}
 
\Ovalbox{\textbf{ Exercic 2}} Une enquête auprès des 48 élèves de deux classes portait sur la durée du trajet pour se rendre au collège.
 \begin{enumerate}
 	\item Reproduire et compléter le tableau ci-dessous.
 	
 	\begin{center}
 		\begin{tabular}{|c|c|c|}
 			\hline
 			Temps en min & Effectifs & Fréquences en \% \\
 			\hline
 			$0 \leq t < 15$ & 6 & \\
 			\hline
 			$15 \leq t < 30$ & 24 & \\
 			\hline
 			$30 \leq t < 45$ & & \\
 			\hline
 			$t > 45$ & 3 & 6{,}25 \\
 			\hline
 			Total & & \\
 			\hline
 		\end{tabular}
 	\end{center}
 	
 	\item Tracer l'histogramme des effectifs.
 	\item Quel est le nombre d'élèves dont la durée du trajet est inférieure à 45 min ?
 	\item Quel est le nombre d'élèves dont la durée du trajet est supérieure ou égale à 30 min ?
 	\item Quel est le pourcentage d'élèves dont la durée du trajet est inférieure à 30 min ?
 \end{enumerate}
 

 \Ovalbox{\textbf{ Exercice 3}} Dame Awa est celle chargée de vendre de la nourriture. Elle souhaite connaître les habitudes alimentaires de ses clients et profiter pour renforcer la vente de son mets le plus consommé. Le tableau incomplet ci-dessous est celui relatif à la fréquentation journalière de son espace de vente par les clients.
 
 \begin{center}
 	\begin{tabular}{|p{1.5cm}|p{0.5cm}|p{0.5cm}|p{1cm}|p{1cm}|p{1cm}|p{0.5cm}|}
 		\hline
 		Produits & Riz & Pâte & Haricot & Couscous & Attiéké & Total \\
 		\hline
 		Nombre de clients & 17 & & 16 & 15 & & \\
 		\hline
 		Fréquence (en \%) & & & 20 & & 10 & \\
 		\hline
 	\end{tabular}
 \end{center}
 
 \begin{enumerate}
 	\item[1.] (a) Donne le caractère étudié et précise sa nature.
 	
 	(b) Détermine l'effectif total de la population de la série statistique étudiée.
 	
 	(c) Recopie puis complète le tableau statistique ci-dessus.
 	\item Construis le diagramme semi-circulaire de cette série statistique en prenant $5\,cm$ pour rayon et en arrondissant les mesures des angles à l'unité près.
 	\item Donne en te justifiant, le mets le plus consommé par les clients.
 \end{enumerate}
 
\Ovalbox{\textbf{ Exercice 4}} Paul, un pépiniériste, décide de répartir ses plants selon leur taille en centimètre. Il a donc établi le tableau ci-dessous.
 
 \begin{center}
 \hspace*{-0.5cm}	\begin{tabular}{|p{1.2cm}|p{0.6cm}|p{0.8cm}|p{0.8cm}|p{0.8cm}|p{0.8cm}|p{0.8cm}|p{0.5cm}|}
 		\hline
 		Tailles (en cm) & $[0;15[$ & $[15;30[$ & $[30;45[$ & $[45;60[$ & $[60;75[$ & $[75;90[$ & Total \\
 		\hline
 		Effectifs & 50 & & 70 & & 60 & & 500 \\
 		\hline
 		Fréquence (en \%) & & 20 & & 25 & & & \\
 		\hline
 	\end{tabular}
 \end{center}
 
 \begin{enumerate}
 	\item Donne le caractère étudié et précise sa nature.
 	\item[2.] (a) Reproduis et complète le tableau ci-dessus.
 	
 	(b) Précise la classe modale.
 	
 	(c) Détermine le nombre de plants ayant une taille supérieure ou égale à $45\,cm$.
 	\item Construis l'histogramme de la série de tailles relevées.
 \end{enumerate}
 
 	\section*{III- RETOUR ET PROJECTION}
 	\doublebox{\textbf{Activité 1 }} \textbf{Objectivation/auto-évaluation}
 	
 	\begin{enumerate}
 		\item  Fais le point de tout ce que tu as appris sur chaque séquence de cette SA. 
 		\item  Fais le point de tes difficultés et de tes réussites au cours de l'apprentissage.
 		\item  Identifie des situations de la vie courante où tu peux utiliser tes acquis
 	\end{enumerate}
 	\doublebox{\textbf{Activité 2}} \textbf{Réinvestissement}
 	
 	(Séance d'exercice)
 	
 	\begin{center}
 		\it\textit{\textbf{Fin de la $SA_{4} $!!!}}
 	\end{center}
 \end{multicols*}
 
 \end{document}